\documentclass[12pt, a4paper]{article}
\usepackage[utf8]{inputenc}
\usepackage{amsmath, amssymb, amsthm, amsfonts,mathrsfs}
\usepackage{CJKutf8}
\usepackage[margin=1in]{geometry}
\usepackage{tikz-cd}
\newcommand{\spec}{\operatorname{Spec}}
\newcommand{\spa}{\operatorname{Span}}
\newcommand{\p}{\mathfrak{p}}
\renewcommand{\a}{\mathfrak{a}}
\renewcommand{\b}{\mathfrak{b}}
\newcommand{\z}{\mathbb{Z}}
\newcommand{\Q}{\mathbb{Q}}
\newcommand{\R}{\mathbb{R}}
\newcommand{\q}{\mathfrak{q}}
\newcommand{\tr}{\text{Tr}}
\newcommand{\disc}{\text{disc}}
\renewcommand{\o}{\mathcal{O}}
\newtheorem{lemma}{Lemma}


\newenvironment{solution}
  {\paragraph{\textit{Solution:}}\par\medskip}
  {\hfill$\blacksquare$\vspace{0.1cm}}

\title{Algebraic Number Theory:Assignment}
\author{
    \begin{CJK*}{UTF8}{gbsn}
    钟皓宇 (12533556)
    \end{CJK*}
}
\date{\today}

\begin{document}

\maketitle
\newcounter{paranum} 
\newcommand{\np}{\stepcounter{paranum}\noindent\textbf{\arabic{paranum}.}\quad}

\section*{Some Examples}
Let $R$ be an integral domian with fraction field $F$. Let $K$ be a field extension $F$ and let $V$ be a finite dimensional $K$-vector space.

\np $L$ is an $R$-submodule of $F\otimes_{R} L$. To anlysis, we can consider 
\begin{align*}
    (v_1,\dots,v_m) = (u_1,\dots,u_n)P
\end{align*} 
Where $(u_i)$ and $(v_i)$ is the basis of $F\otimes_{R} L$ and generators of $L$.
We can also extend $v_1,\dots,v_m$ to a basis of $F\otimes_{R} L$ by some consideration.

\np An example of $R$-lattice not support in $V$. 
Let $K=V=\mathbb{R}$ and $R=\z$. Then
$L= \z+\z\sqrt{2}$ is a $\z$-lattice but not support in $\mathbb{R}$.

\np It is easy to see that from 2., $\dim L\ge \dim_K\spa_K{L}$ and a system of vectors in $L$ is linearly independent over $R$ iff. over $K$.

\np We can clearly see from 2. the equivalence of existence of $K$-free basis and support in $V$.

\np Let $L= 2\cdot \z+\z(1+\sqrt{-5})$, $K = \Q(\sqrt{2})$ and $V = \Q(\sqrt{2},\sqrt{-5})$. Then $L$ is support in $V$.
$L$ is a non-trivial generator of ideal class group of $\Q(\sqrt{-5})$.


\np Let $L= \z+\z\sqrt{-5}$ and $L' = 2\cdot\z+\z(1+\sqrt{-5})$. Then $\det(L/L')=2 \cdot \z^*$ and $2L\subseteq L'$.

\textbf{Norm and discriminant in Algebraic Number Theory}

We consider $L/K$ is a seperable finite algebraic extension. We specifically let $O_K$ is not a PID. 
Let $I$ be a fraction ideal of $O_L$. Then there exists a basis $\{\omega_i\}$ of $L$ over $K$ such that
$$ I = \mathfrak{a}_1 \omega_1 \oplus \mathfrak{a}_2 \omega_2 \oplus \dots \oplus \mathfrak{a}_n \omega_n. $$
Similarly, we have 
$$ \mathcal{O}_L = \mathfrak{b}_1 \nu_1 \oplus \mathfrak{b}_2 \nu_2 \oplus \dots \oplus \mathfrak{b}_n \nu_n.$$
Suppose that $P$ is the transform matrix such that $(\omega_i) = P(\nu_i)$.
In this sense, the norm of $I$ is given by $$\det(P)\cdot \frac{\mathfrak{a}_1\cdots \mathfrak{a}_n}{\mathfrak{b}_1\cdots \mathfrak{b}_n}$$
The discriminant of $I$ is given by 
$$(\mathfrak{a}_1 \mathfrak{a}_2 \dots \mathfrak{a}_n)^2 \cdot \det\left(\text{Tr}_{L/K}(\omega_i \omega_j)\right).$$
The discriminant of field $L$ is the discriminant of $O_L$.
If $L=K(\theta)$ with miminal polynomial $p(x)$ then the discriminant $\text{dics}(f)$ of $f$ is given by the discriminant of lattice generated by $\{\theta^n\}$.

\np Let $L= \z+\z\sqrt{-5}$ and $L' = \z(3-\sqrt{-5})+\z(9+\sqrt{-5})$. Let $P$ be the transition matrix. Here is the triangular process of $P$.
Let $L'^{(1)}= 3\mathbb{Z}$ and $L'^{(2)} = L$. Then $\pi_1(L'^{(1)}) = 3\z$ and $\pi_2(L'^{(2)}) = \sqrt{-5}\z$ is principal.
Let $d=\det(L/L') = 12$. Then $12\mathbb{Z}\subseteq 3\mathbb{Z}$ and $12 \z\sqrt{-5}\subseteq \z\sqrt{-5}$ it implies that $\pi_i(L'^{(i)})$ is not zero.
Choose $f_i\in L'^{(i)}$ such that $\pi_i(f_i)=t_{ii}$ where $t_{ii}$ is the generator of ideal $\pi_i(L'^{(i)})$.

The core of process above is the coefficient ideal $\pi_i(L'^{(i)})$. And the condition $L'\subseteq L$ guarantee the projection is not zero ideal.

\np Resay a trick. If we don't have $L'\subseteq L$. we just need to consider $dL'\subseteq L$.

We consider field extension $L/K$ equipped with bilinear form $\tr_{L/K}$.

\np If $L/K$ is separable then $\tr$ is not singular.

\np If $L/K$ is inseparable extension, then $\tr$ is totally singular.
Let $L=\mathbb{F}_p(t^{1/p})$ and $K = \mathbb{F}_p(t)$ then $\tr(L)=p$.

\np Isotropic case: Let $L=\Q(i)$ and $L=\Q$. Consider $\tr((1+i)^2)= 0$.

\np Anisotropic case: Let $L=\Q(\sqrt{2})$ and $L=\Q$. Consider $\tr(x^2)= x^2+\bar{x}^2\neq 0$ if $x\neq 0$.

\np Alternative case: Twisted Trace. Let $L=\Q(i)$ and $L=\Q$ with $B(x,y)=\tr(i\cdot \bar{x}y)$.


\textbf{Twisted Trace Form and Different}

Let $L/K$ be a finite seperable algebraic extension. The Different of $L/K$ is defined by $(^* \o_L)^{-1}$.

\np Conisder vector space $\mathbb{F}_2x$ with quadratic form $Q(x)= 1$. Then radical is $\mathbb{F}_2x$ but it is not singular.




\section*{A.1 \quad Homework 1}

Throughout what follows, let $R$ denote an integral domain with fraction field $K$.

\subsection*{Exercise A.1.1.} 
Suppose $R$ is a local domain. Let $I$ be a fractional ideal of $R$. Prove that $I$ is invertible if and only if it is principal. \\
\emph{Hint: There exist $x \in I$, $y \in I^{-1}$ such that $xy \in R^{\times}$. Then prove $I = xR$.}
\begin{solution}
    Suppose $x\neq 0$ and $I=xR$ is principal then let $J=x^{-1}R$ and $IJ=R$ implies $I$ is invertible. Conversely, if there exist a fractional ideal $J$ such that $IJ=R$. Then there exist $x_i\in I$ and $y_i\in J$ such that
    \begin{align*}
        \sum x_iy_i = 1.
    \end{align*} 
    Because $R$ is local, without loss of generality, we can suppose $x_1y_1\in R\backslash \mathfrak{m}=R^{\times}$. Since $x_1\in I$, $x_1R\subseteq I$ is obvious. Hence $y_1\in J$ implies $y_1I\subseteq R$, multiplying by $x_1$ on both sides, we obtain that 
    $I\subseteq x_1 R$. Therefore, $I=x_1R$ is principal.
\end{solution}


\subsection*{Exercise A.1.2.} 
Prove that a Dedekind domain with only finitely many prime ideals is a PID.
\begin{solution}
    Let $\mathfrak{p}_1,\dots,\mathfrak{p}_n$ be all prime ideals of Dedekind domain $R$.
    It suffices to prove all prime ideals are principal. Without loss of generality, we prove $\mathfrak{p}_1$ is principal and others are in 
    the same way. Let $x\in \mathfrak{p}_1\backslash \mathfrak{p}_1^2$ and by Chinese Remainder Theorem, there exist $u\in R$ such that 
    \begin{align*}
        u\equiv x \bmod \mathfrak{p}_1^2 \text{ and } u\equiv 1 \bmod \mathfrak{p}_i\text{ for }i\neq 1.
    \end{align*} 
    Let $v_1,\dots,v_n$ be the valuations of local rings $R_{\mathfrak{p}_1},\dots,R_{\mathfrak{p}_n}$ respectively. Then we obtain that 
    \begin{align*}
        v_1(u) = 1 \text{ and } v_i(u) = 0\text{ for }i\neq 1,
    \end{align*}
    because $u-x\in \mathfrak{p}_1^2$ and implies $u\in (x)+\mathfrak{p}_1^2\subseteq \mathfrak{p}_1$ meaning that $v_1(u)\ge 1$. Suppose $v_1(u)\ge 2$, equivalently $u\in \mathfrak{p}_1^2$, then $x\in (u)+\mathfrak{p}_1^2 = \mathfrak{p}_1^2$ contradicting our
    assumption, therefore $v_1(u)=1$. Because $R$ is Dedekind domain, the unique factorization of fractional ideals implies 
    \begin{align*}
        (u) = \mathfrak{p}_1^{v_1(u)}\dots\mathfrak{p}_n^{v_n(u)} = \mathfrak{p}_1,
    \end{align*}
    the principality of $\mathfrak{p}_1$ follows.
\end{solution}



\subsection*{Exercise A.1.3.} 
Suppose $R$ is a DVR with maximal ideal $\mathfrak{p}$. Let $k = R/\mathfrak{p}$ and $U = R^{\times}$. Let $v$ be the normalized discrete valuation of $K$ associated to $R$. For each $n \geq 1$, define
\[
U_n = 1 + \mathfrak{p}^n = \{x \in K \mid v(x - 1) \geq n\} \, .
\]

\begin{enumerate}
    \item Show that each $U_n$ is a subgroup of $U$ and that $\bigcap_{n=1}^{\infty} U_n = 1$.
    \item Prove that the natural quotient map $R \to k = R/\mathfrak{p}$ induces an isomorphism of multiplicative groups $U/U_1 \cong k^*$.
    \item Prove that for each $n \in \mathbb{Z}$, there is an isomorphism of $R$-modules $k = R/\mathfrak{p} \xrightarrow{\sim} \mathfrak{p}^n/\mathfrak{p}^{n+1}$.
    \item Prove that for each $n \geq 1$, the map $u \mapsto u - 1$ induces an isomorphism of abelian groups $U_n/U_{n+1} \xrightarrow{\sim} \mathfrak{p}^n/\mathfrak{p}^{n+1}$.
    \item Let $p$ be the characteristic of $k$. (So $p = 0$ or $p$ is a prime number.)
    \begin{enumerate}
        \item[(a)] Show that $U_n^p := \{a^p \mid a \in U_n\}$ is contained in $U_{n+1}$ for every $n \geq 1$.
        \item[(b)] Let $m \geq 1$ be an integer not divisible by $p$. Prove that for every $n \geq 1$, the group homomorphism
        \[
        U_n \longrightarrow U_n \ ; \quad u \longmapsto u^m
        \]
        is injective.
    \end{enumerate}
\end{enumerate}

\begin{solution}
    \textbf{(1)} Closed under multiplication holds since $(1 + \mathfrak{p}^n)(1 + \mathfrak{p}^n)\subseteq 1 + \mathfrak{p}^n$. 
    Identity element is $1$. Suppose $x\in U_n$, then $v(x^{-1}-1) = v(x^{-1})+v(1-x)$ hence $x$ is invertible meaning that $v(x)=0$ and 
    $v(1-x) = v(x-1)$ implies $v(x^{-1}-1)\ge n$. Suppose $x\in \bigcap_{n=1}^{\infty} U_n $, then we obtain $v(x-1) = \infty$, equivalently, $x=1$.
    \textbf{(2)} The natural quotient map restricted to $U$, denoted by $\pi:U\to k^*$, induces isomorphism $U/U_1 \cong k^*$ because $\ker(\pi) = U_1$.
    \textbf{(3)} Consider commutative graph 
\[
\begin{tikzcd}[row sep=large, column sep=large]
0 \arrow[r] & \mathfrak{p} \arrow[r, hook] \arrow[d, "\cdot u^n"', "\cong"] & R \arrow[r, two heads] \arrow[d, "\cdot u^n"', "\cong"] & k \arrow[r] \arrow[d, dashed, "\bar{\phi}", "\cong"'] & 0 \\
0 \arrow[r] & \mathfrak{p}^{n+1} \arrow[r, hook] & \mathfrak{p}^n \arrow[r, two heads] & \mathfrak{p}^n/\mathfrak{p}^{n+1} \arrow[r] & 0
\end{tikzcd}
\] 
which the third isomorphism follows from the Five Lemma.
\textbf{(4)} Consider a surjective group homomorphism 
\begin{align*}
    \varphi:U_n\to \mathfrak{p}^n/\mathfrak{p}^{n+1} \quad u\mapsto u-1 + \mathfrak{p}^{n+1}
\end{align*}
with kernel $U_{n+1}$, then the first isomorphism theorem yields isomorphism $U_n/U_{n+1} \cong \mathfrak{p}^n/\mathfrak{p}^{n+1}$.
\textbf{(5.a)} For $p=0$, obvious to see that $U_n^0=\{1\}\subseteq U_{n+1}$. For $p\neq 0$ and any $a\in U_n$, 
\begin{align*}
    v(a^p-1) =  v((a-1)^p) = p\cdot v(a-1)\ge n+1
\end{align*} 
implies $a^p\in U_{n+1}$.
\textbf{(5.b)} It suffices to prove for $m=q$ where $q$ is a prime number and coprime to $p$. Suppose that $u\in U_n$ such that $u^q =1$. 
According to Bézout's identity and \textbf{(5.a)}, suppose $ap+bq=1$, then $u = u^{ap+bq} = u^{ap}\in U_{n+1}$. Thus we can consider a new group homomorphism:
$U_{n+1} \to U_{n+1}$. Continuing this process, we obtain that $u\in \bigcap_{n=1}^{\infty} U_n$, then $u=1$ and injectivity holds.
\end{solution}

\subsection*{Exercise A.1.4.} 
Prove that if a domain $R$ is integrally closed, then so is the polynomial ring $R[t]$.
\begin{solution}
Let $f=p/q$ be an arbitrary integral element in $R(t)$. Without loss of generality, we can suppose $q$ is irreducible, because if $q=q_1q_2$ then we consider
integral element $fq_1$, and suppose $p,q$ are coprime. Suppose $f$ satisfies a monic polynomial $f^n+a_{n-1}f^{n-1}+\dots+a_0=0$ where $a_k\in R[t]$. Then, by multiplying both side by $q^n$ 
, we obtain that $p^n+a_{n-1}qp^{n-1}+\dots+a_0q^n=0$, which implies that $q|p^n$. Let $F=\text{Frac}(R)$ be the fraction field of $R$, and 
interprete $p,q$ as elements in $F[t]$, then we obtain that $q|p$ since $F[t]$ is an UFD and hence $q\in F^*$ because $p,q$ coprime.
Therefore, we can write $f=c_mt^m+\dots+c_0$ where $c_k\in F$.  Let $f_1 = t^r - f$, where $r=\deg f+\sum \deg a_k$, then $f_1$ is the root of polynomial
\begin{align*}
    P(X)=(t^r-X)^{n}+a_{n-1}(t^r-X)^{n-1}+\dots + a_0.
\end{align*}
Let $X=0$, we obtain the constant term of this polynomial
\begin{align*}
    t^{rn}+a_{n-1}t^{rn-r}+a_{n-2}t^{rn-2r}+\dots+a_0
\end{align*}
denoted by $P(0)$. Let $P(X)=XQ(X)+P(0)$, then from $P(f_1)=0$ we obtain that 
\begin{align*}
    P(0) = -f_1Q(f_1),
\end{align*} 
the roots of $f_1$ are among the roots of monic polynomial $P(0)\in R[t]$, the coefficients of $f$ as of $f_1$ are elementary symmetric polynomials of its roots, which are integral over $R$. Therefore, we obtain that the coefficients of 
$f$ are contained in $R$, $R[t]$ is integrally closed.  
\end{solution}



\subsection*{Exercise A.1.5.} 
Let $A = \mathbb{Q}[X, Y]/(X^2 - Y^3)$. Show that $A$ is a domain but it is not integrally closed.
\begin{solution}
    $(X^2-Y^3)$ is a prime ideal because $X^2-Y^3$ is irreducible and $\mathbb{Q}[X,Y]$ is UFD, \textit{a fortiori}, $A$ is a domain.
    Hence, $X/Y$ is the root of a monic polynomial $P(Z) = Z^2-Y$ but not contained in $A$. If it is contained, then there exist $f,g\in\mathbb{Q}[X,Y]$ such that 
    \begin{align*}
        X = Yf-(X^2-Y^3)g
    \end{align*}
    implies contradiction by comparing their linear terms.
\end{solution}

\subsection*{Exercise A.1.6.} 
Let $R$ be a Dedekind domain and let $\mathfrak{a} \subseteq R$ be a nonzero ideal. Show that the quotient ring $R/\mathfrak{a}$ has only finitely many ideals and that all of these ideals are principal.
\begin{solution}
    By isomorphism theorem of rings, the natural quotient $R\to R/\mathfrak{a}$ gives one-to-one correspondence of the ideals of $R$ containing $\mathfrak{a}$ and the ideals of $R/\mathfrak{a}$.
    Therefore, all ideals of $R/\mathfrak{a}$ can be written as the form 
    \begin{align*}
        \mathfrak{p}_1^{i_1}\dots \mathfrak{p}_n^{i_n} + \mathfrak{a}
    \end{align*}
    where $\mathfrak{p}_k$ are the prime ideals of $R$ containing $\mathfrak{a}$ and $0\le i_k\le v_k(\mathfrak{a})$ where $v_k$ is the valuation associated to $\mathfrak{p}_k$.
    The principality directly from \textit{Exercise A.1.2.}
\end{solution}
\subsection*{Exercise A.1.7.} 
Show that every integral ideal of a Dedekind domain can be generated by two elements.
\begin{solution}
    Let $\mathfrak{a}$ is an arbitrary ideal of $R$ with prime ideal decomposition
    \begin{align*}
        \mathfrak{a} = \mathfrak{p}^k\mathfrak{b}
    \end{align*}
    Where $\mathfrak{p}$ is a prime ideal coprime to the ideal $\mathfrak{b}$.
    The quotient ring $R/\mathfrak{p}^{k+1}\mathfrak{b}$ is principal from \textit{Exercise A.1.6.} Therefore
    we can assume that $\mathfrak{a} = (a)+\mathfrak{p}^{k+1}\mathfrak{b}$. Hence, consider another quotient ring $R/(a)$, we obtain that $\mathfrak{a}=(a,b)$. 
\end{solution}

\section*{A.2 \quad Homework 2}
\subsection*{Exercise A.2.1}
Let $R$ be the ring of algebraic integers in $\mathbb{C}$, i.e., the integral closure of $\mathbb{Z}$ in $\mathbb{C}$. Consider the ideal $I \subseteq R$ generated by the infinite subset $\{\sqrt[2^n]{2} \mid n \in \mathbb{N}\}$. 

Show that $I$ is not finitely generated. Deduce that $R$ is not Noetherian (and hence not a Dedekind domain).

\begin{solution}
    Suppose that $I$ is finitely generated with generators $\alpha_1,\dots,\alpha_m$.
    According to Chevalley's Extension Theorem for Valuations, we can extend 2-adic valuation on $\mathbb{Q}$ to $\overline{\mathbb{Q}}$, let $v_2$ denote this extended valuation.
    By definition, we obtain that $R\subseteq \overline{\mathbb{Q}}$ and $v_2(R)\ge 0$.
    Since $I$ is generated by $\{\sqrt[2^n]{2} \}$ and hence $v_2(\sqrt[2^n]{2})= 2^{-n}$ is strictly greater than zero, the valuation $v_2(\alpha_i)$ is strictly greater than zero.
    Let $\varepsilon = \min\{v_{2}(\alpha_i)\}$ be the lower bound of valuations of finite generators $\alpha_i$.
    For arbitrary $r \in I$, we can write $r = \sum_{i=1}^m c_i \alpha_i$ for some $c_i \in R$. Using the non-Archimedean property and $v_2(c_i) \ge 0$, we obtain:$$v_2(r) \ge \min_{i} \{v_2(c_i) + v_2(\alpha_i)\} \ge \min_{i} \{v_2(\alpha_i)\} = \varepsilon > 0$$
    However we find that $v_{2}(\sqrt[2^n]{2}) = 2^{-n}$ converging to zero, which is a contradiction.
\end{solution}
\subsection*{Exercise A.2.2}
Let $R$ denote a Dedekind domain with fraction field $K$ and suppose $R \neq K$. 

Prove that there exists a set $\Omega$ of normalized discrete valuations of $K$ such that the following two properties hold:
\begin{itemize}
    \item[(R1)] $R = \{x \in K \mid v(x) \geq 0 \text{ for all } v \in \Omega\}$;
    \item[(R2)] for every $a \in K^\times$, the set $\{v \in \Omega \mid v(a) \neq 0\}$ is finite.
\end{itemize}

\begin{solution}
    Let 
    \begin{align*}
        \Omega = \{v_{\mathfrak{p}} \mid \mathfrak{p} \in \text{Spm}(R)\}
    \end{align*}
    where $v_{\mathfrak{p}}$ is the normalized discrete valuation associated with $R_{\mathfrak{p}}$.
    For (R1), we find that $v_{\mathfrak{p}}(x)\ge 0$ if and only if $x\in R_{\mathfrak{p}}$, then we obtain that
    \begin{align*}
        \{x \in K \mid v(x) \geq 0 \text{ for all } v \in \Omega\} = \bigcap_{\mathfrak{p}\in\text{Spm}(R)}R_{\mathfrak{p}} = R.
    \end{align*}
    For $(R2)$, we first fix an arbitrary element $a\in K^*$ and let $A$ denote the set $\{v \in \Omega \mid v(a) \neq 0\}$.
    Since $R$ is Dedekind, we may suppose that 
    \begin{align*}
        aR = \mathfrak{p}_1^{a_1}\dots \mathfrak{p}_n^{a_n}.
    \end{align*}
    where $a_i\neq 0$.
    I claim that $A = \{v_{\mathfrak{p}_1}, \dots, v_{\mathfrak{p}_n}\}$. For an arbitrary maximal ideal $\mathfrak{q} \notin \{\mathfrak{p}_1, \dots, \mathfrak{p}_n\}$, it does not appear in the factorization of $aR$, yielding $v_{\mathfrak{q}}(a) = 0$ and hence $v_{\mathfrak{q}} \notin A$. This implies $A \subseteq \{v_{\mathfrak{p}_1}, \dots, v_{\mathfrak{p}_n}\}$. Conversely, by definition we have $v_{\mathfrak{p}_i}(a) = a_i \neq 0$, meaning $\{v_{\mathfrak{p}_1}, \dots, v_{\mathfrak{p}_n}\} \subseteq A$. 
    Therefore, $A$ is precisely this finite set. 
\end{solution}

\subsection*{Exercise A.2.3}
Let $R$ denote an integral domain with fraction field $K$. Let $\Omega$ be a set of normalized discrete valuations of $K$ satisfying Properties (R1) and (R2) in Exercise A.2.2. For each $v \in \Omega$, let $R_v$ be its valuation ring and let $\mathfrak{p}_v$ be the maximal ideal of $R_v$.

Suppose further that every nonzero prime ideal of $R$ is maximal.

\begin{enumerate}
    \item Let $\mathfrak{p}$ be a maximal ideal of $R$.
    \begin{enumerate}
        \item Let $z \in K^\times$ be such that $R \cap z^{-1}R \subseteq \mathfrak{p}$. Show that there exists $v \in \Omega$ such that $\mathfrak{p} = R \cap \mathfrak{p}_v$ and $v(z) < 0$.
        \item Prove that $\Omega_{\mathfrak{p}} := \{v \in \Omega \mid R_v \supseteq R_{\mathfrak{p}}\}$ is a finite set.
        \item Prove that $R_{\mathfrak{p}} = R_v$ for a unique $v \in \Omega$. Thus $\Omega_{\mathfrak{p}}$ consists of a single element.
    \end{enumerate}
    
    \item Prove that for every $v \in \Omega$, there exists a unique maximal ideal $\mathfrak{p}$ such that $R_v = R_{\mathfrak{p}}$.
    
    \item Deduce from (1) and (2) that the map
    \[ \Omega \longrightarrow \text{Spm}(R) \ ; \ v \longmapsto \mathfrak{p}_v \cap R \]
    is well defined and bijective. Let $\mathfrak{p} \mapsto v_{\mathfrak{p}}$ denote its inverse map.
    
    \item Let $I \subseteq R$ be a nonzero ideal.
    \begin{enumerate}
        \item Prove that $\{\mathfrak{p} \in \text{Spm}(R) \mid v_{\mathfrak{p}}(I) \neq 0\}$ is finite, where $v_{\mathfrak{p}}(I) := \min\{v_{\mathfrak{p}}(x) \mid x \in I\}$.
        \item Prove that there is a unique factorization $I = \prod_{\mathfrak{p} \in \text{Spm}(R)} \mathfrak{p}^{r_{\mathfrak{p}}}$, where each $r_{\mathfrak{p}} \in \mathbb{N}$ and almost all of them are zero.
        \item If $J \subseteq R$ is another nonzero ideal, show that $I \mid J$ if and only if $I \supseteq J$.
    \end{enumerate}
    
    \item Prove that $R$ is a Dedekind domain.
\end{enumerate}
\begin{solution}
(1.a) Let $\Gamma =\{v\in \Omega|v(z)<0\}$. According to (R2), it is finite. Hence, it cannot be emptyset since if so then $z^{-1}R\cap R = R \subseteq \p$ is impossible.
We have 
\begin{align*}
    R\cap z^{-1}R = \{r\in R|v(r)\ge -v(z)\text{ for any }v\in \Gamma\}.
\end{align*}
If our claim is not hold, that is, for any $v\in \Gamma$ we then have $\p\neq R\cap \p_v$.
We obtain pairwise coprime ideals $\p,\p_i$ (Here we can suppose that $\p_i\cap R\neq \p_j\cap R$ by discarding the same ideals and choose the maximal valuation on $z$)
and consider CRT by 
\begin{align*}
    r \equiv 1 \bmod \p \quad r\equiv  0 \bmod (\p_{v_i}\cap R)^{-v_i(z)}
\end{align*}
Then we find that $r\in R\cap z^{-1}R$ and $r\notin \p$. It is a contradiction.

(1.b)
Suppose that $\Omega_{\p}$ is an infinite set and for arbitrary $v\in\Omega_\p$ we can have
\begin{align*}
    \p_v\cap R_\p = \p R_\p.
\end{align*}
Since $\p_v\cap R $ is a prime ideal and therefore it is maximal. $\p_v\cap R_\p$ is the prime ideal of $R_\p$ lying over $\p_v\cap R$, then we have $\p_v\cap R_\p$ is also maximal implying what we want.
Therefore 
\begin{align*}
    \bigcap_{v\in\Omega_\p}\p_v \supseteq \p R_{\p}
\end{align*}
For arbitrary nonzero element $x\in\p R_\p$ and $v\in\Omega_\p$, we have $v(x)\neq 0$ which contradicts to (R2) since we suppose $\Omega_\p$ is an infinite set.


(1.c)

Choose $z^{-1}\in\p$ and (1.a) implies that there exists $v$ such that $\p=R\cap \p_v$ and $v(z)<0$. 
It implies that $R_v\supseteq R_\p$. Therefore $\Omega_\p$ is not an emptyset.

First I claim that 
\begin{align*}
    R_\p = \bigcap_{v\in\Omega_\p} R_v.
\end{align*}
It is easy to see that $R_\p \subseteq \bigcap_{v\in\Omega_\p} R_v$ by definition. Conversely, suppose that $x\in \bigcap_{v\in\Omega_\p} R_v$. Then we obtain that $v'(x)\ge 0$ for arbitrary $v'\in\Omega_\p$.
For arbitrary $v\in (\Omega\setminus \Omega_\p)$, we have $\p\neq R\cap \p_v$. If not, $v$ must contained in $\Omega_\p$. Therefore, by (1.a), we find that $R\cap x^{-1}R$ is not contained in $\p$.
Suppose that $y\in (R\cap x^{-1}R)-\p$, then we find for arbitrary $v\in\Omega$, $v(y),v(xy)\ge 0$ and for arbitrary $v'\in\Omega_\p$, we have $v'(y)=0$. It means that $xy\in R\subseteq R_\p$ and $y\in R-\p$.
Since $v'(R-\p)= v'((R_{v'}-\p_{v'})\cap R)=0$ and $v'(\p)=v'(\p_{v'}\cap R)\ge 1$, then it must have $y\in R-\p$, therefore $x\in R_\p$ meaning that $R_\p \supseteq \bigcap_{v\in\Omega_\p} R_v$.

Select an arbitrary valuation $v\in \Omega_\p$, I claim that $R_v=R_\p$. From the discussion of (1.b), we obtain that 
\begin{align*}
    v_\p\cap R_\p = \p R_p.
\end{align*}
It means that $v(R_\p-\p R_p)= 0$. Hence, we have 
\begin{align*}
    v^2_\p\cap R_\p = \p^2 R_p.
\end{align*}
We notice that $v^2_\p\cap R_\p\supseteq \p^2 R_p$ since $v(\p^2 R_\p)\ge v(\p^2)\ge 2$. Therefore, $v^2_\p\cap R_\p$ must equals $\p^2 R_\p$ or $\p R_\p$. If it equals $\p R_\p$,
then for arbitrary $x\in K^*$, if $v(x)=1$ then we have $x\notin R_\p$, which is wrong. To see that, we choose an arbitrary element $x\in K^*$ such that $v(x)=1$.
Suppose that $\Omega_\p=\{v,v_1,\dots,v_k\}$.
Then we use the weak approximation theorem, there exists an element $y\in K^*$ such that $v(y)=0$ and $v_i(y)=-v_i(x)$.
Then we find that $v(xy)=1$ and $v_i(xy)=0$ meaning that $xy\in R_\p$. Therefore, it must have $v^2_\p\cap R_\p=\p^2 R_\p$. Use same method, we can prove that $v^j_\p\cap R_\p=\p^j R_\p$ for any positive integer $j$. 
It meaning that $R_v= R_\p$. Since select arbitrarily, the uniqueness is obtained.

(2) Let $\p=\p_v\cap R$. It is a prime ideal and hence is maximal ideal. We find that $v(R-\p)=0$ and therefore $v(R_\p)\ge 0$ implying $R_p\subseteq R_v$. 
By (1), we have $R_\p=R_v$. If there exists a different maximal ideal $\mathfrak{q}$ such that $R_v=R_\mathfrak{q}$. Then it means that $\p_v$ is lying over $\mathfrak{q}$, then it must equal to $\p$.

(3) We just need to prove that $R_\p\subseteq R_v$ according to (1). It is quite clear from the construction of (2).

(4.a) We notice that $v_{\p}(I)\neq 0$ if and only if $I\subseteq \p$. If set $\{\mathfrak{p} \in \text{Spm}(R) \mid v_{\mathfrak{p}}(I) \neq 0\}$ is not finite, It means that 
for arbitrary element $x\in I$, there exists infinite valuations $v$ such that $v(x)\neq 0$.
Then by (R2) we must have $I=(0)$ 
contradicting to our condition.

(4.b) Let $J$ denote $\prod_{i=1}^n \mathfrak{p}^{r_{\mathfrak{p_i}}}$ such that for $\p\notin\{\p_1,\dots,\p_n\}$ we have $v_\p(I)=0$ and
$r_\p=v_\p(I)$.
We notice that $I_\p = \p^{r_\p}R_\p$ since
\begin{align*}
    I_\p \subseteq \p^{r_\p}R_\p \text{ and } I_\p \subsetneq \p^{r_\p+1}R_\p
\end{align*}
Then we have 
\begin{align*}
    I= \bigcap_{\p \in \text{Spm}(R)} I_\p =\bigcap_{i=1}^{n} I_{\p_i}= J.
\end{align*}
The last step is because $\p_i$ are pairwise coprime.
(4.c) It is directly deduced by (4.b).

(5) Noetherian is due to (4.c). For intergrality, we consider ideal $(x)$ generated by an arbitrary element $x\in R$.
\end{solution}


\section*{Homework 3}
\subsection*{Exercise 3.1}
Let $L$ be an $R$-lattice contained in the $K$-vector space $V$.
Prove that $L$ is an $R$-lattice supported in $V$ if and only if $FL$ has a $K$-free basis.
\begin{solution}
    Suppose that $L$ is supported in $V$. It means that $\dim_F(FL)=\dim_K(KL)$.  Choose a basis $\{e_i\}$ of $FL$.
    Since $\{e_i\}$ spans $KL$, then if $\{e_i\}$ is not a basis of $KL$, we then have $\dim_K(KL)<\dim_F(FL)$. This is a contradiction.
    Conversely, it is trivial.
\end{solution}


\section*{Homework 4}

\subsection*{Exercise 4.1}
Let $\mathfrak{p}_1, \cdots, \mathfrak{p}_m$ be prime ideals of a ring $R$. Prove that for any ideal $I \subseteq R$, $I \subseteq \bigcup_{i=1}^m \mathfrak{p}_i$ if and only if $I \subseteq \mathfrak{p}_i$ for some $1 \leq i \leq m$.
\begin{solution}
    Suppose that $I \subseteq \bigcup_{i=1}^m \mathfrak{p}_i$. 
    We use induction on $n$ and starting by $n=2$.
    Suppose that $a_1\in \p_1\cap \p_2^c\cap I$ and $a_2\in \p_1^c\cap\p_2\cap I$ where both $a_1$ and $a_2$ are not equal to zero.
    Since $I$ is an ideal, we then have $a_1+a_2\in I$. Clearly, $a_1+a_2$ must not contained in $\p_1$ and $\p_2$. Otherwise, we then have $a_1\in\p_2$ or $a_2\in\p_1$. It is a contradiction.

    Suppose our claim holds for $n-1$. Suppose that there exist an element $a_i\in I\cap \p_i$ such that $a_i$ is not contained in any $I\cap\p_j$ for $j\neq i$. Otherwise, we can reduce it to $n-1$ case.
    Then we consider the elements $a=\sum_{i}\prod_{k\neq i} a_k$. Since $I$ is an ideal then we have $a\in I$. Suppose that $a\in\p_j\cap I$.
    Then we obtain that $\prod_{k\neq j} a_k\in \p_j$. Since $\p_j$ is a prime ideal and $a_i\notin \p_j$ for $i\neq j$, we then have $\prod_{k\neq j} a_k\notin \p_j$.
    It is a contradiction.
    Therefore, our claim holds.
\end{solution}



\subsection*{Exercise 4.2}
Let $R$ be a Dedekind domain, and let $S \subseteq R \setminus \{0\}$ be a multiplicative subset.
\begin{enumerate}
    \item Prove that the ring of fractions $S^{-1}R$ is a Dedekind domain.
    \item Suppose that $S = R \setminus \bigcup_{i=1}^m \mathfrak{p}_i$ for finitely many maximal ideals $\mathfrak{p}_1, \cdots, \mathfrak{p}_m$ of $R$. Prove that $S^{-1}R$ is a PID.
\end{enumerate}
\begin{solution}
    (1) For each $\p\cap S=\emptyset$, we have 
    \begin{align*}
        (S^{-1}R)_{S^{-1}\p} =  S^{-1} R\otimes_R (R\setminus \p)^{-1}= S^{-1}R_\p = R_\p
    \end{align*}
    Since $R$ is Dedekind domain, we then have $R_\p$ is DVR. Therefore, $S^{-1}R$ is also a Dedekind domain.

    (2)
    Since Dedekind domain $S^{-1}R$ only have finite prime ideal $S^{-1}\p_i$. According to Excercise 1.2, it is a PID.
\end{solution}




\subsection*{Exercise 4.3}
Let $R$ be a Dedekind domain, and let $M \neq 0$ be a finitely generated $R$-module. Let $M_{\mathrm{tors}}$ be the torsion submodule of $M$. The invariants factors of $M$ are defined as those of $M_{\mathrm{tors}}$ (cf. Thm. 1.2.27).
\begin{enumerate}
    \item Prove that there is a torsion-free submodule $N \subseteq M$ such that $M = M_{\mathrm{tors}} \oplus N$, and that the isomorphism class of $N$ is uniquely determined by $M$.
    
    We define the \textbf{dimension} of $M$ and the \textbf{Steinitz class} of $M$ as the dimension of $N$ and the Steinitz class of $N$ respectively.
    
    \item Prove that the isomorphism class of $M$ is determined by its invariant factors, dimension and Steinitz class.
\end{enumerate}
\begin{solution}
    (1) Because $M/M_{\text{tros}}$ is a $R$-lattice. It is projective since $R$ is Dedekind domain.
    We have the exact sequence 
    \begin{align*}
        0\to M_{\text{tros}} \to M \to M_{\text{tros}}\to 0
    \end{align*} 
    split. We then have $M = M/M_{\text{tros}}\oplus M_{\text{tros}}$.

    (2) Let $M\simeq M'$. We have $M_{\text{tros}}\simeq M'_{\text{tros}}$ and $M/M_{\text{tros}}\simeq  M'/M'_{\text{tros}}$.
    The torsion-free part is uniquely determined by its Steinitz class and dimension.
    According to our textbook, the invariant factors are uniquely determined by torsion part $M_{\text{tros}}$.
    We just need to prove that for integral ideal $\mathfrak{a}_i$ and $\mathfrak{b}_j$ such that $\mathfrak{a}_i\supseteq \mathfrak{a}_{i+1}$ and $\mathfrak{b}_j\supseteq \mathfrak{b}_{j+1}$ 
    if 
    \begin{align*}
        \bigoplus_{i=1}^n R/\mathfrak{a}_i \simeq \bigoplus_{j=1}^m R/\mathfrak{b}_j
    \end{align*}
    we then have $\mathfrak{a}_i = \mathfrak{b}_i$.
    It is also trivial, since for every prime ideal $\mathfrak{q}$ of $R$, this two module must agree on those localization.

    If $\mathfrak{q}\supseteq \mathfrak{a}_n$ but $\mathfrak{q}\not\supseteq \mathfrak{b}_m$, then we have 
    \begin{align*}
        \left(\bigoplus_{j=1}^m R/\mathfrak{b}_j\right)_\mathfrak{q}
    \end{align*}
    is a zero $R_\mathfrak{q}$-module.
    However, 
    \begin{align*}
       \left(\bigoplus_{i=1}^n R/\mathfrak{a}_i\right)_\mathfrak{q}
    \end{align*}
    is a torsion $R_\mathfrak{q}$-module. It is impossible.
    Therefore, we must have $\mathfrak{a}_n$ and $\mathfrak{b}_m$ shares the same prime factors.
    Moreover, since $R_\mathfrak{q}$ is a DVR, and thus is a PID. According to the classification of modules over PID and 
    the fact that the prime elements on DVR are exactly uniformizers of $R_\mathfrak{q}$. It implies that 
    \begin{align*}
        (\mathfrak{a}_i)_\mathfrak{q}  = (\mathfrak{b}_i)_\mathfrak{q}.
    \end{align*}  
    and $n=m$.
    Our claim holds.
\end{solution}



\subsection*{Exercise 4.4}
Let $R$ be an integral domain with fraction field $F$. Let $L$ be a nonzero $R$-lattice in a finite dimensional $F$-vector space $V$.
Prove that $\dim L$ is equal to the largest positive integer $n$ such that $L$ contains a linearly independent system of $n$ vectors of $V$.

\begin{solution}
    Suppose that $\omega_1,\dots,\omega_n$ is such a system. Suppose that $\omega\in FL$
    can not be linearly represented by $\omega_i$ over $V$. Therefore, $\{\omega,\omega_1,\dots,\omega_n\}$ is linearly independent over $F$.
    Since $R$ is integral domain and $F$ is the fraction field of $R$, there exist an element $r\in R$ such that $\{r\omega,r\omega_1,\dots,r\omega_n\}$
    is included in $L$ and for each of those elements is not zero.
    Then we find a linearly independent system larger then previous one, which is a contradiction.
\end{solution}

\subsection*{Exercise 4.5}
Let $R$ be a Dedekind integral domain with fraction field $F$. Let $L$ be a nonzero $R$-lattice in a finite dimensional $F$-vector space $V$.

Prove that for any nonzero $x \in L$ the following are equivalent:
\begin{enumerate}
    \item $x$ is a primitive vector in $L$ (in the sense of Definition 1.2.17).
    \item the quotient $R$-module $L/Rx$ is torsion-free.
    \item $x$ is a maximal vector in $L$ (in the sense of (1.2.31)).
\end{enumerate}
\begin{solution}
    We first show that (1) implies (2). $x$ is primitive means the exact sequence 
    \begin{align*}
        0\to Rx \to L \to L/Rx\to 0 
    \end{align*}
    splits and $L\simeq Rx\oplus L/Rx$. If $L/Rx$ is not torsion-free,
    then there exists $r\in R$ and $y\in L/Rx$ such that $ry\in Rx$, the direct sum implies that $ry=0$. 
    then $L$ has a torsion element. Consequently, it cannot be a $R$-lattice.

    (2) also implies (1). Since $L/Rx$ as a $R$-lattice, it is projective $R$-module, and thus the exact sequence splits.

    Next, we show that (2) implies (3). Let $\mathfrak{a}$ be the coefficient ideal of $x$. Then we have $\mathfrak{a}x/Rx$ is torsion-free module.
    It also say that 
    \begin{align*}
        \mathfrak{a}/R \simeq R/\mathfrak{a}^{-1}
    \end{align*}
    is a torsion-free module. Actually, we choose an arbitrary element $d$ in determinant ideal $\det(\mathfrak{a})$, we then have $d\mathfrak{a}\subseteq R$.
    It holds if and only if $\mathfrak{a}=R$.
    Therefore, $x$ is maximal.

    We check (3) implies (2). Let $\mathfrak{a}$ be the coefficient ideal of $x$. Suppose that $L/\mathfrak{a}x$ has torsion element $y+\mathfrak{a}x$.
    It means that there exists $r$ such that $ry \in \mathfrak{a}x$ but $y\notin \mathfrak{a}x$.
    We can suppose that $ry = ax$ and $a/r x\notin \mathfrak{a}x$ means $a/r\notin \mathfrak{a}$. Then we have $y = a/r x\notin L$.
    It is a contradiction.

    Therefore, statements are all equivalent.
\end{solution}

\section*{Homework 5}

\subsection*{Exercise 5.1}
Let $I, J$ be fractional ideals of a Dedekind domain $R$. Prove that $I^{-1} \cap J^{-1} = (I + J)^{-1}$.
\begin{solution}
    Let $I^{-1} = \p_1^{-a_1}\cdots \p_n^{-a_n}$ and $J^{-1}=\p_1^{-b_1}\cdots \p_n^{-b_n}$.
    Then $I+J = \p_1^{c_1}\cdots \p_n^{c_n}$ where $c_i=\min\{a_i,b_i\}$.
    Let $v_i$ be the valuations correspond to $\p_i$ respectively.
    Then we have 
    \begin{align*}
        I^{-1}\cap J^{-1} =& \{r\in F\mid v_i(r)\ge \max\{-a_i,-b_i\}\text{ and $w(r)\ge 0$ for others valuation $w$}\}\\
        =& \{r\in F\mid v_i(r)\ge -\min\{-a_i,-b_i\}\text{ and $w(r)\ge 0$ for others valuation $w$}\} \\
        =& (I+J)^{-1}.
    \end{align*}
\end{solution}



\subsection*{Exercise 5.2}
Let $J$ be an integral ideal of a Dedekind domain $R$. Prove that for every fractional ideal $I$ of $R$, one has $I/IJ \cong R/J$ as $R$-modules.
\begin{solution}
    Let $\p$ be an arbitrary prime ideal. Then we have isomorphim of $R_\p/J_\p$-module
    \begin{align*}
        (I/IJ)_\p = I_\p/(I_\p J_\p) \simeq R_\p/J_\p.
    \end{align*}
    The isomorphism is given by $R_\p\simeq I_\p$. 
    Interpret $I/IJ$ and $R/J$ be the $R/J$-module. 
    According to Structure Theorem for Artinian Modules, we have 
    \begin{align*}
        I/IJ \cong R/J.
    \end{align*}
    Hence it is an isomorphism as $R$-modules.
\end{solution}


\subsection*{Exercise 5.3}
Prove that the existence statement in Thm.1.2.27 can be deduced from Thm.1.2.39.
\begin{solution}
    Let $\{x_1,\cdots,x_n\}$ be the generators of torsion $R$-module $M$. We consider the homomorphism 
    \begin{align*}
        R^{n} &\to M \\
        (r_i)_i&\mapsto \sum r_ix_i
    \end{align*}
    which is a epimorphism. Suppose its kernel is $N$. Clearly, $N$ is a torsion-free module with finite generators since $R$ is Noetherian. Therefore $N$ is a projective $R$-module.
    Consider the short exact sequence 
    \begin{align*}
        0\to N \to R^n \to M \to 0
    \end{align*}
    It implies that 
    \begin{align*}
        K^n \simeq N \otimes_R K.
    \end{align*}
    Since functor $\otimes_R K$ is exact and $M\otimes_R K$ is zero.

    Then there is a basis such that 
    \begin{align*}
        R^n = \a_1y_1 +\a_2y_2 +\cdots+ \a_ny_n\\
        N = \a_1\b_1y_1+\a_2\b_2y_2+\cdots + \a_n\b_ny_n
    \end{align*}
    and $\b_i\supseteq \b_{i+1}$.
    Then we obtain that 
    \begin{align*}
        R^n/N \simeq \bigoplus_i \a_i/(\a_i\b_i) \simeq  \bigoplus_i R/\b_i.
    \end{align*}
\end{solution}

\subsection*{Exercise 5.4}
 Prove the following facts about metabolic bilinear spaces over a field $F$:
\begin{enumerate}
    \item The hyperbolic spaces $m.\mathbb{H}$ are metabolic.
    \item For any symmetric bilinear spaces $V, V'$, if two of the three spaces $V, V'$ and $V \perp V'$ are metabolic, then so is the third. In particular, an orthogonal sum of metabolic spaces is metabolic.
    \item For any nonsingular symmetric bilinear space $(V, b)$, the space $(V, b) \perp (V, -b)$ is metabolic.
    \item Let $V$ be a 2-dimensional nonsingular symmetric space. Then $V$ is metabolic if and only if it is isotropic.
    \item Every 2-dimensional metabolic space is isomorphic to $\mathbb{H}$ or $\langle a, -a \rangle_{\mathrm{bil}}$ for some $a \neq 0$.
    \item If $\mathrm{char}(F) \neq 2$, then every metabolic bilinear space is hyperbolic.
\end{enumerate}
\begin{solution}
    (1) Clearly, $m.\mathbb{H}$ it is not singular of even dimension. Let $\{x_i,y_i\}$ be the basis of $i^{\text{th}}$ hyperbolic subspace of $m.\mathbb{H}$,
    then we let $V=\text{span}(x_i)$. It is easy to check that $V$ is totally isotropic and $\dim(V)=m$.

    (2) Suppose that $V$ and $V'$ are metabolic. Clearly, $V\perp V'$ is not singular of even dimension. Let $W$ and $W'$ be the totally isotropic subspace of $V$ and $V'$ respectively and 
    hance have half dimension. Then $W\perp W'$ also be the totally isotropic subspace of $V\perp V'$ with a half of dimension. Therefore, $V\perp V'$ is totally isotropic.

    Suppose that $V$ and $V\perp V'$ is metabolic. Then we have $\dim(V')= \dim(V\perp V')-\dim(V)$ is even.
    And $V'$ must be not singular. Otherwise, $V\perp V'$ will be singular. Let $W$ be the totally isotropic subspace of $V\perp V'$ with a half dimension.
    Then we have $W = (W\cap V)\perp (W\cap V')$. I claim that $\dim(W\cap V)\le \dim(V)/2$. If $\dim(W\cap V)> \dim(V)/2$, then we choose a basis of $W\cap V$ and extend it to be a basis of $V$.
    Let we denote it by $\{e_1,\dots,e_n,\dots,e_{2n}\}$. Write the Gram matrix over this basis and using Laplace expension, we deduce that $V$ is singular.
    It is a contradiction. Therefore, we have $\dim(V'\cap W)\ge \dim(W)-\dim(V)/2 = \dim(V')/2$.
    
    (3) Let $e_i$ be a standard basis of $(V,b)$. Then it is also a standard basis of $(V,-b)$. 
    Consider the basis $(e_i,e_i),(e_i,-e_i)$ of $V\oplus V$. Check subspace generated by $(e_i,e_i)$ is totally isomorphic subspace with half dimension of $(V,b)\perp (V,-b)$.

    (4) If it is isotropic, then it has Gram matrix $$ \begin{pmatrix} 0 & a \\ a & c \end{pmatrix} $$ with $a\in F^*$. 
    Consider base change matrix 
    
    It is a hyperbolic and thus is a metabolic.
    Conversely, if it is a metabolic, it has the Gram matrix $$ \begin{pmatrix} 0 & a \\ a & c \end{pmatrix} $$
    It implies that $V$ is isotropic.

    (5) If $c=0$, then $V$ is isomorphic to $\mathbb{H}$. If $c\neq 0$, let $\{e_1,e_2\}$ be the basis corresponding to this matrix.
    Then we consider the basis $\{e_2,e_2-a^{-1}ce_1\}$. It has the Gram matrix
    $$ \begin{pmatrix} c & 0 \\ 0 & -c \end{pmatrix} $$
    The statements holds.

    (6)  Let $U$ be the totally isotropic subspace of $V$ with half dimension.
    According to our textbook, metabolic space can be decomposit to $V_1\perp \dots \perp V_n$ where $V_i=Fx_i\otimes Fy_i$ is a subspace of dimension 2 where $(x_i,y_i)=1$ and $x_i\in U$.
    Therefore, $V_i$ is isotropic. 
    Suppose that $V$ is isotropic to the second case of (5). 
    We consider basis $\{x_i+y_i,x_i-y_i\}$. It has the Gram matrix
    $$ \begin{pmatrix} 0 & 2c \\ 2c & 0 \end{pmatrix} $$
    Since $\text{char}(F)\neq 2$, then we have $V_i$ is hyperbolic space and $V\simeq n.\mathbb{H}$.
\end{solution}


\section*{Homework 6}

In the following exercises, $F$ denotes a field of characteristic $\text{char}(F)\neq 2$. 
In principle, the following exercises can be done by using only results stated in \S 2.1.

\subsection*{Exercise 6.1}
Let $V$ be a nonsingular quadratic space over $F$. Prove that $V$ is isotropic if and only if $V$ represents $\mathbb{H}$.
\begin{solution}
    Suppose that $Q(x)=0$. Since $V$ is not singular, then there exists $y\in V$ such that $V = V'\perp(Fx\oplus Fy)$. Since characteristic is not equal 2, then we have isotropic subspace $Fx\oplus Fy$ is hyperbolic.
    Therefore we have representation $\varphi:\mathbb{H} \to Fx\oplus Fy$.
    Converse is trivial.
\end{solution}


\subsection*{Exercise 6.2}
Let $V$ be a quadratic space over $F$ and $\alpha \in F^\times$. Prove that the following statements are equivalent:
\begin{enumerate}
    \item[(i)] $V$ represents the 1-dimensional space $\langle \alpha \rangle$.
    \item[(ii)] There is a splitting $V = \langle \alpha \rangle \perp W$.
\end{enumerate}
Prove that if $V$ is nonsingular, then the above conditions are also equivalent to the following:
\begin{enumerate}
    \item[(iii)] The quadratic space $V \perp \langle -\alpha \rangle$ is isotropic.
\end{enumerate}

\begin{solution}
    We show that (1) implies (2). Since $\alpha\in F^*$, then $\langle \alpha \rangle$ is not singular.
    Then according to the textbook, we have $V\simeq \langle \alpha\rangle \perp W$.
    Converse is trivial.
    
    Suppose that $V$ is nonsingular. Then clearly (ii) implies (iii) by considering the $\langle\alpha\rangle\perp \langle-\alpha\rangle$.
    Conversely, $V\perp \langle -\alpha\rangle$ is also nonsingular. Then according to Excercise 6.1, we have $V\perp \langle -\alpha \rangle$ represents a hyperbolic space $W\simeq\mathbb{H}$ such that $V\perp \langle-\alpha\rangle = W\perp W^{\perp}$.
    Since $\text{char}(F)\neq 2$, $W$ represents $$ \begin{pmatrix} \alpha & 0 \\ 0 & -\alpha \end{pmatrix} $$
    Using Witt cancellation then we have $V\simeq \langle \alpha \rangle\perp W^{\perp}$.
    Therefore, $V$ represents the 1-dimensional space $\langle \alpha \rangle$.
\end{solution}

\subsection*{Exercise 6.3}
Prove that Witt cancellation holds for all quadratic spaces in characteristic $\neq 2$. More precisely, for any (possibly singular) quadratic spaces $V, V_1, V_2$ over $F$, prove that $V \perp V_1 \cong V \perp V_2$ if and only if $V_1 \cong V_2$.
\begin{solution}
The direction ($\Longleftarrow$) is trivial. We prove ($\Longrightarrow$).

Suppose there is an isomorphsim $V \perp V_1 \cong V \perp V_2$.

Since $\mathrm{char}(F) \neq 2$, the relation $B(x,x) = 2Q(x)$ guarantees that the radical of any quadratic space is a totally singular space. This ensures the existence of a valid orthogonal decomposition for any space $U$:
\[ U = \mathrm{rad}(U) \perp U_{nd} \]
where $U_{nd}$ is a strictly nonsingular subspace.

An isometry must preserve the radical. Applying this to our spaces gives:
\[ \mathrm{rad}(V) \perp \mathrm{rad}(V_1) \cong \mathrm{rad}(V) \perp \mathrm{rad}(V_2) \]
For totally singular spaces, isometry is entirely determined by dimension. Taking dimensions on both sides and canceling $\dim(\mathrm{rad}(V))$ yields $\dim(\mathrm{rad}(V_1)) = \dim(\mathrm{rad}(V_2))$. Thus:
\[ \mathrm{rad}(V_1) \cong \mathrm{rad}(V_2) \]

Furthermore, the isometry also preserves the nonsingular orthogonal complements:
\[ V_{nd} \perp (V_1)_{nd} \cong V_{nd} \perp (V_2)_{nd} \]
Because all components here are strictly nonsingular and $\mathrm{char}(F) \neq 2$, we are permitted to invoke the standard Witt's Cancellation Theorem (which relies on non-degenerate reflections) to cancel $V_{nd}$ from both sides. This yields:
\[ (V_1)_{nd} \cong (V_2)_{nd} \]

Finally, we reconstruct $V_1$ and $V_2$ by orthogonally summing their respective isometric parts:
\[ V_1 = \mathrm{rad}(V_1) \perp (V_1)_{nd} \cong \mathrm{rad}(V_2) \perp (V_2)_{nd} = V_2 \]
This completes the proof.
\end{solution}

\section*{Homework 7}

\subsection*{Exercise 7.1}
Let $(V,\langle \cdot,\cdot \rangle)$ be a nonsingular bilinear space and let $x_1,\cdots,x_n$ be a basis of $V$. Let $y_1,\cdots,y_n \in V$ be the dual basis of $x_1,\cdots,x_n$ with respect to the given bilinear form $\langle \cdot,\cdot \rangle$. Prove
\[
    \det(\langle x_i, x_j \rangle) \det(\langle y_i, y_j \rangle) = 1.
\]
\begin{solution}
    Suppose that $(x_i)=A(y_i)$. Then we have $x_i = \sum_{j}a_{ij}y_j$.
    Since $(y_i)$ is the dual basis of $(x_i)$, then we have 
    \begin{align*}
        \langle x_i,x_j\rangle = \langle \sum_{k}a_{ik}y_k,x_j\rangle = a_{ij}.
    \end{align*}
    Next, we consider $A^* (x_i) = \det(A)(y_i)$. We let $A^*=(a^*_{ij})$.
    \begin{align*}
        \det(A)\langle y_i,y_j\rangle = \langle y_i,\sum_{k}a^*_{jk}x_k\rangle = a_{ji}^*.
    \end{align*}
    We find that 
    \begin{align*}
        \det(\langle x_i, x_j \rangle) \det(\langle y_i, y_j \rangle) = \det(A)\det(A^*)\det(A^{-n})=1
    \end{align*}
\end{solution}


\subsection*{Exercise 7.2}
Consider the quadratic $\mathbb{Z}$-lattice $(L, Q)$ with $L = \mathbb{Z}^{\oplus 2}$ and $Q(x,y) = 2xy$.
Prove that $L$ cannot be decomposed as an orthogonal sum of two 1-dimensional lattices.

\begin{solution}
    Let $e_1,e_2$ denotes $(1,0)$ and $(0,1)$ respectively. Then we have 
    \begin{align*}
        \langle ae_1+be_2,c e_1+de_2 \rangle =2(bc+ad).
    \end{align*}
    If $ae_1+be_2,c e_1+de_2$ are orthogonal, then we must have $bc+ad=0$.
    Hence, $ae_1+be_2,c e_1+de_2$ must be the generators of $L$, which means the relative determinant must be $\pm 1$, which is $ad-bc=\pm 1$.
    Then we have 
    \begin{align*}
        2bc = \mp 1.
    \end{align*}
    It is impossible, since we must have $b,c\in\mathbb{Z}$.
\end{solution}

\section*{Homework 9}

\subsection*{Exercise 9.1}
 Let $(V, Q)$ be a quadratic space over a field $F$ of characteristic $\neq 2$.
\begin{enumerate}
    \item Check that for any anisotropic vector $x \in V$, the symmetry $\tau_x$ defined in (2.3.1) is an isometry of $(V, Q)$.
    \item Assume that $(V, Q)$ is nonsingular. Without using the Cartan--Dieudonné theorem, prove that $\det(\tau_x) = -1$ for any anisotropic vector $x \in V$ and that $\det(\sigma) = \pm 1$ for every $\sigma \in \mathrm{O}(V)$.
\end{enumerate}
\begin{solution}
    (1)
    We let $B(x,y)=Q(x+y)-Q(x)-Q(y)$.
    Clearly, $\tau_x$ is invertible since $\tau_x^2 = \text{id}$. 
    On one hand, we have
    \begin{align*}
        B\left(y,\frac{-B(x,y)}{Q(x)}x\right) = Q(\tau_x(y)) - Q(y) - \frac{B^2(x,y)}{Q^2(x)}Q(x).
    \end{align*} 
    On the other hand we have 
    \begin{align*}
        B\left(y,\frac{-B(x,y)}{Q(x)}x\right) = -\frac{B^2(x,y)}{Q(x)}.
    \end{align*}
    Then we obtain that $Q(\tau_x(y))=Q(y)$

    (2) Since $V$ is not singular, there exists an orthogonal decomposition $V = \langle x\rangle+ \langle x\rangle^\perp$. Therefore, the linear transformation $\tau_x$ gives
    matrix of diagonal form $\text{diag}(-1,1,\dots,1)$. Then we obtain that $\det(\tau_x)=-1$.

    To prove that $\det(\sigma)=\pm 1$. Take $\{e_i\}$ be the orthogonal basis whose Gram matrix is given by $G$. Let $M$ be the matrix representing the linear transformation $\sigma$.
    Then we have 
    \begin{align*}
        M^TGM = G
    \end{align*}
    and hence $\det(G)\det^2(M)=\det(G)$. We then find that $\det(M)=\pm 1$.
\end{solution}

\subsection*{Exercise 9.2}
 Let $M$ be a subgroup of the additive group $\mathbb{R}$ endowed with the usual Euclidean topology. Prove that the following assertions are equivalent:
\begin{itemize}
    \item[(i)] $M$ is not dense in $\mathbb{R}$.
    \item[(ii)] There exists a real number $\varepsilon > 0$ such that $M \cap \left]-\varepsilon, \varepsilon\right[ = \{0\}$.\\
    (Here $\left]-\varepsilon, \varepsilon\right[$ denotes the open interval $\{x \in \mathbb{R} \mid -\varepsilon < x < \varepsilon\}$.)
    \item[(iii)] $M$ is discrete in $\mathbb{R}$, i.e., the subspace topology on $M$ is the discrete topology.
    \item[(iv)] $M = r\mathbb{Z}$ for some $r \in \mathbb{R}$.
\end{itemize}
\begin{solution}
    We first show that (1) implies (2). Suppose that $M$ is dense on $0$. Then for arbitrary $x\in M$ and $y\in \mathbb{R}$, there exists $\varepsilon\in M$ such that $|x-y|<\varepsilon$.
    If $x>y$, it implies that $x>x-\varepsilon>y$ and if $x<y$, it implies that $x<x+\varepsilon<y$. In any case, we find an element $x\pm\varepsilon$ is closer than $x$ to $y$.
    Equivalently we write $d(M,y)=0$. It contradicts to (1).

    Clearly, (2) also implies (1) since $M$ is not dense on zero.

    Hence, (2) implies (3). Because $M$ must be nowhere dense. If $M$ is dense on $x$, for arbitrary $\varepsilon>0$ such that sufficiently small, then there exists an $x+\eta_1,x+\eta_2\in M$ such that $|\eta_1-\eta_2|<\varepsilon$.
    It follows that $M$ is dense on zero, which contradicts to (2).

    (3) implies (2) also holds by defenition.

    (4) implies (2) is trivial.

    Finally, we check that (3) implies (4). We consider the set of distence between elements in $M$, denoted by 
    \begin{align*}
        D=\{d(x,y)\mid x,y\in M\}.
    \end{align*}
    Since $M$ is discrete, we can find a miminal value $\delta$ of $D$. Clearly, we have $\delta\mathbb Z\subseteq M$. 
    It there exists an element $y\in M$ is not contained in $\delta \mathbb{Z}$, then we find that $d(y,\delta \mathbb{Z})<\delta$. It contradicts to the minimality of $\delta$.
    Therefore, we then have $M=\delta \mathbb{Z}$.
\end{solution}


\subsection*{Exercise 9.3}
 Let $K$ be a field and let $v$ be an exponential valuation on $K$ with valuation ring $A$. Let $(\hat{K}, \hat{v})$ be the completion of the valued field $(K, v)$.
\begin{enumerate}
    \item Let $\hat{A}$ denote the valuation ring of $\hat{v}$ in $\hat{K}$. Prove that $\hat{A}$ equals the topological closure of $A$ in $\hat{K}$.
    \item Show that $\hat{v}(\hat{K}^\times) = v(K^\times)$ and that the natural map $A \to \hat{A}$ induces an isomorphism on their residue fields. In particular, if $v$ is discrete (resp. discrete and normalized), then so is $\hat{v}$.
\end{enumerate}
\begin{solution}
    \textbf{Maybe we should show stability first.}

    (1) Let $\bar{A}$ denote the topological closure of $A$. Let $||$ denote the absolute valuate determined by $v$.
    Let $(x_i)$ be a Cauchy sequence over $K$. Suppose that $(x_i)\in \hat{A}$, it means that $\lim_{n\to\infty}|x_n|\le 1$. 
    Then it is to said that there exists an integer $N>0$ for any $n>N$ we have $x_n\in A$ since stability.  The we consider constant sequences $(x_{N+k})_i$.
    For arbitrary $\varepsilon>0$, there exists an integer $N$ such that for $n>N$ we have 
    \begin{align*}
        |x_n-x_{N+k}|<\varepsilon
    \end{align*}
    for arbitrary $k>0$. It equivalently say given any $\varepsilon>0$, we find an element $(x_{N+k})_i\in A$ such that $\|(x_{N+k})_i - (x_i)_i\|<\varepsilon$ where $\|\|$ denotes the absolute valuate induced by $\hat{v}$.
    Thus we have $\hat{A}\subseteq \bar{A}$.

    Conversely, since $\hat{A}$ is closed in $\hat{K}$ and it contains $A$, we then have $\bar{A}\subseteq \hat{A}$.
    Therefore, we obtain that $\bar{A}=  \hat{A}$

    (2) We just need to show that any Cauchy sequence $(x_i)$ admits a stable sequence of valuation $v(x_i)$ if $(x_i)$ is not the zero over $\hat{K}$.
    To see that we use the strongly triangular inequality. Let $a = \lim_{n\to \infty}|x_n|$. There exists $n>0$ such that for any $k$ we have 
    \begin{align*}
        |x_n-x_{n+k}|<a/2.
    \end{align*}
    We can also choose sufficiently large of $n$ such that
    \begin{align*}
        ||x_{n+k}|-a|<a/2
    \end{align*} 
    In this sense we have $|x_{n+k}|>a/2>|x_n-x_{n+k}|$, then 
    \begin{align*}
        |x_{n}| = |x_{n}-x_{n+k}+x_{n+k}| = |x_{n+k}|.
    \end{align*}
    Therefore we show that $|x_n|$ is stable and thus $\hat{v}(\hat{K}^\times) = v(K^\times)$. 

    Let $\mathfrak{m}$ and $\hat{\mathfrak{m}}$ denote the maximal prime ideal of $A$ and $\hat{A}$ respectively. 
    Consider the homomorphism 
    \begin{align*}
        A \to \hat{A}/\hat{\mathfrak{m}}.
    \end{align*}
    It is epimorphic. For arbitrary Cauchy sequence $(x_i)$ such that $\lim_{n\to\infty}|x_n| = 1$, we suppose it is stable at $k^{\text{th}}$.
    Hence we can suppose that $|x_{k+n}-x_k|<1/2$, then we find that $(x_n-x_k)_n\in \hat{\mathfrak{m}}$ and therefore $(x_n)$ have preimage $(x_k)_n$.
    Its kernal is $A\cap \hat{\mathfrak{m}} = \mathfrak{m}$.

    Our claim holds.
\end{solution}



\subsection*{Exercise 9.4}
Let $R$ be a proper subring of a field $K$ such that for every $x \in K^\times$, either $x \in R$ or $x^{-1} \in R$. (We say that $R$ is a \textbf{valuation ring} in $K$.)
\begin{enumerate}
    \item Prove that $R$ is integrally closed.
    \item Let $v$ be a nontrivial exponential valuation on a field $K$. Show that the ring $\mathcal{O}_v = \{x \in K \mid v(x) \ge 0\}$ is a valuation ring in the above sense, i.e., for all $x \in K^\times$, either $x \in \mathcal{O}_v$ or $x^{-1} \in \mathcal{O}_v$. Deduce that $\mathcal{O}_v$ is integrally closed.
\end{enumerate}
\begin{solution}
    (1) Suppose that it is not integrally closed. Equivalently, there exits an integral element $x\in K\setminus R$. Suppose that 
    \begin{align*}
        x^n + a_{n-1}x^{n-1}+\dots+a_0 = 0
    \end{align*}
    Then we have 
    \begin{align*}
        a_0x^{-n} + a_1x^{-n+1}+\dots+1 = 0
    \end{align*}
    Hence we have 
    \begin{align*}
        x = -(a_0x^{-n+1}+a_1x^{-n+2}+\dots+a_{n-1}).
    \end{align*}
    It implies that $x\in R$.

    (2) If $x\notin \o_v$, then we have $v(x)<0$. Therefore, we must have $v(x^{-1})>0$ and $x^{-1}\in\o_v$. The intergrality of $\o_v$ follows.
\end{solution}


\subsection*{Exercise 9.5}
Let $|\cdot|_1, \cdots, |\cdot|_n$ be pairwise inequivalent non-archimedean absolute values of a field $K$. Prove that for all $a_1, \cdots, a_n \in K^\times$ there exists an element $x \in K^\times$ such that $|x|_i = |a_i|_i$ for each $i = 1, \cdots, n$.
\begin{solution}
    We use Weak Approximation Theorem. Let $\varepsilon<\min_i\{|a_i|_i\}$, then there exist an element $x\in K$ such that 
    \begin{align*}
        |x-a_i|_i<\varepsilon.
    \end{align*}
    Next, we consider 
    \begin{align*}
        |x|_i = |x-a_i+a_i|_i = |a_i|_i
    \end{align*}
    since $|x-a_i|_i<|a_i|_i$.
\end{solution}



\subsection*{Exercise 9.6}
Let $(K, |\cdot|)$ be a complete non-archimedean valued field. Let $(a_n)_{n \in \mathbb{N}}$ be a sequence in $K$. Prove that the series $\sum_{n=0}^\infty a_n$ converges in $K$ if and only if $\lim_{n \to \infty} a_n = 0$ in $K$.
\begin{solution}
    Suppose that $\sum a_n$ converges. Then $|a_{n+k+1}-a_{n+k}|<\varepsilon$ therefore $|a_{n+k}-a_n|=|a_{n+k}-a_{n+k-1}+\dots+a_{n+1}-a_n|<\varepsilon$ which implies that $a_n$ converges.

    Conversely, $|a_n+\dots+a_{n+k}|\le \max_{k}\{|a_{n+k}|\}<\varepsilon$. $\sum a_n$ converges.
\end{solution}


\subsection*{Exercise 9.7}
Let $(K, |\cdot|)$ be a henselian (non-archimedean valued) field.
\begin{enumerate}
    \item Let $f(T) = a_2T^2 + a_1T + a_0$ be a polynomial of degree $2$ over $K$.\\
    Prove that if $f$ is irreducible over $K$, then $|a_1|^2 \le |a_0| \cdot |a_2|$.
    \item Let $(V, Q)$ be an anisotropic quadratic space over $K$. Prove that
    \[ |\langle x, y \rangle_Q|^2 \le |Q(x)| \cdot |Q(y)| \quad \text{for all } x, y \in V. \]
    (Here $\langle x, y \rangle_Q := Q(x+y) - Q(x) - Q(y)$.)
\end{enumerate}
\begin{solution}
    (1) Suppose that $|a_1|^2>|a_0a_2|$, it equivalently say that $a_0a_1^{-2}a_2\in \p$ where $R$ is the valuation ring induced by absolute value.
    Without loss of generality, let $f(T)$ be the primitive polynomial. 
    Let $T = a_0a_1^{-1}X$, we then obtain that 
    $$f(X) = a_0 \left( \frac{a_2 a_0}{a_1^2} X^2 + X + 1 \right)$$
    Let $$g(X) =\frac{a_2 a_0}{a_1^2} X^2 + X + 1 $$
    which is the polynomial with coefficients in $R$. Consider that
    \begin{align*}
        \bar{g}(X) = X+1 \bmod \p
    \end{align*}
    Then $\bar{g}(X)$ has a simple root since $\bar{g}'(X)= 1$. Then by Hensel's Lemma, it can be lifted to $R$.
    However, $g(X)$ is irreducible and thus it can not have a root.

    (2) Let $a_2= Q(x),a_0=Q(y)$ and $a_1= \langle x, y \rangle_Q$. 
    Then $Q(Tx+y) = a_2T^2+a_1T+a_0$. Since $V$ is anisotropic, then $Q(Tx+y)$ is irreducible over $K$.
    Then according to (1) is what we want.
\end{solution}


\section*{Homework 10}

\subsection*{Exercise 10.1}
Let $K = \mathbb{R}$ or $K = \mathbb{C}$, and let $|\cdot|$ be the usual absolute value on $K$. Let $s$ be a positive real number. Prove that $|\cdot|^s$ is an absolute value on $K$ if and only if $0 < s \leq 1$.

\begin{solution}
    AV1 and AV2 is quite clear. We just show AV3.
    Suppose that $|\cdot|^s$ is an absolute value, we consider that 
    \begin{align*}
        2\ge 2^s
    \end{align*}
    It menas that $0<s\le 1$. Suppose that $0<s\le 1$, we just need to prove that 
    \begin{align*}
        1+x^s\ge (1+x)^s
    \end{align*}
    for $0<x< 1$. Consider function $g(x)=1+x^s-(1+x)^s$ with 
    \begin{align*}
        g'(x) = sx^{s-1} - s(1+x)^{s-1}>0
    \end{align*}
    Therefore $g(x)>g(0)=0$. Our claim holds.
\end{solution}

\section*{Homework 11}

\subsection*{Exercise A.11.1}

Let $\mathbb{F}$ be a finite field.

\begin{enumerate}
    \item Prove that every quadratic space of dimension $\geq 3$ over $\mathbb{F}$ is isotropic.

    You should not use Proposition 4.2.7 for this proof.

    \item Prove that $\mathbb{F}$ has separable $2$-dimension $\leq 1$; cf. Definition 4.2.5.
\end{enumerate}

The above statements hold even if $\operatorname{char}(\mathbb{F})=2$. If you like, you can only treat the case $\operatorname{char}(\mathbb{F})\neq 2$.

\begin{solution}
    We assume that $\operatorname{char}(\mathbb{F})\neq 2$ and $q=|\mathbb{F}|$.

    (1)First we prove a simple counting lemma. Let $a,b\in \mathbb{F}^{\times}$. Then the binary quadratic form $ax^2+by^2$ represents every element of $\mathbb{F}$. Indeed, fix $t\in \mathbb{F}$ and consider \[ A=\{ax^2:x\in \mathbb{F}\},\qquad B=\{t-by^2:y\in \mathbb{F}\}. \] Since $\mathbb{F}$ has odd cardinality, the set of squares in $\mathbb{F}$, including $0$, has cardinality $(q+1)/2$. Hence $|A|=|B|=(q+1)/2$. Since $|A|+|B|=q+1>q$, the two subsets $A$ and $B$ must intersect. Thus there exist $x,y\in \mathbb{F}$ such that $ax^2=t-by^2$, or equivalently $ax^2+by^2=t$.
    
    Now let $(V,Q)$ be a quadratic space over $\mathbb{F}$ with $\dim V\geq 3$. Since $\operatorname{char}(\mathbb{F})\neq 2$, the quadratic form $Q$ is diagonalizable. Thus, for a suitable basis, we may write \[ Q\simeq \langle a_1,\dots,a_r,0,\dots,0\rangle \] with $a_i\in \mathbb{F}^{\times}$. If a zero coefficient occurs, then the corresponding basis vector is a nonzero isotropic vector, so $Q$ is isotropic. Hence we may assume that no zero coefficient occurs. Since $\dim V\geq 3$, the restriction of $Q$ to the subspace spanned by the first three basis vectors has the form \[ ax^2+by^2+cz^2 \] with $a,b,c\in \mathbb{F}^{\times}$. By the counting fact above, the binary form $ax^2+by^2$ represents $-c$. Therefore there exist $x,y\in \mathbb{F}$ such that $ax^2+by^2=-c$. It follows that \[ Q(x,y,1,0,\dots,0)=ax^2+by^2+c=0. \] The vector $(x,y,1,0,\dots,0)$ is nonzero, so $Q$ is isotropic. This proves that every quadratic space of dimension at least $3$ over $\mathbb{F}$ is isotropic. 
    
    
    (2)It remains to prove that $\operatorname{sd}_2(\mathbb F)\leq 1$. Let
$L/\mathbb F$ be a finite separable extension, and let $E/L$ be a quadratic
separable extension. Since $\mathbb F$ is finite, $L$ is again a finite field.
Assume $\operatorname{char}(\mathbb F)\neq 2$. Then $E=L(\sqrt d)$ for some
$d\in L^\times\setminus L^{\times 2}$.

We show that $N_{E/L}:E^\times\to L^\times$ is surjective. Let
$a\in L^\times$. Consider the quadratic form
\[
q(x,y,z)=x^2-dy^2-az^2
\]
over $L$. By the first part, this three-dimensional quadratic form is
isotropic. Hence there exists a nonzero vector $(x,y,z)\in L^3$ such that
\[
x^2-dy^2=az^2.
\]
We claim that $z\neq 0$. If $z=0$, then $x^2=dy^2$. Since
$d\notin L^{\times 2}$, this forces $x=y=0$, contradicting that
$(x,y,z)$ is nonzero. Thus $z\neq 0$, and after dividing by $z^2$ we get
\[
\left(\frac{x}{z}\right)^2
-
d\left(\frac{y}{z}\right)^2
=a.
\]
But the left-hand side is
\[
N_{E/L}\left(\frac{x}{z}+\frac{y}{z}\sqrt d\right).
\]
Therefore $a$ lies in the image of $N_{E/L}$. Since $a\in L^\times$ was
arbitrary, the norm map is surjective. Hence, by Definition 4.2.5,
$\operatorname{sd}_2(\mathbb F)\leq 1$.
\end{solution}

\subsection*{Exercise A.11.2}

Let $K$ be a field of characteristic $\neq 2$, and let $(V,\varphi)$, $(U,\psi)$ be nonsingular quadratic spaces of positive dimension over $K$.

\begin{enumerate}
    \item Prove that $(V,\varphi)$ is hyperbolic if and only if $\dim V$ is even and there exists a totally isotropic subspace $W\subseteq V$ such that $\dim V=2\dim W$.

    \item Prove that the quadratic space $(U,\psi)\perp (U,-\psi)$ is hyperbolic.

    \item Prove that $(V,\varphi)\cong (U,\psi)$ if and only if $(V,\varphi)\perp (U,-\psi)$ is hyperbolic.
\end{enumerate}
\begin{solution}
    (1) Let $B(x,y)=\varphi(x+y)-\varphi(x)-\varphi(y)$. Since $(V,\varphi)$ is nonsingular and $\operatorname{char}K\neq 2$, $B$ is nondegenerate. Suppose first that $(V,\varphi)$ is hyperbolic. Then $V=H_1\perp\cdots\perp H_n$, where each $H_i$ is a hyperbolic plane. Choose a hyperbolic basis $e_i,f_i$ of $H_i$. Then $W=\langle e_1,\dots,e_n\rangle$ is totally isotropic, and $\dim V=2n=2\dim W$. 
    
    Conversely, suppose that $W\subseteq V$ is totally isotropic and $\dim V=2\dim W$. Since $\varphi|_W=0$ and $\operatorname{char}K\neq 2$, one has $B|_{W}=0$, hence $W\subseteq W^\perp$. By nondegeneracy, $\dim W^\perp=\dim V-\dim W=\dim W$, so $W=W^\perp$. We prove that $V$ is hyperbolic by induction on $\dim W$. Choose $0\neq e\in W$. Since $B$ is nondegenerate, there exists $y\in V$ with $B(e,y)\neq 0$. For any $a\in K$, $\varphi(y+ae)=\varphi(y)+aB(y,e)$, because $\varphi(e)=0$. Taking $a=-\varphi(y)/B(y,e)$, we get $f=y+ae$ with $\varphi(f)=0$ and $B(e,f)\neq 0$. After scaling $f$, assume $B(e,f)=1$. Thus $H=\langle e,f\rangle$ is a hyperbolic plane. Since $H$ is nonsingular, $V=H\perp H^\perp$. Set $W'=W\cap H^\perp=\{w\in W:B(w,f)=0\}$. The map $W\to K$, $w\mapsto B(w,f)$, is nonzero because $B(e,f)=1$; hence $\dim W'=\dim W-1$. Also $W'$ is totally isotropic and $\dim H^\perp=\dim V-2=2(\dim W-1)=2\dim W'$. By induction, $H^\perp$ is hyperbolic. Therefore $V=H\perp H^\perp$ is hyperbolic.

    (2)Consider $(U,\psi)\perp (U,-\psi)$ on $U\oplus U$. Its quadratic form is \[ q(u_1,u_2)=\psi(u_1)-\psi(u_2). \] Let \[ \Delta=\{(u,u):u\in U\}\subseteq U\oplus U. \] Then $q(u,u)=0$ for all $u\in U$. Moreover, for the associated bilinear form, \[ B_q((u,u),(v,v))=B_\psi(u,v)-B_\psi(u,v)=0, \] so $\Delta$ is totally isotropic. Also \[ \dim(U\oplus U)=2\dim U=2\dim \Delta. \] By part $(1)$, $(U,\psi)\perp (U,-\psi)$ is hyperbolic.

(3) Assume first that \((V,\varphi)\cong (U,\psi)\). Then
\[
(V,\varphi)\perp (U,-\psi)
\cong
(U,\psi)\perp (U,-\psi),
\]
which is hyperbolic by part (2).

Conversely, suppose that
\[
(V,\varphi)\perp (U,-\psi)
\]
is hyperbolic. Since this space is nonsingular and hyperbolic, its anisotropic
part is zero.

Write the Witt decompositions
\[
(V,\varphi)\cong H^r\perp A,
\qquad
(U,\psi)\cong H^s\perp B,
\]
where \(H\) is a hyperbolic plane, and \(A,B\) are anisotropic. Then
\[
(U,-\psi)\cong H^s\perp (-B),
\]
because the negative of a hyperbolic plane is again hyperbolic. Therefore
\[
(V,\varphi)\perp (U,-\psi)
\cong
H^{r+s}\perp A\perp (-B).
\]
Since this space is hyperbolic, its anisotropic part is zero. Hence
\(A\perp (-B)\) is hyperbolic.

We now show that this implies \(A\cong B\). Since
\[
A\perp (-B)
\]
is hyperbolic, by part (1) it has a totally isotropic subspace of dimension
\[
\frac{1}{2}\dim(A\perp (-B)).
\]
Equivalently, by Witt decomposition, the anisotropic part of
\(A\perp (-B)\) is zero. If we decompose \(A\perp (-B)\), the only possible
anisotropic contribution comes from \(A\) and \(B\). By the uniqueness of the
anisotropic part in the Witt decomposition, \(A\) and \(B\) must be isometric.
Indeed, if
\[
A\perp (-B)\cong H^m,
\]
then
\[
A\perp (-B)\perp B\cong H^m\perp B.
\]
By part (2), \((-B)\perp B\) is hyperbolic. Hence the left-hand side is
isometric to \(A\) plus a hyperbolic space. Cancelling the common hyperbolic
part in the Witt decomposition gives
\[
A\cong B.
\]

Now use the assumption \(\dim V=\dim U\). From
\[
\dim V=2r+\dim A,
\qquad
\dim U=2s+\dim B,
\]
and \(A\cong B\), we get
\[
2r+\dim A=2s+\dim A.
\]
Hence \(r=s\). Therefore
\[
(V,\varphi)\cong H^r\perp A
\cong
H^s\perp B
\cong
(U,\psi).
\]
This proves
\[
(V,\varphi)\cong (U,\psi)
\quad\Longleftrightarrow\quad
(V,\varphi)\perp (U,-\psi)\text{ is hyperbolic}.
\]
\end{solution}


\subsection*{Exercise A.11.3}

Let $K$ be a field of characteristic $\neq 2$.

\begin{enumerate}
    \item Prove that for any two nonsingular quadratic spaces $(U,\psi)$, $(V,\varphi)$ over $K$, one has
    \[
        \operatorname{disc}\bigl((U,\psi)\perp (V,\varphi)\bigr)
        =
        (-1)^{\dim U \dim V}
        \operatorname{disc}(U,\psi)\operatorname{disc}(V,\varphi)
        \in K^{\times}/K^{\times 2}.
    \]

    \item Let $(V,\varphi)$ be a nonsingular quadratic space of dimension $2$ over $K$. Prove that $(V,\varphi)$ is hyperbolic if and only if $\operatorname{disc}(V,\varphi)=1\in K^{\times}/K^{\times 2}$.
\end{enumerate}

\begin{solution}
    We use the convention that, if $(V,\varphi)$ has dimension $n$ and Gram matrix $A$, then \[ \operatorname{disc}(V,\varphi) = (-1)^{n(n-1)/2}\det(A) \in K^\times/K^{\times 2}. \] For the first statement, let $\dim U=m$ and $\dim V=n$. Choose Gram matrices $A$ and $B$ for $(U,\psi)$ and $(V,\varphi)$ respectively. Then the Gram matrix of $(U,\psi)\perp (V,\varphi)$ is \[ \begin{pmatrix} A & 0\\ 0 & B \end{pmatrix}. \] Hence \[ \operatorname{disc}\bigl((U,\psi)\perp (V,\varphi)\bigr) = (-1)^{(m+n)(m+n-1)/2}\det(A)\det(B). \] On the other hand, \[ \operatorname{disc}(U,\psi)\operatorname{disc}(V,\varphi) = (-1)^{m(m-1)/2+n(n-1)/2}\det(A)\det(B). \] Since \[ \frac{(m+n)(m+n-1)}{2} - \frac{m(m-1)}{2} - \frac{n(n-1)}{2} = mn, \] we get \[ \operatorname{disc}\bigl((U,\psi)\perp (V,\varphi)\bigr) = (-1)^{mn} \operatorname{disc}(U,\psi)\operatorname{disc}(V,\varphi). \] For the second statement, since $\operatorname{char}K\neq 2$, we may write \[ (V,\varphi)\cong \langle a,b\rangle \] with $a,b\in K^\times$. Then \[ \operatorname{disc}(V,\varphi)=-ab\in K^\times/K^{\times 2}. \] If $(V,\varphi)$ is hyperbolic, then $(V,\varphi)\cong \langle 1,-1\rangle$, so \[ \operatorname{disc}(V,\varphi)=-(1)(-1)=1. \] Conversely, suppose that $\operatorname{disc}(V,\varphi)=1$. Then $-ab\in K^{\times 2}$, so there exists $c\in K^\times$ such that $-ab=c^2$. Thus \[ b=-\frac{c^2}{a}=-a\left(\frac{c}{a}\right)^2. \] Therefore \[ \langle a,b\rangle = \left\langle a,-\frac{c^2}{a}\right\rangle \cong \langle a,-a\rangle. \] The vector $(1,1)$ is isotropic for $\langle a,-a\rangle$. Since $\langle a,-a\rangle$ is nonsingular of dimension $2$, it is hyperbolic. Hence $(V,\varphi)$ is hyperbolic.
\end{solution}

\subsection*{Exercise A.11.4}

Let $F$ be a henselian discrete valuation field with residue field $k$. Assume $\operatorname{char}(k)\neq 2$. Let $u_1,\cdots,u_r,w_1,\cdots,w_r\in \mathcal{O}_F^{\times}$, and put
\[
    \varphi=\langle u_1,\cdots,u_r\rangle,
    \qquad
    \psi=\langle w_1,\cdots,w_r\rangle.
\]

\begin{enumerate}
    \item Prove that $\varphi\perp -\psi$ is hyperbolic over $F$ if and only if $\overline{\varphi}\perp -\overline{\psi}$ is hyperbolic over $k$.

    \item Prove that $\varphi$ and $\psi$ are isomorphic over $F$ if and only if $\overline{\varphi}$ and $\overline{\psi}$ are isomorphic over $k$.
\end{enumerate}
\begin{solution}
    We first deduce the following consequence of Springer's theorem. Let
\(q=\langle a_1,\dots,a_n\rangle\) be a unit form over \(F\), i.e.
\(a_i\in\mathcal O_F^\times\). Then
\[
q \text{ is hyperbolic over }F
\quad\Longleftrightarrow\quad
\bar q \text{ is hyperbolic over }k.
\]

The $\Rightarrow$ is clear. For $\Leftarrow$, let $(V,q) = m\mathbb{H}\perp W$, then we have $(\bar{V},\bar{q})= m\mathbb{H}\perp \bar{W}$. If $\bar{q}$ is hyperbolic, then 
it menas that $\bar{W}$ is isotropic. However, the springer theroem implies that $\bar{W}$ must be anisotropic, it is a contradiction.


Now apply this to
\[
\varphi\perp(-\psi)
=
\langle u_1,\dots,u_r,-w_1,\dots,-w_r\rangle .
\]
This is a unit form, and its reduction is
\[
\bar\varphi\perp(-\bar\psi)
=
\langle \bar u_1,\dots,\bar u_r,-\bar w_1,\dots,-\bar w_r\rangle .
\]
Hence
\[
\varphi\perp(-\psi)\text{ is hyperbolic over }F
\quad\Longleftrightarrow\quad
\bar\varphi\perp(-\bar\psi)\text{ is hyperbolic over }k.
\]
This proves part (1).

For part (2), since \(\dim\varphi=\dim\psi\), Exercise A.11.2(3) gives
\[
\varphi\cong\psi
\quad\Longleftrightarrow\quad
\varphi\perp(-\psi)\text{ is hyperbolic over }F.
\]
By part (1), this is equivalent to
\[
\bar\varphi\perp(-\bar\psi)\text{ is hyperbolic over }k.
\]
Again by Exercise A.11.2(3), this is equivalent to
\[
\bar\varphi\cong\bar\psi.
\]
Therefore
\[
\varphi\cong\psi
\quad\Longleftrightarrow\quad
\bar\varphi\cong\bar\psi .
\]
\end{solution}



\subsection*{Exercise A.11.5}

Let $\mathbb{F}$ be a finite field of characteristic $\neq 2$. Prove that up to isomorphism there is one and only one anisotropic binary quadratic space over $\mathbb{F}$.

\begin{solution}

Let $\delta\in \mathbb F^\times$ be a nonsquare. We claim that
\[
\langle 1,-\delta\rangle
\]
is the unique anisotropic binary quadratic space over $\mathbb F$ up to
isomorphism.

First, $\langle 1,-\delta\rangle$ is anisotropic. Indeed, if
$x^2-\delta y^2=0$ and $y\neq 0$, then $\delta=(x/y)^2$, contradicting that
$\delta$ is a nonsquare. Hence $y=0$, and then $x=0$.

Now let $(V,q)$ be any anisotropic binary quadratic space over $\mathbb F$.
Since $\operatorname{char}\mathbb F\neq 2$, we may diagonalize
\[
(V,q)\cong \langle a,b\rangle
\]
with $a,b\in \mathbb F^\times$. By Exercise A.11.3(2), a nonsingular binary
form is hyperbolic if and only if its discriminant is $1$. Since $(V,q)$ is
anisotropic, it is not hyperbolic, so
\[
\operatorname{disc}(V,q)=-ab
\]
is the nontrivial square class in $\mathbb F^\times/\mathbb F^{\times 2}$.
Thus $-ab=\delta$ in $\mathbb F^\times/\mathbb F^{\times 2}$, and hence
\[
\langle a,b\rangle \cong \left\langle a,-\frac{\delta}{a}\right\rangle .
\]

It remains to compare this form with $\langle 1,-\delta\rangle$. Since
$\mathbb F(\sqrt{\delta})/\mathbb F$ is a quadratic extension of finite fields,
its norm map is surjective. Hence the norm form
\[
x^2-\delta y^2
\]
represents $a$. Therefore
\[
\langle 1,-\delta\rangle \cong \langle a,c\rangle
\]
for some $c\in\mathbb F^\times$. Comparing discriminants gives
\[
-ac=\delta \quad \text{in } \mathbb F^\times/\mathbb F^{\times 2},
\]
so $c$ differs from $-\delta/a$ by a square factor. Hence
\[
\langle a,c\rangle \cong \left\langle a,-\frac{\delta}{a}\right\rangle .
\]
Therefore
\[
(V,q)\cong \langle 1,-\delta\rangle .
\]
Thus, up to isomorphism, there is exactly one anisotropic binary quadratic
space over $\mathbb F$.
\end{solution}

\subsection*{Exercise A.11.6}

Let $F$ be a henselian discrete valuation field of characteristic $\neq 2$. Prove that $F^{\times 2}$ is an open and closed subgroup of $F^{\times}$. Here of course, the topology on $F^{\times}$ is the subspace topology induced from $F$.
\begin{solution}
Let \(\pi\) be a uniformizer of \(\mathcal O_F\). By the local square theorem,
for every \(\alpha\in\mathcal O_F\), there exists \(\beta\in\mathcal O_F\) such
that
\[
1+4\pi\alpha=(1+2\pi\beta)^2.
\]
Hence
\[
1+4\pi\mathcal O_F\subset F^{\times 2}.
\]

Since \(4\pi\mathcal O_F\) is an open ideal of \(\mathcal O_F\), the set
\[
1+4\pi\mathcal O_F
\]
is an open neighborhood of \(1\) in \(F^\times\). Thus \(F^{\times 2}\) contains
an open neighborhood of the identity. Since \(F^{\times 2}\) is a subgroup of
the topological group \(F^\times\), it follows that \(F^{\times 2}\) is open.

Indeed, for any \(a\in F^{\times 2}\), the translate
\[
a(1+4\pi\mathcal O_F)
\]
is an open neighborhood of \(a\) contained in \(F^{\times 2}\). Therefore
\(F^{\times 2}\) is a union of open sets, hence open.

Finally, every open subgroup of a topological group is closed, because its
complement is a union of the other cosets:
\[
F^\times\setminus F^{\times 2}
=
\bigcup_{g\notin F^{\times 2}} gF^{\times 2},
\]
and each coset \(gF^{\times 2}\) is open. Hence the complement is open, so
\(F^{\times 2}\) is closed.

Therefore \(F^{\times 2}\) is both open and closed in \(F^\times\).
\end{solution}


\section*{Homework 12}

In this section, let $F$ be a henselian discrete valuation field of
characteristic $\neq 2$, let $\pi\in \mathcal O_F$ be a uniformizer and
$k=\mathcal O_F/\pi\mathcal O_F$. The normalized discrete valuation on $F$ is
denoted by $v$ or $v_F$.

\subsection*{Exercise A.12.1} Let $\mathfrak a$ be a fractional ideal of $F$. 
\begin{enumerate} 
    \item For any $\alpha,\beta\in F^\times$, let us write $\alpha\simeq \beta$ if $\alpha\beta^{-1}\in \mathcal O_F^{\times 2}$. 
        \begin{enumerate} 
            \item Prove that the above relation $\simeq$ is an equivalence relation on the set $F^\times$. 
            \item Show that for any $\alpha,\beta\in F^\times$, we have \[ \alpha\simeq \beta \Longleftrightarrow \alpha\beta^{-1}\in \mathcal O_F^\times \text{ and } \mathfrak d(\alpha\beta^{-1})=0. \] 
        \end{enumerate} 
        \item For any $\alpha,\beta\in F^\times$, let us write $\alpha\simeq \beta \pmod{\mathfrak a}$ if $\alpha\beta^{-1}\in \mathcal O_F^\times$ and \[ \alpha=\beta\varepsilon^2 \pmod{\mathfrak a} \] for some $\varepsilon\in \mathcal O_F^\times$. 
            \begin{enumerate} 
                \item Prove that the above congruence relation modulo $\mathfrak a$ is an equivalence relation on $F^\times$. 
                \item Show that if $\alpha,\beta\in F^\times$ satisfies $\alpha\simeq \beta\pmod{\mathfrak a}$, then \[ \mathfrak d(\alpha\beta^{-1})\subseteq \beta^{-1}\mathfrak a. \] 
            \end{enumerate} 
        \item Assume $k$ is perfect of characteristic $2$. Let $\alpha,\beta\in F^\times$. 
            \begin{enumerate} 
                \item Prove that the relation $\alpha\simeq \beta\pmod{\mathfrak a}$ holds if and only if $\alpha\beta^{-1}\in \mathcal O_F^\times$ and \[ \mathfrak d(\alpha\beta^{-1})\subseteq \beta^{-1}\mathfrak a. \] 
                \item Show that if $\alpha\beta^{-1}\in \mathcal O_F^\times$, then \[ \alpha\simeq \beta \pmod{\alpha\pi\mathcal O_F}. \] 
            \end{enumerate} 
\end{enumerate}
\begin{solution}

Write $\mathcal O=\mathcal O_F$. We use the following description of the
quadratic defect:
\[
    \mathfrak d(c)=\bigcap_{x\in F}(c-x^2)\mathcal O .
\]
Thus $\mathfrak d(c)\subseteq (c-x^2)\mathcal O$ for every $x\in F$.

\paragraph*{(1.a)}
The relation $\simeq$ is the equivalence relation induced by the subgroup
$\mathcal O^{\times 2}\subseteq F^\times$. Indeed,
$\alpha\alpha^{-1}=1\in\mathcal O^{\times 2}$, if
$\alpha\beta^{-1}\in\mathcal O^{\times 2}$ then
$\beta\alpha^{-1}=(\alpha\beta^{-1})^{-1}\in\mathcal O^{\times 2}$, and if
$\alpha\beta^{-1},\beta\gamma^{-1}\in\mathcal O^{\times 2}$, then
$\alpha\gamma^{-1}=(\alpha\beta^{-1})(\beta\gamma^{-1})\in
\mathcal O^{\times 2}$.

\paragraph*{(1.b)}
Put $c=\alpha\beta^{-1}$. If $\alpha\simeq\beta$, then
$c\in\mathcal O^{\times 2}$, hence $c\in\mathcal O^\times$ and
$\mathfrak d(c)=0$.

Conversely, suppose $c\in\mathcal O^\times$ and $\mathfrak d(c)=0$. Then $c$
can be approximated arbitrarily closely by squares. Choose $x\in F$ such that
$v(c-x^2)>2v(2)$. Since $c\in\mathcal O^\times$, this forces $x\in\mathcal
O^\times$. Applying Hensel's lemma to $T^2-c$ at $x$, we get $c\in F^{\times
2}$. Since $v(c)=0$, its square root is a unit. Hence
$c\in\mathcal O^{\times 2}$, so $\alpha\simeq\beta$.

\paragraph*{(2.a)}
Reflexivity is clear by taking $\varepsilon=1$. For symmetry, suppose
$\alpha=\beta\varepsilon^2 \pmod{\mathfrak a}$. Then
\[
    \beta=\alpha\varepsilon^{-2}\pmod{\mathfrak a},
\]
because $\varepsilon^{-2}\in\mathcal O^\times$ and $\mathfrak a$ is an
$\mathcal O$-submodule of $F$. Also $\beta\alpha^{-1}\in\mathcal O^\times$.
Hence $\beta\simeq\alpha\pmod{\mathfrak a}$.

For transitivity, suppose
\[
    \alpha=\beta\varepsilon^2\pmod{\mathfrak a},
    \qquad
    \beta=\gamma\eta^2\pmod{\mathfrak a}.
\]
Then
\[
    \alpha=\gamma(\eta\varepsilon)^2\pmod{\mathfrak a},
\]
and $\alpha\gamma^{-1}\in\mathcal O^\times$. Hence
$\alpha\simeq\gamma\pmod{\mathfrak a}$.

\paragraph*{(2.b)}
Put $c=\alpha\beta^{-1}$. Since
$\alpha\simeq\beta\pmod{\mathfrak a}$, there exists
$\varepsilon\in\mathcal O^\times$ such that
\[
    \alpha-\beta\varepsilon^2\in\mathfrak a.
\]
Equivalently,
\[
    c-\varepsilon^2\in\beta^{-1}\mathfrak a.
\]
By the definition of quadratic defect,
\[
    \mathfrak d(c)\subseteq (c-\varepsilon^2)\mathcal O
    \subseteq \beta^{-1}\mathfrak a.
\]
Thus
\[
    \mathfrak d(\alpha\beta^{-1})\subseteq\beta^{-1}\mathfrak a.
\]

\paragraph*{(3.a)}
Put $c=\alpha\beta^{-1}$. The implication
\[
    \alpha\simeq\beta\pmod{\mathfrak a}
    \Longrightarrow
    c\in\mathcal O^\times
    \text{ and }
    \mathfrak d(c)\subseteq\beta^{-1}\mathfrak a
\]
is exactly (2.b).

Conversely, suppose $c\in\mathcal O^\times$ and
$\mathfrak d(c)\subseteq\beta^{-1}\mathfrak a$. Since $k$ is perfect of
characteristic $2$, every element of $k^\times$ is a square. Hence $c$ is
congruent modulo $\pi\mathcal O$ to a unit square. Therefore, by discreteness
of the valuation, there exists $\varepsilon\in\mathcal O^\times$ such that
\[
    \mathfrak d(c)=(c-\varepsilon^2)\mathcal O .
\]
Thus
\[
    c-\varepsilon^2\in \mathfrak d(c)\subseteq \beta^{-1}\mathfrak a.
\]
Multiplying by $\beta$, we get
\[
    \alpha-\beta\varepsilon^2\in\mathfrak a.
\]
Hence
\[
    \alpha\simeq\beta\pmod{\mathfrak a}.
\]

\paragraph*{(3.b)}
Put $c=\alpha\beta^{-1}\in\mathcal O^\times$. Since $k$ is perfect of
characteristic $2$, choose $\varepsilon\in\mathcal O^\times$ such that
\[
    c\equiv \varepsilon^2\pmod{\pi\mathcal O}.
\]
Then
\[
    c-\varepsilon^2\in\pi\mathcal O=c\pi\mathcal O,
\]
because $c$ is a unit. Multiplying by $\beta$, we obtain
\[
    \alpha-\beta\varepsilon^2\in \beta c\pi\mathcal O
    =
    \alpha\pi\mathcal O.
\]
Therefore
\[
    \alpha\simeq\beta\pmod{\alpha\pi\mathcal O}.
\]
\end{solution}


\subsection*{Exercise A.12.2}

For a nonempty subset $S\subseteq F$, define
\[
    \operatorname{ord}(S):=
    \inf\{v(x)\mid x\in S\}\in \mathbb Z\cup\{\pm\infty\}.
\]

For any $a\in F^\times$ we define its \emph{relative quadratic defect} as
\[
    \operatorname{dft}(a):=
    \operatorname{ord}\bigl(a^{-1}\mathfrak d(a)\bigr)
    \in \mathbb Z\cup\{+\infty\}.
\]

Prove that $\operatorname{dft}(a)\in \mathbb N\cup\{+\infty\}$ for all
$a\in F^\times$, and that there is a well-defined map
\[
    \operatorname{dft}:F^\times/F^{\times 2}
    \longrightarrow
    \mathbb N\cup\{+\infty\},
    \qquad
    aF^{\times 2}\longmapsto \operatorname{dft}(a).
\]

Let $a\in F^\times$ and put $\varepsilon=a\pi^{-v(a)}$. Prove
\[
    \operatorname{dft}(a)=
    \begin{cases}
        0, & \text{if } v(a)=R \text{ is odd},\\
        \operatorname{dft}(\varepsilon)
        =\operatorname{ord}\bigl(\mathfrak d(\varepsilon)\bigr),
        & \text{if } v(a)=R \text{ is even}.
    \end{cases}
\]

\begin{solution}
We use the definition
\[
    \mathfrak d(a)=\bigcap_{x\in F}(a-x^2)\mathcal O_F .
\]
Since the term with $x=0$ occurs in the intersection, we have
\[
    \mathfrak d(a)\subseteq a\mathcal O_F.
\]
Hence
\[
    a^{-1}\mathfrak d(a)\subseteq \mathcal O_F,
\]
so
\[
    \operatorname{dft}(a)
    =
    \operatorname{ord}\bigl(a^{-1}\mathfrak d(a)\bigr)
    \in \mathbb N\cup\{+\infty\}.
\]

We next show that $\operatorname{dft}(a)$ only depends on the square class of
$a$. Let $b\in F^\times$. Then
\[
    \mathfrak d(ab^2)
    =
    \bigcap_{x\in F}(ab^2-x^2)\mathcal O_F.
\]
Writing $x=by$, we get
\[
    ab^2-x^2=b^2(a-y^2).
\]
Since $y$ also runs through all of $F$, it follows that
\[
    \mathfrak d(ab^2)=b^2\mathfrak d(a).
\]
Therefore
\[
    (ab^2)^{-1}\mathfrak d(ab^2)
    =
    a^{-1}\mathfrak d(a),
\]
and hence
\[
    \operatorname{dft}(ab^2)=\operatorname{dft}(a).
\]
Thus
\[
    \operatorname{dft}:F^\times/F^{\times 2}
    \longrightarrow
    \mathbb N\cup\{+\infty\}
\]
is well-defined.

Now let $a\in F^\times$ and put
\[
    R=v(a),\qquad \varepsilon=a\pi^{-R}.
\]
Then $\varepsilon\in\mathcal O_F^\times$.

If $R$ is odd, then for every $x\in F$, the two valuations $v(a)=R$ and
$v(x^2)=2v(x)$ are distinct. Hence
\[
    v(a-x^2)=\min\{R,2v(x)\}\leq R.
\]
Taking $x=0$ gives equality, so
\[
    \operatorname{ord}(\mathfrak d(a))=R.
\]
Therefore
\[
    \operatorname{dft}(a)
    =
    \operatorname{ord}\bigl(a^{-1}\mathfrak d(a)\bigr)
    =
    R-R
    =
    0.
\]

If $R$ is even, write $R=2s$. Then
\[
    a=\pi^{2s}\varepsilon.
\]
For $x\in F$, write $x=\pi^s y$. Then
\[
    a-x^2
    =
    \pi^{2s}(\varepsilon-y^2).
\]
Since $y$ runs through all of $F$, we get
\[
    \mathfrak d(a)=\pi^{2s}\mathfrak d(\varepsilon).
\]
Hence
\[
    a^{-1}\mathfrak d(a)
    =
    \varepsilon^{-1}\mathfrak d(\varepsilon).
\]
Because $\varepsilon$ is a unit,
\[
    \operatorname{dft}(a)
    =
    \operatorname{dft}(\varepsilon)
    =
    \operatorname{ord}\bigl(\mathfrak d(\varepsilon)\bigr).
\]
Thus
\[
    \operatorname{dft}(a)=
    \begin{cases}
        0, & \text{if } v(a) \text{ is odd},\\
        \operatorname{dft}(\varepsilon)
        =
        \operatorname{ord}\bigl(\mathfrak d(\varepsilon)\bigr),
        & \text{if } v(a) \text{ is even}.
    \end{cases}
\]
\end{solution}

\subsection*{Exercise A.12.3}

With notation as in Exercise A.12.2, now assume further that the residue field
$k$ is a perfect field of characteristic $2$ which has a unique quadratic
extension. Let $\Delta\in \mathcal O_F^\times$ be chosen such that
\[
    \mathfrak d(\Delta)=4\mathcal O_F
\]
(cf. Lemma 4.2.12). Let $e=v_F(2)$.

Prove
\[
    \operatorname{dft}(F^\times)
    =
    \{0,2e,+\infty\}
    \cup
    \{2j-1\mid j\in [1,e]\}.
\]

Show that for any $a\in F^\times$, we have the following equivalences:
\[
    \operatorname{dft}(a)=0
    \Longleftrightarrow
    v(a)\text{ is odd},
\]
\[
    \operatorname{dft}(a)\geq 1
    \Longleftrightarrow
    v(a)\text{ is even},
\]
\[
    \operatorname{dft}(a)=2e
    \Longleftrightarrow
    aF^{\times 2}=\Delta F^{\times 2},
\]
\[
    \operatorname{dft}(a)=+\infty
    \Longleftrightarrow
    a\in F^{\times 2}.
\]

For any $a,b\in F^\times$, prove that
\[
    \begin{cases}
        \operatorname{dft}(ab)=0
        =
        \min\{\operatorname{dft}(a),\operatorname{dft}(b)\},
        & \text{if } v(a)\not\equiv v(b)\pmod 2,\\
        \operatorname{dft}(ab)=0
        =
        \operatorname{dft}(a)=\operatorname{dft}(b),
        & \text{if } v(a)\equiv v(b)\equiv 1\pmod 2,
    \end{cases}
\]
and that we have the \emph{domination principle}
\[
    \operatorname{dft}(ab)
    \geq
    \min\{\operatorname{dft}(a),\operatorname{dft}(b)\},
\]
and if $\operatorname{dft}(a)\neq \operatorname{dft}(b)$, then
\[
    \operatorname{dft}(ab)
    =
    \min\{\operatorname{dft}(a),\operatorname{dft}(b)\}.
\]
\begin{solution}
    Let $\mathcal O=\mathcal O_F$ and let $e=v(2)$. 
    By the exercise A.12.2, we reduce to the situation: $\text{dft}(\mathcal{O}^\times)$.
    According to our textbook, for non-squre unit, we obtain the following statement
    \begin{align*}
        \{\text{dft}(\varepsilon)|\varepsilon\text{ is nonsquare}\}=\{1,3,\cdots,2e-1\}.
    \end{align*}
    Hence, the unit suqre give $\text{dft}(\varepsilon)=+\infty$ and $\Delta$ gives $\text{dft}(\Delta)=2e$. Hence, if $v(a)$ is odd we then have $\text{dft}(a)=0$.
    Therefore, we obtain that 
    \[
    \operatorname{dft}(F^\times)
    =
    \{0,2e,+\infty\}
    \cup
    \{2j-1\mid j\in [1,e]\}.
\]
    Next we prove the stated equivalences.
    The first holds since $k$ is a perfect field of characteristic $2$, the frob map is surjective means every element in $k$ is squre, therefore we have $u-c^2 \in \pi$ implies that if $v(a=u\pi^{2r})$ is even then ord is not zero.
    The first two terms is quite clear. To give third, since it has unique quadratic extension, then the squre class in set 
    \begin{align*}
        \{\varepsilon| v(\varepsilon)= v(4)\}
    \end{align*}
    is only one, which menas that $\varepsilon F^{\times 2} =\Delta F^{\times 2}$
    
    The first two follow immediately from Exercise A.12.2 and the unit classification above: \[ \operatorname{dft}(a)=0 \Longleftrightarrow v(a)\text{ is odd}, \] and \[ \operatorname{dft}(a)\geq 1 \Longleftrightarrow v(a)\text{ is even}. \] If $v(a)$ is even, then $aF^{\times 2}=uF^{\times 2}$ for some $u\in\mathcal O^\times$. Hence \[ \operatorname{dft}(a)=2e \Longleftrightarrow uF^{\times 2}=\Delta F^{\times 2} \Longleftrightarrow aF^{\times 2}=\Delta F^{\times 2}. \] Similarly, \[ \operatorname{dft}(a)=+\infty \Longleftrightarrow a\in F^{\times 2}. \] 
    It remains to prove the domination principle. Since $\operatorname{dft}$ only depends on square classes, we may multiply $a$ and $b$ by squares whenever convenient. 
    If $v(a)\not\equiv v(b)\pmod 2$, then $v(ab)$ is odd. Hence, by Exercise A.12.2, \[ \operatorname{dft}(ab)=0 = \min\{\operatorname{dft}(a),\operatorname{dft}(b)\}. \] 
    If $v(a)\equiv v(b)\equiv 1\pmod 2$, then \[ \operatorname{dft}(a)=\operatorname{dft}(b)=0. \]
    Hence, it menas the unit part of $a,b$ is square unit, then we have $\text{dft}(ab)=0$.

    Therefore, the domination principal only remians both $v(a)$ and $v(b)$ are even.



    After multiplying by squares, we may assume $a,b\in\mathcal O^\times$. Put \[ r=\operatorname{dft}(a), \qquad s=\operatorname{dft}(b), \qquad m=\min\{r,s\}. \] If one of $r,s$ is $+\infty$, the assertion is immediate from the fact that squares do not change square classes. Thus assume $r,s<+\infty$. By Exercise A.12.1(3.a), there exist $\varepsilon,\eta\in\mathcal O^\times$ such that \[ a\equiv \varepsilon^2 \pmod{a\pi^r\mathcal O}, \qquad b\equiv \eta^2 \pmod{b\pi^s\mathcal O}. \] Since $a,b,\varepsilon,\eta$ are units, multiplying the two congruences gives \[ ab\equiv (\varepsilon\eta)^2 \pmod{ab\pi^m\mathcal O}. \] Again by Exercise A.12.1(3.a), \[ \operatorname{dft}(ab)\geq m. \] Thus \[ \operatorname{dft}(ab) \geq \min\{\operatorname{dft}(a),\operatorname{dft}(b)\}. \] Finally assume $\operatorname{dft}(a)\neq\operatorname{dft}(b)$. Without loss of generality, let $r<s$. The domination principle gives \[ \operatorname{dft}(ab)\geq r. \] If $\operatorname{dft}(ab)>r$, then applying the domination principle to $a=(ab)b^{-1}$ gives \[ r=\operatorname{dft}(a) = \operatorname{dft}((ab)b^{-1}) \geq \min\{\operatorname{dft}(ab),\operatorname{dft}(b^{-1})\}. \] Since $\operatorname{dft}(b^{-1})=\operatorname{dft}(b)=s$ and both $\operatorname{dft}(ab)$ and $s$ are strictly larger than $r$, this is a contradiction. Hence \[ \operatorname{dft}(ab)=r = \min\{\operatorname{dft}(a),\operatorname{dft}(b)\}. \]
\end{solution}
\section*{Homework 13}
\subsection*{Exercise A.13.1}
Let $F$ be a non-archimedean local field of characteristic $\neq 2$. If $\operatorname{char}(k) \neq 2$, prove that for any $\varepsilon, \delta \in \mathcal{O}_F^\times$ we have
\[ (\varepsilon, \delta)_F = 1 \quad \text{and} \quad (\pi, \delta)_F = \begin{cases} 1 & \text{if } \delta \in \mathcal{O}_F^{\times 2} \,, \\ -1 & \text{if } \delta \notin \mathcal{O}_F^{\times 2} \,. \end{cases} \]
\begin{solution}
We prove the two assertions by the following steps.

\emph{Step 1.} Use the quadratic-form criterion for the Hilbert symbol:
\((a,b)_F=1\) if and only if the ternary quadratic form
\(z^2-ax^2-by^2\) has a non-trivial zero over \(F\).

\emph{Step 2.} For \((\epsilon,\delta)_F\), reduce the form modulo \(\pi\).
Since the residue field \(k\) is finite of characteristic different from \(2\),
every quadratic form of dimension at least \(3\) over \(k\) is isotropic.
Then lift a non-trivial zero from \(k\) to \(\mathcal O_F\) by Hensel's lemma.

\emph{Step 3.} For \((\pi,\delta)_F\), the square case is immediate. In the
non-square case, assume a non-trivial zero exists, choose it primitive, and
reduce modulo \(\pi\). This forces all coordinates to be divisible by \(\pi\),
contradicting primitivity.

Now we give the proof.

First let \(\epsilon,\delta\in\mathcal O_F^\times\). Consider
\[
Q(x,y,z)=z^2-\epsilon x^2-\delta y^2.
\]
By the Hilbert-symbol criterion, it is enough to show that \(Q\) has a
non-trivial zero over \(F\). Reducing modulo \(\pi\), we obtain
\[
\overline Q(\bar x,\bar y,\bar z)
=\bar z^2-\bar\epsilon \bar x^2-\bar\delta \bar y^2
\]
over \(k\). Since \(k\) is finite and \(\operatorname{char}(k)\neq 2\), every
ternary quadratic form over \(k\) is isotropic. Hence there exists
\((\bar x_0,\bar y_0,\bar z_0)\neq (0,0,0)\) such that
\(\overline Q(\bar x_0,\bar y_0,\bar z_0)=0\).

At this point the gradient is
\[
\nabla\overline Q(\bar x_0,\bar y_0,\bar z_0)
=(-2\bar\epsilon\bar x_0,-2\bar\delta\bar y_0,2\bar z_0).
\]
Since \(\bar\epsilon,\bar\delta\neq 0\) and \(\operatorname{char}(k)\neq 2\),
these three partial derivatives cannot all vanish. Therefore at least one
partial derivative is non-zero. Fix the other two variables and apply the usual
one-variable Hensel lemma to the remaining variable. Thus the zero
\((\bar x_0,\bar y_0,\bar z_0)\) lifts to a zero
\((x_0,y_0,z_0)\in\mathcal O_F^3\) of \(Q\), and this zero is non-trivial.
Therefore
\[
(\epsilon,\delta)_F=1.
\]

It remains to compute \((\pi,\delta)_F\). If
\(\delta\in\mathcal O_F^{\times 2}\), say \(\delta=u^2\), then
\[
(\pi,\delta)_F=(\pi,u^2)_F=1.
\]

Now assume \(\delta\notin\mathcal O_F^{\times 2}\). Suppose, for contradiction,
that \((\pi,\delta)_F=1\). Then the form
\[
z^2-\pi x^2-\delta y^2
\]
has a non-trivial zero \((x,y,z)\in F^3\). Multiplying by a suitable power of
\(\pi\), we may assume that \(x,y,z\in\mathcal O_F\) and that at least one of
them is a unit. Such a solution is called primitive.

Reducing the equation modulo \(\pi\), we get
\[
\bar z^2-\bar\delta \bar y^2=0.
\]
If \(\bar y\neq 0\), then
\[
\bar\delta=\left(\frac{\bar z}{\bar y}\right)^2\in k^{\times 2}.
\]
By Hensel's lemma, this would imply
\(\delta\in\mathcal O_F^{\times 2}\), contradicting our assumption. Hence
\(\bar y=0\), and then the reduced equation gives \(\bar z=0\).

Thus \(y,z\in\pi\mathcal O_F\). Write \(y=\pi y_1\) and \(z=\pi z_1\). Substituting
back gives
\[
\pi^2 z_1^2-\pi x^2-\delta\pi^2 y_1^2=0.
\]
Dividing by \(\pi\), we obtain
\[
\pi z_1^2-x^2-\delta\pi y_1^2=0.
\]
Reducing modulo \(\pi\) gives \(-\bar x^2=0\), hence \(\bar x=0\). Therefore
\(x,y,z\) are all divisible by \(\pi\), contradicting the primitivity of the
solution.

So no non-trivial zero exists when
\(\delta\notin\mathcal O_F^{\times 2}\). Hence \((\pi,\delta)_F\neq 1\). Since
the Hilbert symbol takes values in \(\{\pm 1\}\), we conclude that
\[
(\pi,\delta)_F=-1.
\]

Therefore
\[
(\epsilon,\delta)_F=1
\]
for all \(\epsilon,\delta\in\mathcal O_F^\times\), and
\[
(\pi,\delta)_F=
\begin{cases}
1,& \delta\in\mathcal O_F^{\times 2},\\
-1,& \delta\notin\mathcal O_F^{\times 2}.
\end{cases}
\]
\end{solution}


\subsection*{Exercise A.13.2}
Let $K$ be a field of characteristic $\neq 2$, and let $a, b \in K^\times$. \\
Prove that the ternary quadratic space $\langle 1, -a, -b \rangle_K$ is isotropic if and only if the quaternary quadratic space $\langle 1, -a, -b, ab \rangle_K$ is isotropic.
\begin{solution}
We prove the equivalence by the following steps.

\emph{Step 1.} The implication from the ternary form to the quaternary form is
immediate by adding the fourth coordinate \(0\).

\emph{Step 2.} For the converse, assume that
\(\langle 1,-a,-b,ab\rangle_K\) is isotropic. If \(a\) is already a square in
\(K\), then \(\langle 1,-a,-b\rangle_K\) is clearly isotropic.

\emph{Step 3.} If \(a\notin K^{\times 2}\), pass to the quadratic extension
\(E=K(\sqrt a)\). Rewrite the isotropy equation as an equality of norms.

\emph{Step 4.} Divide the two norms to show that \(b\) is a norm from
\(E/K\). This gives an explicit non-trivial zero of
\(\langle 1,-a,-b\rangle_K\).

Now we give the proof.

First suppose that \(\langle 1,-a,-b\rangle_K\) is isotropic. Then there exist
\(x,y,z\in K\), not all zero, such that
\[
x^2-ay^2-bz^2=0.
\]
Hence \((x,y,z,0)\) is a non-zero isotropic vector for
\(\langle 1,-a,-b,ab\rangle_K\). Thus the quaternary form is isotropic.

Conversely, suppose that \(\langle 1,-a,-b,ab\rangle_K\) is isotropic. Then
there exist \(x,y,z,w\in K\), not all zero, such that
\[
x^2-ay^2-bz^2+ab w^2=0.
\]
If \(a\in K^{\times 2}\), say \(a=c^2\), then
\(c^2-a\cdot 1^2=0\), so \(\langle 1,-a,-b\rangle_K\) is already isotropic.
Thus we may assume that \(a\notin K^{\times 2}\).

Let \(E=K(\sqrt a)\). Then \(E/K\) is a quadratic field extension, and
\[
N_{E/K}(r+s\sqrt a)=r^2-a s^2.
\]
The equation above can be rewritten as
\[
x^2-ay^2=b(z^2-aw^2),
\]
or equivalently
\[
N_{E/K}(x+y\sqrt a)=b\,N_{E/K}(z+w\sqrt a).
\]
We claim that \(z+w\sqrt a\neq 0\). Indeed, if \(z+w\sqrt a=0\), then
\(z=w=0\). The equation then gives \(x^2-ay^2=0\), so
\(x+y\sqrt a=0\), hence \(x=y=0\), a contradiction.

Therefore we may divide by \(z+w\sqrt a\). We get
\[
b=N_{E/K}\left(\frac{x+y\sqrt a}{z+w\sqrt a}\right).
\]
Write
\[
\frac{x+y\sqrt a}{z+w\sqrt a}=u+v\sqrt a
\]
with \(u,v\in K\). Then
\[
b=u^2-av^2.
\]
Hence
\[
u^2-av^2-b=0.
\]
Thus \((u,v,1)\) is a non-zero isotropic vector for
\(\langle 1,-a,-b\rangle_K\). Therefore the ternary form is isotropic.

This proves that
\[
\langle 1,-a,-b\rangle_K
\]
is isotropic if and only if
\[
\langle 1,-a,-b,ab\rangle_K
\]
is isotropic.
\end{solution}


\subsection*{Exercise A.13.3}
Let $K$ be a field of characteristic $\neq 2$, and let $(U, q), (V, Q)$ be non-singular binary quadratic spaces over $K$. \\
Prove that the following statements are equivalent:
\begin{itemize}
    \item[(i)] $(U, q) \cong (V, Q)$ as quadratic spaces over $K$.
    \item[(ii)] $q(U)^\times \cap Q(V)^\times \neq \emptyset$ and $\det(U) = \det(V) \in K^\times / K^{\times 2}$. \\
    Here 
    \[ Q(V)^\times := K^\times \cap \{ Q(v) \mid v \in V \}. \]
    and similarly for $q(U)^\times$.
\end{itemize}
\begin{solution}
We prove the equivalence by the following steps.

\emph{Step 1.} The implication \((i)\Rightarrow (ii)\) is immediate: an
isometry preserves represented values and determinants.

\emph{Step 2.} For the converse, choose a common non-zero represented value
\(c\in q(U)^*\cap Q(V)^*\).

\emph{Step 3.} Use the vector representing \(c\) to split each binary quadratic
space into an orthogonal sum of two one-dimensional spaces.

\emph{Step 4.} Since the determinants are equal, the remaining one-dimensional
summands are isometric. Hence the two binary spaces are isometric.

Now we give the proof.

First suppose that \((U,q)\cong (V,Q)\). Then clearly \(q(U)^*=Q(V)^*\), so
\(q(U)^*\cap Q(V)^*\neq \emptyset\). Also an isometry preserves the determinant,
hence \(\det(U)=\det(V)\) in \(K^\times/K^{\times 2}\). Therefore \((ii)\)
holds.

Conversely, assume that \((ii)\) holds. Choose
\[
c\in q(U)^*\cap Q(V)^*.
\]
Then there exist \(u\in U\) and \(v\in V\) such that
\[
q(u)=c,\qquad Q(v)=c.
\]
Since \(c\neq 0\) and \(\operatorname{char}(K)\neq 2\), the one-dimensional
subspaces \(Ku\subset U\) and \(Kv\subset V\) are non-singular. Hence we have
orthogonal decompositions
\[
U=Ku\perp u^\perp,\qquad V=Kv\perp v^\perp .
\]
Since \(U\) and \(V\) are binary, both \(u^\perp\) and \(v^\perp\) are
one-dimensional. Therefore, for some \(d,e\in K^\times\), we may write
\[
(U,q)\cong \langle c,d\rangle,\qquad (V,Q)\cong \langle c,e\rangle .
\]

By assumption,
\[
\det(U)=\det(V)\quad \text{in } K^\times/K^{\times 2}.
\]
Using the above diagonal forms, this says
\[
cd=ce \quad \text{in } K^\times/K^{\times 2}.
\]
Since \(c\in K^\times\), we get
\[
d=e \quad \text{in } K^\times/K^{\times 2}.
\]
Thus \(d/e\in K^{\times 2}\), so the one-dimensional forms
\(\langle d\rangle\) and \(\langle e\rangle\) are isometric. Hence
\[
\langle c,d\rangle\cong \langle c,e\rangle.
\]
Therefore
\[
(U,q)\cong (V,Q).
\]

This proves the equivalence of \((i)\) and \((ii)\).
\end{solution}

\noindent In the rest of this section, let $F$ be a non-archimedean local field of characteristic $\neq 2$, let $\pi \in \mathcal{O}_F$ be a uniformier and $k = \mathcal{O}_F / \pi \mathcal{O}_F$. The normalized discrete valuation on $F$ is denoted by $v$ or $v_F$.

\subsection*{Exercise A.13.4}
Let $\delta \in \mathcal{O}_F^\times$ and assume that $\delta \notin F^{\times 2} \cup \Delta F^{\times 2}$. \\
Prove that there exist $\varepsilon_1, \varepsilon_2 \in \mathcal{O}_F^\times$ such that
\[ (\delta, \varepsilon_1)_F = (\delta, \varepsilon_2 \pi)_F = -1. \]
\begin{solution}
We prove the statement by the following steps.

\emph{Step 1.} Use the non-degeneracy of the Hilbert symbol on
\(F^\times/F^{\times 2}\).

\emph{Step 2.} Determine the orthogonal complement of the unit subgroup
\(\mathcal O_F^\times/F^{\times 2}\). The given property of \(\Delta\) implies
that this orthogonal complement is exactly \(\{1,\Delta\}\).

\emph{Step 3.} Since
\(\delta\notin F^{\times 2}\cup \Delta F^{\times 2}\), the class of \(\delta\)
is not orthogonal to all unit classes. Hence there exists
\(\epsilon_1\in\mathcal O_F^\times\) such that
\((\delta,\epsilon_1)_F=-1\).

\emph{Step 4.} Use bilinearity to find
\(\epsilon_2\in\mathcal O_F^\times\) such that
\((\delta,\epsilon_2\pi)_F=-1\).

Now we give the proof.

Let
\[
\overline F=F^\times/F^{\times 2},
\qquad
\overline U=\mathcal O_F^\times/\mathcal O_F^{\times 2}.
\]
The Hilbert symbol induces a non-degenerate bilinear pairing
\[
\overline F\times \overline F\longrightarrow \{\pm 1\}.
\]
Since \(F^\times=\pi^{\mathbb Z}\mathcal O_F^\times\), the subgroup
\(\overline U\subset \overline F\) has codimension \(1\). Hence
\(\overline U^\perp\) has order \(2\).

By the defining property of \(\Delta\), we have
\[
(\epsilon,\Delta)_F=1
\qquad
\text{for all }\epsilon\in\mathcal O_F^\times.
\]
Thus the class of \(\Delta\) lies in \(\overline U^\perp\). Also the class of
\(1\) lies in \(\overline U^\perp\). Since
\[
(\pi,\Delta)_F=-1,
\]
the class of \(\Delta\) is non-trivial. Therefore
\[
\overline U^\perp=\{1,\Delta\}.
\]

Now assume
\[
\delta\notin F^{\times 2}\cup \Delta F^{\times 2}.
\]
Equivalently, the square class of \(\delta\) is neither \(1\) nor \(\Delta\).
Hence the class of \(\delta\) does not lie in \(\overline U^\perp\). Therefore
\(\delta\) is not orthogonal to all unit classes. Thus there exists
\(\epsilon_1\in\mathcal O_F^\times\) such that
\[
(\delta,\epsilon_1)_F=-1.
\]

It remains to find \(\epsilon_2\in\mathcal O_F^\times\) such that
\[
(\delta,\epsilon_2\pi)_F=-1.
\]
If \((\delta,\pi)_F=-1\), take \(\epsilon_2=1\). If
\((\delta,\pi)_F=1\), take \(\epsilon_2=\epsilon_1\). In the second case,
bilinearity gives
\[
(\delta,\epsilon_1\pi)_F
=
(\delta,\epsilon_1)_F(\delta,\pi)_F
=
(-1)(1)
=
-1.
\]
Thus in both cases there exists
\(\epsilon_2\in\mathcal O_F^\times\) such that
\[
(\delta,\epsilon_2\pi)_F=-1.
\]

Therefore we have found \(\epsilon_1,\epsilon_2\in\mathcal O_F^\times\) with
\[
(\delta,\epsilon_1)_F=(\delta,\epsilon_2\pi)_F=-1.
\]
\end{solution}


\subsection*{Exercise A.13.5}
Prove the following formulas for the Hasse symbol of quadratic spaces over $F$:
\begin{align*}
    s(U \perp W) &= s(U)s(W) \cdot (\det(U), \det(W))_F \,, \\
    s(c\varphi) &= s(\varphi) \cdot \left( c, (-1)^{\frac{n(n+1)}{2}} \det(\varphi)^{n+1} \right)_F \,, \\
    s(m\mathbb{H}) &= \left( -1, -1 \right)_F^{\frac{m(m+1)}{2}} .
\end{align*}
Here $U = (U, \varphi)$ and $W$ denote nonsingular quadratic spaces, $n = \dim(\varphi)$, $c \in F^\times$, $m \in \mathbb{Z}_+$, and $\mathbb{H}$ denotes the quadratic hyperbolic plane.
\begin{solution}
We use the convention that, if
\(\varphi=\langle a_1,\dots,a_n\rangle\), then
\[
s(\varphi)=\prod_{1\le i\le j\le n}(a_i,a_j)_F,
\qquad
\det(\varphi)=a_1\cdots a_n\in F^\times/F^{\times 2}.
\]
We shall also use the standard identities
\[
(ab,c)_F=(a,c)_F(b,c)_F,\qquad
(a,bc)_F=(a,b)_F(a,c)_F,\qquad
(a,a)_F=(a,-1)_F.
\]

First let
\[
U=\langle a_1,\dots,a_r\rangle,\qquad
W=\langle b_1,\dots,b_s\rangle.
\]
Then
\[
U\perp W=\langle a_1,\dots,a_r,b_1,\dots,b_s\rangle.
\]
Hence
\[
s(U\perp W)
=
\prod_{1\le i\le j\le r}(a_i,a_j)_F
\prod_{1\le p\le q\le s}(b_p,b_q)_F
\prod_{i,p}(a_i,b_p)_F.
\]
The first two factors are \(s(U)\) and \(s(W)\). For the last factor, by
bilinearity of the Hilbert symbol,
\[
\prod_{i,p}(a_i,b_p)_F
=
\left(\prod_i a_i,\prod_p b_p\right)_F
=
(\det(U),\det(W))_F.
\]
Thus
\[
s(U\perp W)
=
s(U)s(W)(\det(U),\det(W))_F.
\]

Next let \(\varphi=\langle a_1,\dots,a_n\rangle\) and put
\(D=\det(\varphi)=a_1\cdots a_n\). Then
\[
c\varphi=\langle ca_1,\dots,ca_n\rangle,
\]
so
\[
s(c\varphi)=\prod_{1\le i\le j\le n}(ca_i,ca_j)_F.
\]
For each pair \(i\le j\), bilinearity gives
\[
(ca_i,ca_j)_F
=
(c,c)_F(c,a_j)_F(a_i,c)_F(a_i,a_j)_F.
\]
Multiplying over all \(i\le j\), the factors \((a_i,a_j)_F\) give
\(s(\varphi)\), while \((c,c)_F\) appears \(n(n+1)/2\) times. Since
\((c,c)_F=(c,-1)_F\), this contributes
\[
(c,-1)_F^{n(n+1)/2}
=
\left(c,(-1)^{n(n+1)/2}\right)_F.
\]
It remains to count the factors involving the \(a_i\)'s. For a fixed \(a_k\),
it appears \(k\) times among the factors \((c,a_j)_F\), and \(n-k+1\) times
among the factors \((a_i,c)_F\). Hence \(a_k\) appears \(n+1\) times in total.
Therefore these factors contribute
\[
(c,D^{n+1})_F.
\]
Combining the contributions, we obtain
\[
s(c\varphi)
=
s(\varphi)\cdot
\left(c,(-1)^{n(n+1)/2}D^{n+1}\right)_F.
\]
Thus
\[
s(c\varphi)
=
s(\varphi)\cdot
\left(c,(-1)^{n(n+1)/2}\det(\varphi)^{n+1}\right)_F.
\]

Finally, take the hyperbolic plane
\[
\mathbb H=\langle 1,-1\rangle.
\]
Then
\[
m\mathbb H
=
\langle 1,-1,\dots,1,-1\rangle
\]
with \(m\) copies of \(1\) and \(m\) copies of \(-1\). Since every Hilbert
symbol involving \(1\) is equal to \(1\), the only non-trivial contributions
to \(s(m\mathbb H)\) come from pairs of entries both equal to \(-1\). With the
convention \(i\le j\), there are
\[
m+\binom m2=\frac{m(m+1)}2
\]
such pairs. Hence
\[
s(m\mathbb H)
=
(-1,-1)_F^{m(m+1)/2}.
\]
This proves all three formulas.
\end{solution}


\subsection*{Exercise A.13.6}
Let $U, V$ be nonsingular quadratic spaces over $F$ with $d := \dim V - \dim U \ge 0$. Prove that $U$ is represented by $V$ (cf. (2.1.15)) if and only if one of the following conditions holds:
\begin{enumerate}
    \item $d = 0$ and $U \cong V$.
    \item $d = 1$ and $V \cong U \perp \langle \det(U)\det(V) \rangle$.
    \item $d = 2$ and $V \cong U \perp \mathbb{H}$.
    \item $d = 2$ and $\det(U) \neq -\det(V)$.
    \item $d \ge 3$.
\end{enumerate}
\begin{solution}
Recall that \(U\) is represented by \(V\) if and only if there is an isometric
embedding \(U\hookrightarrow V\). Since \(U\) is nonsingular, this is equivalent
to
\[
V\simeq U\perp W
\]
for some nonsingular quadratic space \(W\) of dimension
\(d=\dim V-\dim U\).

We first prove necessity. Suppose that \(U\) is represented by \(V\). Then
\(V\simeq U\perp W\), where \(W\) is nonsingular and \(\dim W=d\).

If \(d=0\), then \(W=0\), so \(U\simeq V\). This gives condition \(1\).

If \(d=1\), then \(W\simeq \langle c\rangle\) for some \(c\in F^\times\). Since
\(\det(V)=\det(U)\det(W)=\det(U)c\) in \(F^\times/F^{\times 2}\), we get
\(c=\det(U)\det(V)\). Hence
\[
V\simeq U\perp \langle \det(U)\det(V)\rangle,
\]
which is condition \(2\).

If \(d=2\), then \(W\) is binary and
\[
\det(W)=\det(U)\det(V).
\]
If \(\det(W)=-1\), then \(W\) is hyperbolic. Indeed, after diagonalizing
\(W=\langle a,b\rangle\), the condition \(ab=-1\) in
\(F^\times/F^{\times 2}\) implies \(W\simeq \langle a,-a^{-1}\rangle\), which
is isotropic, hence hyperbolic. Thus \(V\simeq U\perp \mathbb H\), giving
condition \(3\). If \(\det(W)\neq -1\), then
\[
\det(U)\det(V)\neq -1,
\]
or equivalently
\[
\det(U)\neq -\det(V),
\]
which is condition \(4\). Finally, if \(d\ge 3\), condition \(5\) holds.

We now prove sufficiency. Conditions \(1\), \(2\), and \(3\) directly give an
orthogonal decomposition \(V\simeq U\perp W\), so \(U\) is represented by \(V\).

Assume next that \(d=2\) and
\[
\det(U)\neq -\det(V).
\]
Put
\[
D=\det(U)\det(V).
\]
Then \(D\neq -1\). We shall construct a binary space \(W\) such that
\[
\det(W)=D
\]
and
\[
s(U\perp W)=s(V).
\]
By the formula for the Hasse symbol of an orthogonal sum,
\[
s(U\perp W)=s(U)s(W)(\det(U),D)_F.
\]
Thus it is enough to choose \(W\) with
\[
s(W)=s(V)s(U)(\det(U),D)_F.
\]
Let
\[
\eta=s(V)s(U)(\det(U),D)_F.
\]
For \(a\in F^\times\), set
\[
W_a=\langle a,aD\rangle.
\]
Then
\[
\det(W_a)=a^2D=D
\]
in \(F^\times/F^{\times 2}\). Using the convention
\(s(\langle x,y\rangle)=(x,x)_F(x,y)_F(y,y)_F\), we compute
\[
s(W_a)
=(a,a)_F(a,aD)_F(aD,aD)_F
=(a,-D)_F(D,-1)_F.
\]
Since \(D\neq -1\), the class of \(-D\) is nontrivial in
\(F^\times/F^{\times 2}\). By non-degeneracy of the Hilbert symbol, the
character
\[
a\longmapsto (a,-D)_F
\]
takes both values \(1\) and \(-1\). Hence we can choose \(a\in F^\times\) such
that
\[
s(W_a)=\eta.
\]
For this choice of \(a\), we have
\[
\dim(U\perp W_a)=\dim V,\qquad
\det(U\perp W_a)=\det(U)D=\det(V),
\]
and
\[
s(U\perp W_a)
=s(U)s(W_a)(\det(U),D)_F
=s(V).
\]
Thus \(U\perp W_a\) and \(V\) have the same dimension, the same determinant
and the same Hasse symbol. Therefore, by Theorem 4.2.27,
\[
V\simeq U\perp W_a.
\]
So \(U\) is represented by \(V\).

It remains to prove the case \(d\ge 3\). Again put
\[
D=\det(U)\det(V),
\qquad
\eta=s(V)s(U)(\det(U),D)_F.
\]
We shall construct a nonsingular space \(W\) of dimension \(d\) with
\[
\det(W)=D,\qquad s(W)=\eta.
\]
Let \(\pi\) be a uniformizer. Since \(-\pi\) is not a square class, the
character
\[
a\longmapsto (a,-\pi)_F
\]
is nontrivial by non-degeneracy of the Hilbert symbol.

For \(a\in F^\times\), define a three-dimensional space
\[
W_{3,a}:=\langle a,a\pi\rangle\perp \langle D\pi\rangle.
\]
Then
\[
\det(W_{3,a})=a^2\pi\cdot D\pi=D
\]
in \(F^\times/F^{\times 2}\). Its Hasse symbol is
\[
s(W_{3,a})
=s(\langle a,a\pi\rangle)s(\langle D\pi\rangle)(\pi,D\pi)_F.
\]
As above,
\[
s(\langle a,a\pi\rangle)=(a,-\pi)_F(\pi,-1)_F,
\]
and
\[
s(\langle D\pi\rangle)=(D\pi,-1)_F.
\]
Hence
\[
s(W_{3,a})
=(a,-\pi)_F(\pi,-1)_F(D\pi,-1)_F(\pi,D\pi)_F.
\]
All factors except \((a,-\pi)_F\) are fixed. Since
\(a\mapsto (a,-\pi)_F\) takes both values \(1\) and \(-1\), we may choose
\(a\) so that
\[
s(W_{3,a})=\eta.
\]
If \(d=3\), put \(W=W_{3,a}\). If \(d>3\), put
\[
W=W_{3,a}\perp \langle 1,\dots,1\rangle
\]
with \(d-3\) copies of \(\langle 1\rangle\). Adding these copies does not
change determinant or Hasse symbol. Thus in all cases \(d\ge 3\), we have
constructed \(W\) with
\[
\dim W=d,\qquad \det(W)=D,\qquad s(W)=\eta.
\]

Therefore
\[
\dim(U\perp W)=\dim V,
\]
and
\[
\det(U\perp W)=\det(U)D=\det(V).
\]
Moreover, by the orthogonal-sum formula,
\[
s(U\perp W)
=s(U)s(W)(\det(U),D)_F
=s(U)\eta(\det(U),D)_F
=s(V).
\]
Hence \(U\perp W\) and \(V\) have the same dimension, determinant and Hasse
symbol. By Theorem 4.2.27,
\[
V\simeq U\perp W.
\]
Thus \(U\) is represented by \(V\).

This proves the equivalence.
\end{solution}
\subsection*{Exercise A.13.7}
Let $V$ be a nonsingular quadratic space of dimension $n > 0$ over $F$. Prove that the following assertions are equivalent:
\begin{itemize}
    \item[(i)] There exists a nonsingular quadratic space $V'$ over $F$ such that
    \[ \dim V' = n \,,\; \operatorname{disc}(V') = \operatorname{disc}(V) \quad \text{and} \quad s(V') = -s(V). \]
    \item[(ii)] One of the following conditions holds:
    \begin{enumerate}
        \item $n = 2$ and $\operatorname{disc}(V) \neq 1 \in F^\times / F^{\times 2}$.
        \item $n \ge 3$.
    \end{enumerate}
\end{itemize}
\begin{solution}
We prove the equivalence by the following steps.

\emph{Step 1.} Since the dimension is fixed, having the same discriminant is
equivalent to having the same determinant. We use the classification theorem
of nonsingular quadratic spaces over \(F\): two nonsingular quadratic spaces
are isomorphic if and only if they have the same dimension, discriminant, and
Hasse symbol.

\emph{Step 2.} We show that the assertion is impossible when \(n=1\), and also
impossible when \(n=2\) and \(\disc(V)=1\).

\emph{Step 3.} If \(n=2\) and \(\disc(V)\neq 1\), we construct a binary form
with the same discriminant as \(V\) and the opposite Hasse symbol.

\emph{Step 4.} If \(n\ge 3\), we construct an \(n\)-dimensional form with the
same discriminant as \(V\) and the opposite Hasse symbol.

Now we give the proof.

First assume that there exists a nonsingular quadratic space \(V'\) over \(F\)
such that
\[
\dim V'=n,\qquad \disc(V')=\disc(V),\qquad s(V')=-s(V).
\]
If \(n=1\), then \(V=\langle a\rangle\) and \(V'=\langle b\rangle\). Having the
same discriminant means that \(a=b\) in \(F^\times/F^{\times 2}\). Hence
\(V\simeq V'\), and therefore \(s(V')=s(V)\), a contradiction. Thus \(n\neq 1\).

Now suppose \(n=2\) and \(\disc(V)=1\). Since
\[
\disc(V)=(-1)\det(V),
\]
we have \(\det(V)=-1\). Hence \(V\) is hyperbolic. Indeed, after diagonalizing
\(V=\langle a,b\rangle\), the condition \(ab=-1\) in
\(F^\times/F^{\times 2}\) gives
\[
V\simeq \langle a,-a^{-1}\rangle,
\]
which is isotropic, hence hyperbolic. The same argument applies to \(V'\).
Thus \(V'\simeq V\), so again \(s(V')=s(V)\), a contradiction. Therefore, if
the assertion holds, then either
\[
n=2 \text{ and } \disc(V)\neq 1,
\]
or
\[
n\ge 3.
\]

Conversely, first assume that \(n=2\) and \(\disc(V)\neq 1\). Put
\[
D=\det(V).
\]
Then \(\disc(V)=-D\), so the assumption \(\disc(V)\neq 1\) means
\[
D\neq -1.
\]
For \(a\in F^\times\), define
\[
V'_a:=\langle a,aD\rangle.
\]
Then
\[
\det(V'_a)=a^2D=D,
\]
so \(V'_a\) has the same discriminant as \(V\). We compute its Hasse symbol.
Using
\[
s(\langle x,y\rangle)=(x,x)_F(x,y)_F(y,y)_F
\]
and \((t,t)_F=(t,-1)_F\), we get
\[
s(V'_a)
=(a,a)_F(a,aD)_F(aD,aD)_F
=(a,-D)_F(D,-1)_F.
\]
Since \(D\neq -1\), the class of \(-D\) is nontrivial in
\(F^\times/F^{\times 2}\). By non-degeneracy of the Hilbert symbol, the
character
\[
a\longmapsto (a,-D)_F
\]
takes both values \(1\) and \(-1\). Therefore we can choose \(a\in F^\times\)
such that
\[
s(V'_a)=-s(V).
\]
For this choice of \(a\), the space \(V'_a\) has dimension \(2\), the same
discriminant as \(V\), and the opposite Hasse symbol. Thus condition \((i)\)
holds.

It remains to treat the case \(n\ge 3\). Put again
\[
D=\det(V).
\]
We construct an \(n\)-dimensional nonsingular quadratic space with determinant
\(D\) and Hasse symbol \(-s(V)\). Let \(\pi\) be a uniformizer. For
\(a\in F^\times\), set
\[
W_a:=\langle a,a\pi,D\pi\rangle.
\]
Then
\[
\det(W_a)=a^2\pi\cdot D\pi=D
\]
in \(F^\times/F^{\times 2}\). By the orthogonal-sum formula for the Hasse
symbol,
\[
s(W_a)
=
s(\langle a,a\pi\rangle)
s(\langle D\pi\rangle)
(\pi,D\pi)_F.
\]
As in the binary computation,
\[
s(\langle a,a\pi\rangle)=(a,-\pi)_F(\pi,-1)_F,
\]
and
\[
s(\langle D\pi\rangle)=(D\pi,-1)_F.
\]
Hence
\[
s(W_a)
=
(a,-\pi)_F
(\pi,-1)_F
(D\pi,-1)_F
(\pi,D\pi)_F.
\]
All factors except \((a,-\pi)_F\) are independent of \(a\). Since
\(-\pi\notin F^{\times 2}\), the character
\[
a\longmapsto (a,-\pi)_F
\]
takes both values \(1\) and \(-1\). Therefore we may choose \(a\) such that
\[
s(W_a)=-s(V).
\]

If \(n=3\), set \(V'=W_a\). If \(n>3\), set
\[
V'=W_a\perp \langle 1,\dots,1\rangle,
\]
where there are \(n-3\) copies of \(\langle 1\rangle\). Adding copies of
\(\langle 1\rangle\) changes neither determinant nor Hasse symbol. Hence
\[
\dim V'=n,\qquad \det(V')=D,\qquad s(V')=-s(V).
\]
Since the dimension is fixed, \(\det(V')=D=\det(V)\) is equivalent to
\(\disc(V')=\disc(V)\). Therefore \(V'\) satisfies
\[
\dim V'=n,\qquad \disc(V')=\disc(V),\qquad s(V')=-s(V).
\]

This proves the equivalence.
\end{solution}

\subsection*{Exercise A.13.8}
Suppose given a positive integer $n_0$, an element $d_0 \in F^\times / F^{\times 2}$ and $s_0 \in \{ \pm 1 \}$. Prove that the following assertions are equivalent:
\begin{itemize}
    \item[(i)] There exists a nonsingular quadratic space $V$ over $F$ such that
    \[ \dim V = n_0 \,,\; \det(V) = d_0 \quad \text{and} \quad s(V) = s_0. \]
    \item[(ii)] One of the following conditions holds:
    \begin{enumerate}
        \item $n_0 = 1$ and $s_0 = (d_0, -1)_F$.
        \item $n_0 = 2$, $s_0 = (d_0, -1)_F$ and $d_0 = -1 \in F^\times / F^{\times 2}$.
        \item $n_0 = 2$ and $d_0 \neq -1 \in F^\times / F^{\times 2}$.
        \item $n_0 \ge 3$.
    \end{enumerate}
\end{itemize}
\begin{solution}
We use the convention
\[
s(\langle a_1,\dots,a_n\rangle)
=
\prod_{1\le i\le j\le n}(a_i,a_j)_F.
\]
Also, for fixed dimension, equality of discriminants is equivalent to equality
of determinants.

We prove the equivalence by the following steps.

\emph{Step 1.} In dimensions \(1\) and \(2\), identify the cases where the
Hasse symbol is forced.

\emph{Step 2.} In the remaining cases, construct a form \(V_0\) with the
required dimension and determinant.

\emph{Step 3.} If \(s(V_0)\neq s_0\), apply Exercise A.13.7 to replace \(V_0\)
by another form with the same dimension and determinant but opposite Hasse
symbol.

Now suppose first that such a quadratic space \(V\) exists.

If \(n_0=1\), then \(V\simeq \langle d_0\rangle\). Hence
\[
s_0=s(V)=(d_0,d_0)_F=(d_0,-1)_F.
\]
Thus condition \(1\) holds.

If \(n_0=2\) and \(d_0=-1\), then \(V\) is hyperbolic, since a binary form of
determinant \(-1\) is isotropic. Hence
\[
s_0=s(V)=s(\mathbb H)=(-1,-1)_F=(d_0,-1)_F.
\]
Thus condition \(2\) holds. If \(n_0=2\) and \(d_0\neq -1\), then condition
\(3\) holds. If \(n_0\ge 3\), then condition \(4\) holds. This proves the
necessity.

Conversely, assume that one of the listed conditions holds.

If \(n_0=1\), take
\[
V=\langle d_0\rangle .
\]
Then
\[
\dim V=1,\qquad \det(V)=d_0,\qquad
s(V)=(d_0,d_0)_F=(d_0,-1)_F=s_0.
\]

If \(n_0=2\), \(d_0=-1\), and \(s_0=(d_0,-1)_F\), take
\[
V=\mathbb H=\langle 1,-1\rangle .
\]
Then
\[
\dim V=2,\qquad \det(V)=-1=d_0,\qquad
s(V)=(-1,-1)_F=(d_0,-1)_F=s_0.
\]

If \(n_0=2\) and \(d_0\neq -1\), take first
\[
V_0=\langle d_0,1\rangle .
\]
Then
\[
\dim V_0=2,\qquad \det(V_0)=d_0.
\]
If \(s(V_0)=s_0\), take \(V=V_0\). If \(s(V_0)\neq s_0\), then
\(s_0=-s(V_0)\). Since \(d_0\neq -1\), we have
\[
\operatorname{disc}(V_0)=-d_0\neq 1.
\]
By Exercise A.13.7, there exists a nonsingular binary quadratic space \(V\)
such that
\[
\dim V=2,\qquad
\operatorname{disc}(V)=\operatorname{disc}(V_0),
\qquad
s(V)=-s(V_0)=s_0.
\]
Since the dimension is fixed, \(\operatorname{disc}(V)=\operatorname{disc}(V_0)\)
is equivalent to \(\det(V)=\det(V_0)=d_0\).

Finally, suppose \(n_0\ge 3\). Take
\[
V_0=\langle d_0,1,\dots,1\rangle ,
\]
with \(n_0-1\) copies of \(1\). Then
\[
\dim V_0=n_0,\qquad \det(V_0)=d_0.
\]
If \(s(V_0)=s_0\), take \(V=V_0\). If \(s(V_0)\neq s_0\), then
\(s_0=-s(V_0)\). By Exercise A.13.7, there exists a nonsingular quadratic
space \(V\) such that
\[
\dim V=n_0,\qquad
\operatorname{disc}(V)=\operatorname{disc}(V_0),
\qquad
s(V)=-s(V_0)=s_0.
\]
Again, since the dimension is fixed, this gives
\[
\det(V)=\det(V_0)=d_0.
\]

Therefore there exists a nonsingular quadratic space \(V\) over \(F\) with
\[
\dim V=n_0,\qquad \det(V)=d_0,\qquad s(V)=s_0
\]
if and only if one of the listed conditions holds.
\end{solution}

\section*{Homework 14}
\subsection*{Exercise A.14.1}
Prove Lemma 3.2.2: Let $(V, \| \cdot \|)$ be a normed vector space over $(K, | \cdot |)$. Then the maps
\begin{align*}
    \| \cdot \| : V \longrightarrow \mathbb{R} \ &; \quad x \longmapsto \|x\| \,, \\
    V \times V \longrightarrow V \ &; \quad (x, y) \longmapsto x + y \,, \\
    K \times V \longrightarrow V \ &; \quad (\alpha, x) \longmapsto \alpha x \,, \\
    V \longrightarrow V \ &; \quad x \longmapsto -x
\end{align*}
are all continuous (for the metric topologies on $\mathbb{R}, K$ and $V$).
\begin{solution}
We prove the four assertions separately. Put
\[
d_V(x,y)=\|x-y\|,\qquad d_K(\alpha,\beta)=|\alpha-\beta|.
\]
On product spaces we use the max metric, which induces the product topology.

First, the norm map \(V\to \mathbb R\), \(x\mapsto \|x\|\), is continuous.
Indeed, by the triangle inequality,
\[
\|x\|\le \|y\|+\|x-y\|,
\qquad
\|y\|\le \|x\|+\|x-y\|.
\]
Hence
\[
\bigl|\|x\|-\|y\|\bigr|\le \|x-y\|.
\]
Thus the norm map is \(1\)-Lipschitz, hence continuous.

Next, consider addition
\[
V\times V\longrightarrow V,\qquad (x,y)\longmapsto x+y.
\]
Let \((x_0,y_0)\in V\times V\). Then
\[
\|(x+y)-(x_0+y_0)\|
=
\|(x-x_0)+(y-y_0)\|
\le \|x-x_0\|+\|y-y_0\|.
\]
Therefore, if \(\|x-x_0\|<\varepsilon/2\) and \(\|y-y_0\|<\varepsilon/2\), then
\[
\|(x+y)-(x_0+y_0)\|<\varepsilon.
\]
So addition is continuous.

Now consider scalar multiplication
\[
K\times V\longrightarrow V,\qquad (\alpha,x)\longmapsto \alpha x.
\]
Fix \((\alpha_0,x_0)\in K\times V\). We have
\[
\alpha x-\alpha_0x_0
=
(\alpha-\alpha_0)x_0+\alpha(x-x_0),
\]
hence
\[
\|\alpha x-\alpha_0x_0\|
\le
|\alpha-\alpha_0|\|x_0\|+|\alpha|\|x-x_0\|.
\]
Choose \(\delta\le 1\). If \(|\alpha-\alpha_0|<\delta\), then
\[
|\alpha|\le |\alpha_0|+1.
\]
Thus, whenever
\[
|\alpha-\alpha_0|<\delta,\qquad \|x-x_0\|<\delta,
\]
we get
\[
\|\alpha x-\alpha_0x_0\|
\le
\delta\|x_0\|+(|\alpha_0|+1)\delta
=
\delta(\|x_0\|+|\alpha_0|+1).
\]
Taking
\[
\delta<\frac{\varepsilon}{\|x_0\|+|\alpha_0|+1}
\]
shows that scalar multiplication is continuous at \((\alpha_0,x_0)\). Since
\((\alpha_0,x_0)\) was arbitrary, scalar multiplication is continuous.

Finally, the map
\[
V\longrightarrow V,\qquad x\longmapsto -x
\]
is continuous because
\[
\|(-x)-(-y)\|=\|x-y\|.
\]
Thus it is an isometry.

Therefore the norm map, addition, scalar multiplication, and negation are all
continuous for the metric topologies on \(\mathbb R\), \(K\), and \(V\).
\end{solution}
\subsection*{Exercise A.14.2}
With notation as in (3.2.5), prove that the following maps are all continuous:
\begin{align*}
    \operatorname{End}(V) \times \operatorname{End}(V) \longrightarrow \operatorname{End}(V) \ &; \quad (\sigma, \tau) \longmapsto \sigma\tau \,, \\
    \operatorname{End}(V) \times V \longrightarrow V \ &; \quad (\sigma, x) \longmapsto \sigma(x) \,, \\
    \operatorname{End}(V) \longrightarrow K \ &; \quad \sigma \longmapsto \det(\sigma)
\end{align*}
Show also that if the absolute value on $K$ is non-archimedean, then $\| \sigma\tau \| \le \| \sigma \| \cdot \| \tau \|$ for all $\sigma, \tau \in \operatorname{End}(V)$.
\begin{solution}
Let \(\mathcal E=(e_1,\dots,e_n)\) be the fixed ordered basis of \(V\). Write
\[
\sigma(e_j)=\sum_{i=1}^n \alpha_{ij}e_i,\qquad
\tau(e_j)=\sum_{i=1}^n \beta_{ij}e_i.
\]
Thus
\[
\|\sigma\|_{\mathcal E}=\max_{i,j}|\alpha_{ij}|,
\qquad
\|\tau\|_{\mathcal E}=\max_{i,j}|\beta_{ij}|.
\]

We prove the continuity of the three maps separately.

First consider composition
\[
\operatorname{End}(V)\times \operatorname{End}(V)\longrightarrow
\operatorname{End}(V),\qquad
(\sigma,\tau)\longmapsto \sigma\tau.
\]
The matrix of \(\sigma\tau\) has entries
\[
(\sigma\tau)_{ij}=\sum_{k=1}^n \alpha_{ik}\beta_{kj}.
\]
Each entry is a finite sum of products of the matrix entries of \(\sigma\) and
\(\tau\). Since addition and multiplication in \(K\) are continuous, each
coordinate function \((\sigma,\tau)\mapsto (\sigma\tau)_{ij}\) is continuous.
Because the maximum norm topology on \(\operatorname{End}(V)\simeq K^{n^2}\)
is the product topology, it follows that composition is continuous.

Next consider the evaluation map
\[
\operatorname{End}(V)\times V\longrightarrow V,\qquad
(\sigma,x)\longmapsto \sigma(x).
\]
Write
\[
x=\sum_{j=1}^n x_je_j.
\]
Then
\[
\sigma(x)=\sum_{i=1}^n\left(\sum_{j=1}^n\alpha_{ij}x_j\right)e_i.
\]
Again each coordinate of \(\sigma(x)\) is a finite sum of products
\(\alpha_{ij}x_j\). Hence each coordinate function is continuous, and therefore
the evaluation map is continuous.

Finally, the determinant map
\[
\det:\operatorname{End}(V)\longrightarrow K
\]
is given by the polynomial formula
\[
\det(\sigma)=
\sum_{\rho\in S_n}\operatorname{sgn}(\rho)
\alpha_{1\rho(1)}\cdots \alpha_{n\rho(n)}.
\]
It is therefore a polynomial in the matrix entries of \(\sigma\), hence is
continuous.

It remains to prove the norm inequality in the non-archimedean case. Assume
that the absolute value on \(K\) is non-archimedean. Then, for each \(i,j\),
\[
|(\sigma\tau)_{ij}|
=
\left|\sum_{k=1}^n \alpha_{ik}\beta_{kj}\right|
\le
\max_{1\le k\le n}|\alpha_{ik}\beta_{kj}|
\le
\|\sigma\|_{\mathcal E}\|\tau\|_{\mathcal E}.
\]
Taking the maximum over all \(i,j\), we get
\[
\|\sigma\tau\|_{\mathcal E}
\le
\|\sigma\|_{\mathcal E}\|\tau\|_{\mathcal E}.
\]
This proves the desired inequality.
\end{solution}
\subsection*{Exercise A.14.3}
Let $F$ be a field of characteristic $\neq 2$ and let $\mathcal{O}_F$ be a Dedekind domain with fraction field $F$. \\
Let $L$ be a nonsingular quadratic lattice over $\mathcal{O}_F$ and let $\mathfrak{a}$ be a fractional ideal of $\mathcal{O}_F$. Define
\[ L^{\mathfrak{a}\#} := \{ x \in L \mid B(x, L) \subseteq \mathfrak{a} \}. \]
\begin{enumerate}
    \item Prove that $L^{\mathfrak{a}\#} = L \cap \mathfrak{a}L^\#$ and that $L^{\mathfrak{a}\#}$ spans the same $F$-vector space as $L$.
    \item Let $\mathfrak{b}$ be another fractional ideal of $\mathcal{O}_F$. Prove that if $L$ is $\mathfrak{b}$-modular, then
    \[ L^{\mathfrak{a}\#} = \begin{cases} L & \text{if } \mathfrak{b} \subseteq \mathfrak{a} \,, \\ \mathfrak{a}\mathfrak{b}^{-1}L & \text{if } \mathfrak{a} \subseteq \mathfrak{b} \,. \end{cases} \]
    \item Prove that the following conditions are equivalent:
    \begin{itemize}
        \item[(i)] $\mathfrak{a} \subseteq \mathfrak{b}$.
        \item[(ii)] $\mathfrak{a}\mathfrak{b}^{-1}L^{\mathfrak{b}\#} \subseteq L^{\mathfrak{a}\#} \subseteq L^{\mathfrak{b}\#}$.
    \end{itemize}
\end{enumerate}
\begin{solution}
Write \(\mathcal O=\mathcal O_F\), and let
\[
L^\#:=\{x\in F L\mid B(x,L)\subseteq \mathcal O\}
\]
be the usual dual lattice. We prove the assertions by the following steps.

\emph{Step 1.} Express the \(\mathfrak a\)-dual as
\[
L^{\mathfrak a\#}=L\cap \mathfrak a L^\#.
\]
Then \(L^{\mathfrak a\#}\) is an intersection of two full lattices, hence is
again a full lattice in \(F L\).

\emph{Step 2.} If \(L\) is \(\mathfrak b\)-modular, use
\(L^\#=\mathfrak b^{-1}L\) to compute
\[
L^{\mathfrak a\#}=L\cap \mathfrak a\mathfrak b^{-1}L.
\]
The two cases \(\mathfrak b\subseteq \mathfrak a\) and
\(\mathfrak a\subseteq \mathfrak b\) then follow from ideal inclusion.

\emph{Step 3.} For the final equivalence, prove one direction directly from
\(\mathfrak a\subseteq \mathfrak b\). For the converse, use the fact that if
\(M\) is a full \(\mathcal O\)-lattice, then
\[
\{\lambda\in F\mid \lambda M\subseteq M\}=\mathcal O.
\]

We now give the proof.

First, by definition,
\[
L^{\mathfrak a\#}
=
\{x\in L\mid B(x,L)\subseteq \mathfrak a\}.
\]
On the other hand,
\[
\mathfrak a L^\#
=
\{x\in F L\mid B(x,L)\subseteq \mathfrak a\}.
\]
Indeed, this may be checked after localizing at every prime ideal of
\(\mathcal O\), where \(\mathfrak a\) becomes principal. Hence
\[
L^{\mathfrak a\#}=L\cap \mathfrak a L^\#.
\]
Since \(L\) is nonsingular, \(L^\#\) is a full \(\mathcal O\)-lattice in
\(F L\). Therefore \(\mathfrak a L^\#\) is also a full lattice, and the
intersection of two full lattices is again a full lattice. Thus
\(L^{\mathfrak a\#}\) spans the same \(F\)-vector space as \(L\).

Now assume that \(L\) is \(\mathfrak b\)-modular, so
\[
L^\#=\mathfrak b^{-1}L.
\]
Using the formula just proved, we get
\[
L^{\mathfrak a\#}
=
L\cap \mathfrak a L^\#
=
L\cap \mathfrak a\mathfrak b^{-1}L.
\]
If \(\mathfrak b\subseteq \mathfrak a\), then
\[
\mathcal O\subseteq \mathfrak a\mathfrak b^{-1},
\]
so \(L\subseteq \mathfrak a\mathfrak b^{-1}L\). Hence
\[
L^{\mathfrak a\#}=L.
\]
If \(\mathfrak a\subseteq \mathfrak b\), then
\[
\mathfrak a\mathfrak b^{-1}\subseteq \mathcal O,
\]
so \(\mathfrak a\mathfrak b^{-1}L\subseteq L\). Hence
\[
L^{\mathfrak a\#}=\mathfrak a\mathfrak b^{-1}L.
\]
This proves the second assertion.

Finally we prove the equivalence
\[
\mathfrak a\subseteq \mathfrak b
\quad\Longleftrightarrow\quad
\mathfrak a\mathfrak b^{-1}L^{\mathfrak b\#}
\subseteq
L^{\mathfrak a\#}
\subseteq
L^{\mathfrak b\#}.
\]

Assume first that \(\mathfrak a\subseteq \mathfrak b\). Then
\[
\mathfrak a\mathfrak b^{-1}\subseteq \mathcal O.
\]
If \(x\in L^{\mathfrak a\#}\), then
\[
B(x,L)\subseteq \mathfrak a\subseteq \mathfrak b,
\]
so \(x\in L^{\mathfrak b\#}\). Thus
\[
L^{\mathfrak a\#}\subseteq L^{\mathfrak b\#}.
\]
Moreover, if \(t\in \mathfrak a\mathfrak b^{-1}\) and
\(x\in L^{\mathfrak b\#}\), then \(t\in\mathcal O\), so \(tx\in L\), and
\[
B(tx,L)=tB(x,L)\subseteq
\mathfrak a\mathfrak b^{-1}\mathfrak b
=
\mathfrak a.
\]
Therefore \(tx\in L^{\mathfrak a\#}\). Hence
\[
\mathfrak a\mathfrak b^{-1}L^{\mathfrak b\#}
\subseteq
L^{\mathfrak a\#}
\subseteq
L^{\mathfrak b\#}.
\]

Conversely, suppose that
\[
\mathfrak a\mathfrak b^{-1}L^{\mathfrak b\#}
\subseteq
L^{\mathfrak a\#}
\subseteq
L^{\mathfrak b\#}.
\]
Then for every \(t\in \mathfrak a\mathfrak b^{-1}\), we have
\[
tL^{\mathfrak b\#}\subseteq L^{\mathfrak b\#}.
\]
By the first part, \(L^{\mathfrak b\#}\) is a full \(\mathcal O\)-lattice.
For any full \(\mathcal O\)-lattice \(M\), one has
\[
\{\lambda\in F\mid \lambda M\subseteq M\}=\mathcal O.
\]
Indeed, after localizing at any prime \(\mathfrak p\), the module
\(M_{\mathfrak p}\) is free over the DVR \(\mathcal O_{\mathfrak p}\), and
\(\lambda M_{\mathfrak p}\subseteq M_{\mathfrak p}\) forces
\(\lambda\in\mathcal O_{\mathfrak p}\). Intersecting over all
\(\mathfrak p\) gives \(\lambda\in\mathcal O\).

Applying this to \(M=L^{\mathfrak b\#}\), we get
\[
\mathfrak a\mathfrak b^{-1}\subseteq \mathcal O.
\]
Since \(\mathfrak b\) is invertible in the Dedekind domain \(\mathcal O\), this
is equivalent to
\[
\mathfrak a\subseteq \mathfrak b.
\]
This proves the equivalence.
\end{solution}
\begin{solution}
We prove \((ii)\Rightarrow (i)\). Suppose that
\[
\mathfrak a\mathfrak b^{-1}L^{\mathfrak b\#}
\subseteq
L^{\mathfrak a\#}
\subseteq
L^{\mathfrak b\#}.
\]
Then in particular
\[
\mathfrak a\mathfrak b^{-1}L^{\mathfrak b\#}
\subseteq
L^{\mathfrak b\#}.
\]
By part (1), \(L^{\mathfrak b\#}\) is a full \(\mathcal O_F\)-lattice in
\(FL\). For any full lattice \(M\), one has
\[
\{\lambda\in F\mid \lambda M\subseteq M\}=\mathcal O_F.
\]
Indeed, if \(\lambda M\subseteq M\), then after localizing at every maximal
ideal \(\mathfrak p\), we get
\[
\lambda M_{\mathfrak p}\subseteq M_{\mathfrak p}.
\]
Since \(M_{\mathfrak p}\) is a free \(\mathcal O_{\mathfrak p}\)-lattice, this
forces \(\lambda\in\mathcal O_{\mathfrak p}\) for every \(\mathfrak p\), hence
\(\lambda\in\mathcal O_F\).

Applying this to \(M=L^{\mathfrak b\#}\), the inclusion
\[
\mathfrak a\mathfrak b^{-1}L^{\mathfrak b\#}
\subseteq
L^{\mathfrak b\#}
\]
implies
\[
\mathfrak a\mathfrak b^{-1}\subseteq \mathcal O_F.
\]
Therefore
\[
\mathfrak a\subseteq \mathfrak b.
\]
This proves \((ii)\Rightarrow(i)\).
\end{solution}
\subsection*{Exercise A.14.4}
Let $\mathcal{O}_F$ be a discrete valuation ring with fraction field $F$ and suppose $\operatorname{char}(F) \neq 2$. Let $\mathfrak{a}$ be a fractional ideal of $\mathcal{O}_F$ and let $L$ be a quadratic $\mathcal{O}_F$-lattice with a Jordan splitting
\[ L = L_1 \perp \cdots \perp L_t \perp \operatorname{rad}(L). \]
Let $L^{\mathfrak{a}\#}$ be defined as in Exercise A.14.3.
\begin{enumerate}
    \item Show that $\mathfrak{s}(L^{\mathfrak{a}\#}) \subseteq \mathfrak{a}$.
    \item Prove that $\mathfrak{s}(L^{\mathfrak{a}\#}) = \mathfrak{a}$ if and only if there exists $1 \le i \le t$ such that $\mathfrak{s}(L_i) = \mathfrak{a}$.
\end{enumerate}
\begin{solution}
Write \(\mathcal O=\mathcal O_F\). Let
\[
L=L_1\perp\cdots\perp L_t\perp \operatorname{rad}(L)
\]
be a Jordan splitting, and put
\[
\mathfrak b_i=\mathfrak s(L_i).
\]
Then each \(L_i\) is \(\mathfrak b_i\)-modular.

We prove the result by the following steps.

\emph{Step 1.} Use the orthogonal decomposition to write
\[
L^{\mathfrak a\#}
=
L_1^{\mathfrak a\#}\perp\cdots\perp L_t^{\mathfrak a\#}
\perp \operatorname{rad}(L).
\]

\emph{Step 2.} Compute each \(L_i^{\mathfrak a\#}\) using Exercise A.14.3.

\emph{Step 3.} Compute the scale ideal of each component and then take their
sum.

Since the decomposition is orthogonal, for
\(x=x_1+\cdots+x_t+r\in L\), with \(x_i\in L_i\) and
\(r\in \operatorname{rad}(L)\), we have
\[
B(x,L)\subseteq \mathfrak a
\]
if and only if
\[
B(x_i,L_i)\subseteq \mathfrak a
\qquad
\text{for all } i.
\]
Moreover \(B(r,L)=0\). Hence
\[
L^{\mathfrak a\#}
=
L_1^{\mathfrak a\#}\perp\cdots\perp L_t^{\mathfrak a\#}
\perp \operatorname{rad}(L).
\]

Now fix \(i\). Since \(L_i\) is \(\mathfrak b_i\)-modular, Exercise A.14.3
gives
\[
L_i^{\mathfrak a\#}
=
\begin{cases}
L_i, & \text{if } \mathfrak b_i\subseteq \mathfrak a,\\
\mathfrak a\mathfrak b_i^{-1}L_i, & \text{if } \mathfrak a\subseteq \mathfrak b_i.
\end{cases}
\]
Since \(\mathcal O\) is a discrete valuation ring, any two fractional ideals
are comparable, so these two cases cover all possibilities.

If \(\mathfrak b_i\subseteq \mathfrak a\), then
\[
\mathfrak s(L_i^{\mathfrak a\#})
=
\mathfrak s(L_i)
=
\mathfrak b_i
\subseteq \mathfrak a.
\]
If \(\mathfrak a\subseteq \mathfrak b_i\), then
\[
L_i^{\mathfrak a\#}
=
\mathfrak a\mathfrak b_i^{-1}L_i,
\]
and therefore
\[
\mathfrak s(L_i^{\mathfrak a\#})
=
(\mathfrak a\mathfrak b_i^{-1})^2\mathfrak s(L_i)
=
\mathfrak a^2\mathfrak b_i^{-1}.
\]
Since \(\mathfrak a\subseteq \mathfrak b_i\), we have
\(\mathfrak a\mathfrak b_i^{-1}\subseteq \mathcal O\). Hence
\[
\mathfrak a^2\mathfrak b_i^{-1}
\subseteq
\mathfrak a.
\]
Thus in all cases,
\[
\mathfrak s(L_i^{\mathfrak a\#})\subseteq \mathfrak a.
\]

Since the radical contributes no scale, we get
\[
\mathfrak s(L^{\mathfrak a\#})
=
\sum_{i=1}^t \mathfrak s(L_i^{\mathfrak a\#})
\subseteq \mathfrak a.
\]
This proves the first assertion.

It remains to prove the criterion for equality. From the computation above,
\[
\mathfrak s(L_i^{\mathfrak a\#})
=
\begin{cases}
\mathfrak b_i, & \text{if } \mathfrak b_i\subseteq \mathfrak a,\\
\mathfrak a^2\mathfrak b_i^{-1}, & \text{if } \mathfrak a\subseteq \mathfrak b_i.
\end{cases}
\]
In the first case, this ideal equals \(\mathfrak a\) if and only if
\(\mathfrak b_i=\mathfrak a\). In the second case, it equals
\(\mathfrak a\) if and only if
\[
\mathfrak a^2\mathfrak b_i^{-1}=\mathfrak a,
\]
equivalently \(\mathfrak a=\mathfrak b_i\). Therefore
\[
\mathfrak s(L_i^{\mathfrak a\#})=\mathfrak a
\quad\Longleftrightarrow\quad
\mathfrak s(L_i)=\mathfrak a.
\]

Since \(\mathcal O\) is a discrete valuation ring, fractional ideals are totally
ordered. Hence a finite sum of ideals contained in \(\mathfrak a\) equals
\(\mathfrak a\) if and only if at least one of them equals \(\mathfrak a\).
Thus
\[
\mathfrak s(L^{\mathfrak a\#})=\mathfrak a
\]
if and only if there exists \(1\le i\le t\) such that
\[
\mathfrak s(L_i^{\mathfrak a\#})=\mathfrak a.
\]
By the equivalence just proved, this is the same as saying that there exists
\(1\le i\le t\) such that
\[
\mathfrak s(L_i)=\mathfrak a.
\]
This proves the second assertion.
\end{solution}
\subsection*{Exercise A.14.5}
Let $F = \mathbb{Q}_2$ be the field of $2$-adic numbers, i.e., the $2$-adic completion of $\mathbb{Q}$. Let $\mathcal{O}_F$ be the valuation ring of $F$.
\begin{enumerate}
    \item Find a unit $\Delta \in \mathcal{O}_F^\times$ with quadratic defect $\mathfrak{d}(\Delta) = 4\mathcal{O}_F$.
    \item Find a complete system $S$ of representatives for $F^\times / F^{\times 2}$ and compute the Hilbert symbol $(a, b)_F$ for all $a, b \in S$.
\end{enumerate}
\begin{solution}
We write \(F=\mathbb Q_2\) and \(\mathcal O_F=\mathbb Z_2\).

We prove the result by the following steps.

\emph{Step 1.} Find a unit \(\Delta\) with quadratic defect \(4\mathbb Z_2\).

\emph{Step 2.} Determine a complete system of representatives for
\(\mathbb Q_2^\times/\mathbb Q_2^{\times 2}\).

\emph{Step 3.} Use the standard formula for the Hilbert symbol over
\(\mathbb Q_2\) to compute the table.

First, take
\[
\Delta=5.
\]
Every odd square in \(\mathbb Z_2\) is congruent to \(1\) modulo \(8\). Hence
\(5\) is congruent to a square modulo \(4\), since \(5\equiv 1\pmod 4\), but
not congruent to a square modulo \(8\), since \(5\not\equiv 1\pmod 8\).
Therefore
\[
\mathfrak d(5)=4\mathbb Z_2.
\]
Thus we may take
\[
\Delta=5.
\]

Next, we determine the square classes. Since
\[
\mathbb Q_2^\times=2^{\mathbb Z}\mathbb Z_2^\times,
\]
the valuation contributes one factor represented by \(1\) and \(2\). For units,
one has
\[
u\in \mathbb Z_2^{\times 2}
\quad\Longleftrightarrow\quad
u\equiv 1\pmod 8.
\]
Thus the unit square classes are represented by
\[
1,\ -1,\ 5,\ -5.
\]
Hence a complete set of representatives for
\(\mathbb Q_2^\times/\mathbb Q_2^{\times 2}\) is
\[
S=\{1,-1,2,-2,5,-5,10,-10\}.
\]

We now compute the Hilbert symbol. For
\[
a=2^\alpha u,\qquad b=2^\beta v,
\]
where \(u,v\in\mathbb Z_2^\times\), the Hilbert symbol over \(\mathbb Q_2\) is
given by
\[
(a,b)_2
=
(-1)^{
\frac{u-1}{2}\frac{v-1}{2}
+\alpha\frac{v^2-1}{8}
+\beta\frac{u^2-1}{8}
}.
\]
The exponent is understood modulo \(2\). Evaluating this formula on the above
representatives gives the following table:
\[
\begin{array}{c|rrrrrrrr}
(a,b)_2
& 1 & -1 & 2 & -2 & 5 & -5 & 10 & -10\\
\hline
1
& 1 & 1 & 1 & 1 & 1 & 1 & 1 & 1\\
-1
& 1 & -1 & 1 & -1 & 1 & -1 & 1 & -1\\
2
& 1 & 1 & 1 & 1 & -1 & -1 & -1 & -1\\
-2
& 1 & -1 & 1 & -1 & -1 & 1 & -1 & 1\\
5
& 1 & 1 & -1 & -1 & 1 & 1 & -1 & -1\\
-5
& 1 & -1 & -1 & 1 & 1 & -1 & -1 & 1\\
10
& 1 & 1 & -1 & -1 & -1 & -1 & 1 & 1\\
-10
& 1 & -1 & -1 & 1 & -1 & 1 & 1 & -1
\end{array}
\]
This gives the complete Hilbert symbol table on
\(\mathbb Q_2^\times/\mathbb Q_2^{\times 2}\).

Acutally, we just need to compute 
\begin{align*}
    (2,5) = (-1,-1) = -1\quad (2,-1)=(5,-1) = 1.
\end{align*}
\end{solution}
\subsection*{Exercise A.14.6}
Let $F = \mathbb{Q}_5$ and let $\mathcal{O}_F$ be the valuation ring of $F$. Consider the quadratic $\mathcal{O}_F$-lattices
\[ L = \langle 1, 25 \rangle \quad \text{and} \quad M = \begin{pmatrix} 3 & 6 \\ 6 & 87 \end{pmatrix}. \]
\begin{enumerate}
    \item Prove that there is an isomorphism of quadratic spaces $FL \cong FM$.
    \item Find a Jordan splitting for each of the two lattices $L, M$.
    \item Show that $L$ and $M$ are not isomorphic as quadratic $\mathcal{O}_F$-lattices.
\end{enumerate}
\begin{solution}
We write \(\mathcal O=\mathbb Z_5\) and \(F=\mathbb Q_5\). We prove the
assertions in the following steps.

\emph{Step 1.} Diagonalize \(M\) over \(\mathcal O\).

\emph{Step 2.} Compare the quadratic spaces \(FL\) and \(FM\) by dimension,
determinant and Hasse symbol.

\emph{Step 3.} Read off the Jordan splittings from the diagonal forms.

\emph{Step 4.} Show that \(L\) and \(M\) are not isomorphic as
\(\mathcal O\)-lattices by comparing unit values modulo \(5\).

First diagonalize \(M\). Let \(f_1=e_1\) and \(f_2=e_2-2e_1\). Since the
change-of-basis matrix has determinant \(1\), this is an
\(\mathcal O\)-basis. We compute
\[
B(f_1,f_1)=3,\qquad B(f_1,f_2)=6-2\cdot 3=0,
\]
and
\[
B(f_2,f_2)=87-4\cdot 6+4\cdot 3=75.
\]
Hence
\[
M\simeq \langle 3,75\rangle
\]
as an \(\mathcal O\)-lattice.

Now compare \(FL\) and \(FM\). We have
\[
FL\simeq \langle 1,25\rangle \simeq \langle 1,1\rangle
\]
over \(F\), since \(25\in F^{\times 2}\). Similarly,
\[
FM\simeq \langle 3,75\rangle \simeq \langle 3,3\rangle
\]
over \(F\), since \(75=3\cdot 25\). Both forms have dimension \(2\) and
determinant \(1\) in \(F^\times/F^{\times 2}\).

Using the convention
\[
s(\langle a_1,\dots,a_n\rangle)
=
\prod_{1\le i\le j\le n}(a_i,a_j)_F,
\]
we get
\[
s(\langle 1,1\rangle)=1.
\]
For \(\langle 3,3\rangle\),
\[
s(\langle 3,3\rangle)
=(3,3)_F^3=(3,3)_F=(3,-1)_F.
\]
Since \(-1\) is a square in \(\mathbb Q_5\), we have
\[
(3,-1)_F=1.
\]
Thus \(FL\) and \(FM\) have the same dimension, the same determinant, and the
same Hasse symbol. By the classification theorem for quadratic spaces over a
local field,
\[
FL\simeq FM.
\]

Next we find the Jordan splittings. For \(L\), we already have
\[
L=\langle 1\rangle\perp \langle 25\rangle.
\]
The first component has scale \(\mathcal O\), and the second has scale
\(25\mathcal O\). Hence this is a Jordan splitting.

For \(M\), from the diagonalization above,
\[
M\simeq \langle 3\rangle\perp \langle 75\rangle.
\]
Since \(3\in\mathcal O^\times\) and \(75=3\cdot 25\), the first component has
scale \(\mathcal O\), and the second has scale \(25\mathcal O\). Hence
\[
M=\langle 3\rangle\perp \langle 75\rangle
\]
is a Jordan splitting.

It remains to show that \(L\) and \(M\) are not isomorphic as
\(\mathcal O\)-lattices. The lattice \(L\) represents \(1\), since
\[
1=q(1,0).
\]
We show that \(M\) does not represent \(1\). In the diagonal basis of \(M\),
a vector has norm
\[
3x^2+75y^2,\qquad x,y\in\mathbb Z_5.
\]
If this were equal to \(1\), then reducing modulo \(5\) would give
\[
3\bar x^2\equiv 1\pmod 5.
\]
Thus
\[
\bar x^2\equiv 2\pmod 5,
\]
because \(3^{-1}\equiv 2\pmod 5\). But the nonzero squares modulo \(5\) are
\[
1,\ 4,
\]
so \(2\) is not a square modulo \(5\). This is impossible. Therefore \(M\)
does not represent \(1\).

Since an isomorphism of \(\mathcal O\)-lattices preserves represented values,
\(L\) and \(M\) cannot be isomorphic as quadratic \(\mathcal O\)-lattices.
\end{solution}


\subsection*{Exercise A.14.7}
Let $F = \mathbb{Q}_2$ and let $\mathcal{O}_F$ be the valuation ring of $F$. Find a Jordan splitting for the quadratic $\mathcal{O}_F$-lattice
\[ L = \langle 3 \rangle \perp \begin{pmatrix} 6 & 5 \\ 5 & 8 \end{pmatrix} \perp \begin{pmatrix} 4 & 3 \\ 3 & 2 \end{pmatrix}. \]
\begin{solution}
We write \(\mathcal O=\mathbb Z_2\) and \(F=\mathbb Q_2\). We use the following
criterion: if a lattice has Gram matrix \(A\in M_r(\mathcal O)\) with
\(\det(A)\in\mathcal O^\times\), then it is unimodular, i.e. \(\mathcal O\)-modular.

The lattice is given as an orthogonal direct sum
\[
L=\langle 3\rangle
\perp
\begin{pmatrix}
6&5\\
5&8
\end{pmatrix}
\perp
\begin{pmatrix}
4&3\\
3&2
\end{pmatrix}.
\]
We check each summand.

First,
\[
\det\langle 3\rangle =3\in\mathbb Z_2^\times,
\]
so \(\langle 3\rangle\) is unimodular.

For the second summand,
\[
\det
\begin{pmatrix}
6&5\\
5&8
\end{pmatrix}
=
6\cdot 8-5^2
=
48-25
=
23\in\mathbb Z_2^\times.
\]
Hence this binary lattice is also unimodular.

For the third summand,
\[
\det
\begin{pmatrix}
4&3\\
3&2
\end{pmatrix}
=
4\cdot 2-3^2
=
8-9
=
-1\in\mathbb Z_2^\times.
\]
Thus this binary lattice is also unimodular.

Therefore each orthogonal summand is \(\mathcal O\)-modular. Hence
\[
L=\langle 3\rangle
\perp
\begin{pmatrix}
6&5\\
5&8
\end{pmatrix}
\perp
\begin{pmatrix}
4&3\\
3&2
\end{pmatrix}
\]
is a Jordan splitting of \(L\), with all Jordan components of scale
\(\mathcal O\).

Equivalently, since all components have the same scale, one may combine them
and say that \(L\) itself is unimodular:
\[
L^\#=L.
\]
Thus, under the convention that Jordan components with the same scale are
combined, the Jordan splitting is simply
\[
L=L_0,
\qquad L_0=L,
\]
where \(L_0\) is \(\mathcal O\)-modular.
\end{solution}
\subsection*{Exercise A.14.8}
Let $F = \mathbb{Q}_2$ and let $\mathcal{O}_F$ be the valuation ring of $F$. Prove that as quadratic $\mathcal{O}_F$-lattices
\[ \langle 1 \rangle \perp \begin{pmatrix} 0 & 1 \\ 1 & 0 \end{pmatrix} \cong \langle 1 \rangle \perp \begin{pmatrix} 1 & 1 \\ 1 & 0 \end{pmatrix} \quad \text{but} \quad \begin{pmatrix} 0 & 1 \\ 1 & 0 \end{pmatrix} \not\cong \begin{pmatrix} 1 & 1 \\ 1 & 0 \end{pmatrix}. \]
\begin{solution}
Write \(\mathcal O=\mathbb Z_2\). Put
\[
H=
\begin{pmatrix}
0&1\\
1&0
\end{pmatrix},
\qquad
A=
\begin{pmatrix}
1&1\\
1&0
\end{pmatrix}.
\]
We prove the statement in two steps.

\emph{Step 1.} We construct an explicit \(\mathcal O\)-basis change showing
\[
\langle 1\rangle\perp H \simeq \langle 1\rangle\perp A.
\]

\emph{Step 2.} We show that \(H\not\simeq A\) by comparing the values represented
by the two lattices.

Let \(L=\langle 1\rangle\perp H\), with basis \(e,f,g\), so that
\[
B(e,e)=1,\qquad B(f,g)=1,
\]
and all other pairings among basis vectors are zero. Thus the Gram matrix of
\(L\) is
\[
G=
\begin{pmatrix}
1&0&0\\
0&0&1\\
0&1&0
\end{pmatrix}.
\]
Define
\[
u=e-g,\qquad v=e+f,\qquad w=g.
\]
Then \((u,v,w)\) is an \(\mathcal O\)-basis, since the change-of-basis matrix
from \((e,f,g)\) to \((u,v,w)\) is
\[
P=
\begin{pmatrix}
1&1&0\\
0&1&0\\
-1&0&1
\end{pmatrix},
\qquad \det(P)=1.
\]
Now compute the pairings:
\[
B(u,u)=1,\qquad B(u,v)=0,\qquad B(u,w)=0,
\]
and
\[
B(v,v)=1,\qquad B(v,w)=1,\qquad B(w,w)=0.
\]
Therefore the Gram matrix in the basis \((u,v,w)\) is
\[
\begin{pmatrix}
1&0&0\\
0&1&1\\
0&1&0
\end{pmatrix}
=
\langle 1\rangle\perp
\begin{pmatrix}
1&1\\
1&0
\end{pmatrix}.
\]
Hence
\[
\langle 1\rangle\perp
\begin{pmatrix}
0&1\\
1&0
\end{pmatrix}
\simeq
\langle 1\rangle\perp
\begin{pmatrix}
1&1\\
1&0
\end{pmatrix}.
\]

It remains to show that
\[
\begin{pmatrix}
0&1\\
1&0
\end{pmatrix}
\not\simeq
\begin{pmatrix}
1&1\\
1&0
\end{pmatrix}.
\]
For the lattice \(H\), a vector \((x,y)\in\mathcal O^2\) has norm
\[
q_H(x,y)=2xy.
\]
Thus every value represented by \(H\) lies in \(2\mathcal O\). In particular,
\(H\) does not represent any unit of \(\mathcal O\).

For the lattice \(A\), we have
\[
q_A(x,y)=x^2+2xy.
\]
Hence
\[
q_A(1,0)=1.
\]
So \(A\) represents the unit \(1\).

Since an isomorphism of quadratic \(\mathcal O\)-lattices preserves represented
values, \(H\) and \(A\) cannot be isomorphic. Therefore
\[
\begin{pmatrix}
0&1\\
1&0
\end{pmatrix}
\not\simeq
\begin{pmatrix}
1&1\\
1&0
\end{pmatrix}.
\]
\end{solution}
\section*{Homework 15}
\noindent Throughout what follows, let $F$ and $\mathcal{O}_F$ be as in \S 5.2, and let $(V,Q)$ be a nonsingular quadratic space of dimension $n \ge 1$ over $F$.

\subsection*{Exercise A.15.1}
Prove that the groups $\operatorname{O}_{\mathbf{A}}(V)$ and $\operatorname{O}_{\mathbf{A}}^+(V)$ defined in (5.2.3.1) and (5.2.3.2) are independent of the choice of the basis of $V$.
\begin{solution}
We prove the statement for \(O_{\mathbb A}(V)\); the proof for
\(O_{\mathbb A}^{+}(V)\) is identical.

Let \(e=(e_1,\dots,e_n)\) and \(e'=(e'_1,\dots,e'_n)\) be two \(F\)-bases of
\(V\). Let \(C\in \operatorname{GL}_n(F)\) be the change-of-basis matrix from
\(e'\) to \(e\). If \(\sigma_{\mathfrak p}\in O(V_{\mathfrak p})\), write
\(M_{\mathfrak p}\) for its matrix in the basis \(e\), and
\(M'_{\mathfrak p}\) for its matrix in the basis \(e'\). Then
\[
M'_{\mathfrak p}=C^{-1}M_{\mathfrak p}C .
\]

Since \(C\) and \(C^{-1}\) have entries in \(F\), there is a finite set of
places \(S\) such that, for every finite place \(\mathfrak p\notin S\), all
entries of \(C\) and \(C^{-1}\) are \(\mathfrak p\)-adic integers. Hence
\(\|C\|_{\mathfrak p}\le 1\) and \(\|C^{-1}\|_{\mathfrak p}\le 1\) for
\(\mathfrak p\notin S\).

Using the submultiplicativity of the matrix norm, for
\(\mathfrak p\notin S\) we get
\[
\|M'_{\mathfrak p}\|_{\mathfrak p}
=\|C^{-1}M_{\mathfrak p}C\|_{\mathfrak p}
\le \|M_{\mathfrak p}\|_{\mathfrak p}.
\]
Conversely, since \(M_{\mathfrak p}=CM'_{\mathfrak p}C^{-1}\), we also have
\[
\|M_{\mathfrak p}\|_{\mathfrak p}
\le \|M'_{\mathfrak p}\|_{\mathfrak p}.
\]
Therefore, for all \(\mathfrak p\notin S\),
\[
\|M_{\mathfrak p}\|_{\mathfrak p}=1
\quad\Longleftrightarrow\quad
\|M'_{\mathfrak p}\|_{\mathfrak p}=1 .
\]

Thus the condition \(\|\sigma_{\mathfrak p}\|_{\mathfrak p}=1\) for almost all
\(\mathfrak p\) is independent of the chosen basis. Since \(O(V_{\mathfrak p})\)
and \(O^{+}(V_{\mathfrak p})\) are intrinsically defined from the quadratic
space \((V_{\mathfrak p},Q)\), the adelic groups
\(O_{\mathbb A}(V)\) and \(O_{\mathbb A}^{+}(V)\) are independent of the choice
of basis of \(V\).
\end{solution}
\subsection*{Exercise A.15.2}
Let $L$ be a full $\mathcal{O}_F$-lattice in $V$. For any fractional ideal $\mathfrak{a}$ of $\mathcal{O}_F$ and any $\mathfrak{p} \in \operatorname{Spm}(\mathcal{O}_F)$, let $\mathfrak{a}_{\mathfrak{p}} = \mathfrak{a}\mathcal{O}_{\mathfrak{p}}$ be the fractional ideal of $\mathcal{O}_{\mathfrak{p}}$ generated by $\mathfrak{a}$.
\begin{enumerate}
    \item Prove that for every $\mathfrak{p} \in \operatorname{Spm}(\mathcal{O}_F)$ one has
    \[ \mathfrak{s}(L)_{\mathfrak{p}} = \mathfrak{s}(L_{\mathfrak{p}}) \,,\; \mathfrak{n}(L)_{\mathfrak{p}} = \mathfrak{n}(L_{\mathfrak{p}}) \quad \text{and} \quad \mathfrak{v}(L)_{\mathfrak{p}} = \mathfrak{v}(L_{\mathfrak{p}}) . \]
    \item Let $\mathfrak{a}$ be a fractional ideal of $\mathcal{O}_F$. \\
    Prove that $L$ is $\mathfrak{a}$-modular if and only if $L_{\mathfrak{p}}$ is $\mathfrak{a}_{\mathfrak{p}}$-modular for every $\mathfrak{p} \in \operatorname{Spm}(\mathcal{O}_F)$.
    \item Prove that $L$ is $\mathfrak{a}$-maximal in $V$ if and only if $L_{\mathfrak{p}}$ is $\mathfrak{a}_{\mathfrak{p}}$-maximal in $V_{\mathfrak{p}}$ for every $\mathfrak{p} \in \operatorname{Spm}(\mathcal{O}_F)$.
\end{enumerate}
\begin{solution}
We use the following standard facts in our textbook

\begin{enumerate}
    \item[(a)] For full lattices \(M,N\subset V\), one has \(M=N\) if and only if
    \(M_{\mathfrak p}=N_{\mathfrak p}\) for every \(\mathfrak p\in \operatorname{Spm}(\mathcal O_F)\).
    \item[(b)] If \(M_{\mathfrak p}\subset V_{\mathfrak p}\) is an \(\mathcal O_{\mathfrak p}\)-lattice for each
    \(\mathfrak p\), and \(M_{\mathfrak p}=L_{\mathfrak p}\) for almost all \(\mathfrak p\), then there is a unique
    full \(\mathcal O_F\)-lattice \(M\subset V\) such that \(M_{\mathfrak p}\) is its localization at every
    \(\mathfrak p\).
\end{enumerate}

\begin{enumerate}
    \item We first prove \( \mathfrak s(L)_{\mathfrak p}=\mathfrak s(L_{\mathfrak p}) \).
    By definition, \(\mathfrak s(L)\) is generated over \(\mathcal O_F\) by the values
    \(B(x,y)\) with \(x,y\in L\). After localizing at \(\mathfrak p\), \(\mathfrak s(L)_{\mathfrak p}\)
    is generated over \(\mathcal O_{\mathfrak p}\) by the same values, so
    \(\mathfrak s(L)_{\mathfrak p}\subseteq \mathfrak s(L_{\mathfrak p})\).

    Conversely, every element of \(L_{\mathfrak p}\) has the form \(x/u\) with
    \(x\in L\) and \(u\in \mathcal O_F\setminus \mathfrak p\). Hence every generator of
    \(\mathfrak s(L_{\mathfrak p})\) has the form
    \(B(x/u,y/v)=(uv)^{-1}B(x,y)\), where \(uv\) is a unit in \(\mathcal O_{\mathfrak p}\).
    Thus it lies in \(\mathfrak s(L)_{\mathfrak p}\). So
    \(\mathfrak s(L_{\mathfrak p})=\mathfrak s(L)_{\mathfrak p}\).

    The proof for \( \mathfrak n(L)_{\mathfrak p}=\mathfrak n(L_{\mathfrak p}) \) is the same:
    \(\mathfrak n(L)\) is generated by \(Q(x)\) for \(x\in L\), and
    \(Q(x/u)=u^{-2}Q(x)\) with \(u\) a unit in \(\mathcal O_{\mathfrak p}\).

    For the volume ideal, we cannot assume that \(L\) is free over
\(\mathcal O_F\). Choose a pseudo-basis of \(L\), say
\[
L=\mathfrak a_1 e_1+\cdots+\mathfrak a_n e_n,
\]
where \(e_1,\dots,e_n\) is an \(F\)-basis of \(V\), and
\(\mathfrak a_i\) are fractional ideals of \(F\). Then the volume ideal is
\[
\mathfrak v(L)
=
\mathfrak a_1^2\cdots \mathfrak a_n^2
\det(B(e_i,e_j))\mathcal O_F .
\]

Localizing at \(\mathfrak p\), we get
\[
L_{\mathfrak p}
=
(\mathfrak a_1)_{\mathfrak p}e_1+\cdots+
(\mathfrak a_n)_{\mathfrak p}e_n .
\]
Since \(\mathcal O_{\mathfrak p}\) is a DVR, every fractional ideal
\((\mathfrak a_i)_{\mathfrak p}\) is principal. Choose
\[
(\mathfrak a_i)_{\mathfrak p}=\alpha_i\mathcal O_{\mathfrak p}.
\]
Then
\[
\alpha_1e_1,\dots,\alpha_ne_n
\]
is an \(\mathcal O_{\mathfrak p}\)-basis of \(L_{\mathfrak p}\). Hence
\[
\mathfrak v(L_{\mathfrak p})
=
\det\bigl(B(\alpha_i e_i,\alpha_j e_j)\bigr)\mathcal O_{\mathfrak p}.
\]
But
\[
\det\bigl(B(\alpha_i e_i,\alpha_j e_j)\bigr)
=
(\alpha_1\cdots \alpha_n)^2\det(B(e_i,e_j)).
\]
Therefore
\[
\mathfrak v(L_{\mathfrak p})
=
(\mathfrak a_1)_{\mathfrak p}^2\cdots
(\mathfrak a_n)_{\mathfrak p}^2
\det(B(e_i,e_j))\mathcal O_{\mathfrak p}
=
\mathfrak v(L)_{\mathfrak p}.
\]

    \item We use the criterion that a full lattice \(M\) of rank \(n\) is
\(\mathfrak b\)-modular if and only if
\[
\mathfrak s(M)=\mathfrak b
\qquad\text{and}\qquad
\mathfrak v(M)=\mathfrak b^n .
\]

Suppose first that \(L\) is \(\mathfrak a\)-modular. Then
\[
\mathfrak s(L)=\mathfrak a
\qquad\text{and}\qquad
\mathfrak v(L)=\mathfrak a^n.
\]
Localizing at \(\mathfrak p\), and using part (1), we get
\[
\mathfrak s(L_{\mathfrak p})
=
\mathfrak s(L)_{\mathfrak p}
=
\mathfrak a_{\mathfrak p},
\]
and
\[
\mathfrak v(L_{\mathfrak p})
=
\mathfrak v(L)_{\mathfrak p}
=
(\mathfrak a^n)_{\mathfrak p}
=
\mathfrak a_{\mathfrak p}^n.
\]
Hence \(L_{\mathfrak p}\) is \(\mathfrak a_{\mathfrak p}\)-modular for every
\(\mathfrak p\).

Conversely, suppose that \(L_{\mathfrak p}\) is
\(\mathfrak a_{\mathfrak p}\)-modular for every
\(\mathfrak p\in \operatorname{Spm}(\mathcal O_F)\). Then
\[
\mathfrak s(L_{\mathfrak p})=\mathfrak a_{\mathfrak p}
\qquad\text{and}\qquad
\mathfrak v(L_{\mathfrak p})=\mathfrak a_{\mathfrak p}^n
\]
for every \(\mathfrak p\). By part (1),
\[
\mathfrak s(L)_{\mathfrak p}=\mathfrak a_{\mathfrak p}
\qquad\text{and}\qquad
\mathfrak v(L)_{\mathfrak p}=(\mathfrak a^n)_{\mathfrak p}
\]
for every \(\mathfrak p\).

Since two fractional ideals of \(\mathcal O_F\) are equal if and only if their
localizations at every maximal ideal are equal, we obtain
\[
\mathfrak s(L)=\mathfrak a
\qquad\text{and}\qquad
\mathfrak v(L)=\mathfrak a^n.
\]
Therefore \(L\) is \(\mathfrak a\)-modular.

    \item We prove that \(L\) is \(\mathfrak a\)-maximal in \(V\) if and only if
\(L_{\mathfrak p}\) is \(\mathfrak a_{\mathfrak p}\)-maximal in
\(V_{\mathfrak p}\) for every
\(\mathfrak p\in \operatorname{Spm}(\mathcal O_F)\).

Recall that \(L\) is \(\mathfrak a\)-maximal means that
\(\mathfrak n(L)\subseteq \mathfrak a\), and that there is no full lattice
\(M\subset V\) with
\[
L\subsetneq M
\qquad\text{and}\qquad
\mathfrak n(M)\subseteq \mathfrak a .
\]

Suppose first that \(L\) is \(\mathfrak a\)-maximal. Fix
\(\mathfrak p\in \operatorname{Spm}(\mathcal O_F)\). Let
\(M_{\mathfrak p}\subset V_{\mathfrak p}\) be a lattice such that
\[
L_{\mathfrak p}\subseteq M_{\mathfrak p}
\qquad\text{and}\qquad
\mathfrak n(M_{\mathfrak p})\subseteq \mathfrak a_{\mathfrak p}.
\]
For every \(\mathfrak q\neq \mathfrak p\), set
\[
M_{\mathfrak q}:=L_{\mathfrak q}.
\]
Then \(M_{\mathfrak q}=L_{\mathfrak q}\) for almost all
\(\mathfrak q\), so these local lattices come from a full lattice
\(M\subset V\). Moreover, since
\[
L_{\mathfrak q}\subseteq M_{\mathfrak q}
\]
for every \(\mathfrak q\), we get \(L\subseteq M\).

By part (1), for every \(\mathfrak q\),
\[
\mathfrak n(M)_{\mathfrak q}
=
\mathfrak n(M_{\mathfrak q}).
\]
For \(\mathfrak q=\mathfrak p\), this is contained in
\(\mathfrak a_{\mathfrak p}\) by assumption. For
\(\mathfrak q\neq \mathfrak p\), we have
\(M_{\mathfrak q}=L_{\mathfrak q}\), and since
\(\mathfrak n(L)\subseteq \mathfrak a\), part (1) gives
\[
\mathfrak n(M_{\mathfrak q})
=
\mathfrak n(L_{\mathfrak q})
=
\mathfrak n(L)_{\mathfrak q}
\subseteq
\mathfrak a_{\mathfrak q}.
\]
Thus
\[
\mathfrak n(M)_{\mathfrak q}\subseteq \mathfrak a_{\mathfrak q}
\]
for every \(\mathfrak q\). Since containment of fractional ideals can be checked
after localizing at every maximal ideal, it follows that
\[
\mathfrak n(M)\subseteq \mathfrak a.
\]
By the \(\mathfrak a\)-maximality of \(L\), we get \(M=L\). Localizing at
\(\mathfrak p\), we obtain
\[
M_{\mathfrak p}=L_{\mathfrak p}.
\]
Hence \(L_{\mathfrak p}\) is \(\mathfrak a_{\mathfrak p}\)-maximal.

Conversely, suppose that \(L_{\mathfrak p}\) is
\(\mathfrak a_{\mathfrak p}\)-maximal for every \(\mathfrak p\). First, by
part (1),
\[
\mathfrak n(L)_{\mathfrak p}
=
\mathfrak n(L_{\mathfrak p})
\subseteq
\mathfrak a_{\mathfrak p}
\]
for every \(\mathfrak p\). Hence
\[
\mathfrak n(L)\subseteq \mathfrak a.
\]

Now let \(M\subset V\) be a full lattice such that
\[
L\subseteq M
\qquad\text{and}\qquad
\mathfrak n(M)\subseteq \mathfrak a.
\]
Then for every \(\mathfrak p\),
\[
L_{\mathfrak p}\subseteq M_{\mathfrak p}.
\]
Again by part (1),
\[
\mathfrak n(M_{\mathfrak p})
=
\mathfrak n(M)_{\mathfrak p}
\subseteq
\mathfrak a_{\mathfrak p}.
\]
Since \(L_{\mathfrak p}\) is \(\mathfrak a_{\mathfrak p}\)-maximal, we get
\[
M_{\mathfrak p}=L_{\mathfrak p}
\]
for every \(\mathfrak p\). Since two full lattices are equal if and only if
their localizations at all maximal ideals are equal, it follows that
\[
M=L.
\]
Therefore \(L\) is \(\mathfrak a\)-maximal. 
\end{enumerate}
\end{solution}
\subsection*{Exercise A.15.3}
Let $\mathfrak{p}$ be a non-archimedean place of $F$. Prove that if $n = \dim V \ge 3$, then $\theta_{\mathfrak{p}}(\operatorname{O}^+(V_{\mathfrak{p}})) = F_{\mathfrak{p}}^\times$.
\begin{solution}
Put \(K=F_{\mathfrak p}\), and write \(q\) for the quadratic form on
\(V_{\mathfrak p}\). Let \(\tau_x\) denote the orthogonal symmetry attached to
a vector \(x\) with \(q(x)\neq 0\). Then \(\det(\tau_x)=-1\), so
\(\tau_x\tau_y\in O^+(V_{\mathfrak p})\), and the definition of the spinor norm
gives
\[
\theta_{\mathfrak p}(\tau_x\tau_y)
= q(x)q(y) \quad \text{in } K^\times/K^{\times 2}.
\]
Thus it is enough to prove that every square class in \(K^\times\) can be
written as a product of two nonzero values of \(q\).

Let \(a\in K^\times\). First suppose that \(q\) is isotropic. Since \(q\) is
nondegenerate, it contains a hyperbolic plane. A hyperbolic plane represents
every element of \(K^\times\). Hence we may choose \(x,y\in V_{\mathfrak p}\)
with \(q(x)=a\) and \(q(y)=1\). Then
\[
\theta_{\mathfrak p}(\tau_x\tau_y)=a
\quad \text{in } K^\times/K^{\times 2}.
\]

Now suppose that \(q\) is anisotropic. Since \(K\) is a non-archimedean local
field, every quadratic form over \(K\) of dimension at least \(5\) is isotropic.
The form
\[
q\perp (-a)q
\]
has dimension \(2n\ge 6\), hence is isotropic. Therefore there exist
\(x,y\in V_{\mathfrak p}\), not both zero, such that
\[
q(x)-a q(y)=0.
\]
Because \(q\) is anisotropic, we have \(x\neq 0\), \(y\neq 0\), and
\(q(x),q(y)\neq 0\). Thus \(q(x)=a q(y)\), so in the square-class group,
\[
a=\frac{q(x)}{q(y)}\equiv q(x)q(y)
\pmod{K^{\times 2}}.
\]
Hence
\[
\theta_{\mathfrak p}(\tau_x\tau_y)
=q(x)q(y)
\equiv a
\pmod{K^{\times 2}}.
\]

Since \(a\in K^\times\) was arbitrary, every square class occurs as the spinor
norm of an element of \(O^+(V_{\mathfrak p})\). Therefore
\[
\theta_{\mathfrak p}\bigl(O^+(V_{\mathfrak p})\bigr)
=K^\times/K^{\times 2}
=F_{\mathfrak p}^{\times}/F_{\mathfrak p}^{\times 2}.
\]
\end{solution}
\subsection*{Exercise A.15.4}
Let $\mathfrak{p}$ be a non-dyadic place of $F$. Prove that if $n = \dim V \ge 2$ and $L$ is a modular full lattice in $V$, then $\theta_{\mathfrak{p}}(\operatorname{O}^+(L_{\mathfrak{p}})) = \mathcal{O}_{\mathfrak{p}}^\times F_{\mathfrak{p}}^{\times 2}$.
\begin{solution}
Put \(K=F_{\mathfrak p}\), \(\mathcal O=\mathcal O_{\mathfrak p}\), and
\(\Lambda=L_{\mathfrak p}\). Let \(\pi\) be a uniformizer of \(\mathcal O\).
Since \(\mathfrak p\) is non-dyadic, \(2\in\mathcal O^\times\). Hence the norm
ideal and the scale ideal of \(\Lambda\) are equal:
\[
\mathfrak n(\Lambda)=\mathfrak s(\Lambda).
\]
Write
\[
\mathfrak n(\Lambda)=\mathfrak s(\Lambda)=\pi^r\mathcal O.
\]

We first prove that
\[
\theta_{\mathfrak p}(O^+(\Lambda))
\subseteq \mathcal O^\times K^{\times 2}.
\]
Let \(\sigma\in O^+(\Lambda)\). By Theorem 4.3.9, \(\sigma\) is a product of
symmetries in \(O(\Lambda)\):
\[
\sigma=\tau_{x_1}\cdots \tau_{x_m}.
\]
Since each symmetry has determinant \(-1\), and \(\sigma\in O^+(\Lambda)\), the
integer \(m\) is even.

We claim that for each symmetry \(\tau_x\in O(\Lambda)\), after replacing \(x\)
by a primitive generator of the line \(Kx\cap\Lambda\), one has
\[
Q(x)\mathcal O=\mathfrak n(\Lambda)=\pi^r\mathcal O.
\]
Indeed, since \(\tau_x\) preserves \(\Lambda\), for every \(y\in\Lambda\),
\[
\tau_x(y)=y-\frac{B(y,x)}{Q(x)}x\in \Lambda.
\]
Thus \(\frac{B(y,x)}{Q(x)}x\in\Lambda\). Since \(x\) is primitive, this implies
\[
\frac{B(y,x)}{Q(x)}\in\mathcal O.
\]
Hence \(B(x,\Lambda)\subseteq Q(x)\mathcal O\).

On the other hand, \(\Lambda\) is modular, so the pairing
\(\pi^{-r}B\) is unimodular on \(\Lambda\). Therefore, for primitive \(x\), one
has
\[
B(x,\Lambda)=\pi^r\mathcal O.
\]
So
\[
\pi^r\mathcal O=B(x,\Lambda)\subseteq Q(x)\mathcal O.
\]
But \(x\in\Lambda\), so \(Q(x)\in\mathfrak n(\Lambda)=\pi^r\mathcal O\). Hence
also
\[
Q(x)\mathcal O\subseteq \pi^r\mathcal O.
\]
Therefore
\[
Q(x)\mathcal O=\pi^r\mathcal O.
\]

Now return to \(\sigma=\tau_{x_1}\cdots\tau_{x_m}\). Since \(m\) is even, say
\(m=2s\), and since every \(Q(x_i)\) lies in \(\pi^r\mathcal O^\times\), we get
\[
\theta_{\mathfrak p}(\sigma)
=Q(x_1)\cdots Q(x_{2s})
\in (\pi^r\mathcal O^\times)^{2s}K^{\times 2}
=\mathcal O^\times K^{\times 2}.
\]
Thus
\[
\theta_{\mathfrak p}(O^+(\Lambda))
\subseteq \mathcal O^\times K^{\times 2}.
\]

Conversely, let \(u\in \mathcal O^\times\). Since \(\Lambda\) is non-dyadic
modular and \(\dim\Lambda\ge 2\), we may choose an orthogonal binary sublattice
\(\mathcal O e_1\oplus \mathcal O e_2\subseteq \Lambda\) such that
\[
Q(e_1),Q(e_2)\in \pi^r\mathcal O^\times .
\]
Write
\[
Q(e_1)=\pi^r a_1,\qquad a_1\in \mathcal O^\times.
\]
On the binary sublattice, set
\[
Q'=\pi^{-r}Q.
\]
Then \(Q'\) is a unimodular binary quadratic form over \(\mathcal O\). Reducing
modulo \(\pi\), we get a nondegenerate binary quadratic form \(\overline{Q'}\)
over the residue field \(k\).

Let
\[
t=ua_1^{-1}\in \mathcal O^\times.
\]
Since \(k\) has odd characteristic, every nondegenerate binary quadratic form
over \(k\) represents every nonzero element of \(k\). Hence there exists
\(\bar x\in k^2\) such that
\[
\overline{Q'}(\bar x)=\bar t.
\]
Since \(\bar t\neq 0\), the vector \(\bar x\) is nonzero. Moreover, because
\(\overline{Q'}\) is nondegenerate, the gradient of
\(\overline{Q'}-\bar t\) at \(\bar x\) is nonzero. Therefore Hensel's lemma
lifts \(\bar x\) to a vector \(x\in \mathcal O e_1\oplus \mathcal O e_2\) such
that
\[
Q'(x)=t.
\]
Thus
\[
Q(x)=\pi^r t\in \pi^r\mathcal O^\times.
\]
In particular, \(x\) is primitive. By the claim above, both \(\tau_x\) and
\(\tau_{e_1}\) preserve \(\Lambda\). Hence
\[
\tau_x\tau_{e_1}\in O^+(\Lambda).
\]
Its spinor norm is
\[
\theta_{\mathfrak p}(\tau_x\tau_{e_1})
=
Q(x)Q(e_1)
=
(\pi^r t)(\pi^r a_1)
=
\pi^{2r}u
\equiv u
\pmod{K^{\times 2}}.
\]
Therefore every class in \(\mathcal O^\times K^{\times 2}\) occurs as the
spinor norm of an element of \(O^+(\Lambda)\).
\end{solution}

\subsection*{Exercise A.15.5}
Put
\[ \mathbf{I}_{F, S_0} := \prod_{\mathfrak{p} \in S_0} F_{\mathfrak{p}}^\times \times \prod_{\mathfrak{p} \notin S_0} \mathcal{O}_{\mathfrak{p}}^\times . \]
Prove that there is a group isomorphism
\[ \mathbf{I}_F / F^\times \mathbf{I}_{F, S_0} \cong \mathcal{C}l(\mathcal{O}_F) . \]
(Recall that $\mathcal{C}l(\mathcal{O}_F)$ denotes the ideal class group of $\mathcal{O}_F$.)
\begin{solution}
Define a map
\[
\Phi:I_F\longrightarrow \operatorname{Cl}(\mathcal O_F)
\]
as follows. For an idele \(a=(a_{\mathfrak p})_{\mathfrak p}\in I_F\), set
\[
\Phi(a)=\left[\prod_{\mathfrak p\notin S_0}
\mathfrak p^{v_{\mathfrak p}(a_{\mathfrak p})}\right].
\]
This product is finite, because by definition of the idele group,
\(a_{\mathfrak p}\in\mathcal O_{\mathfrak p}^{\times}\) for almost all finite
places \(\mathfrak p\). Hence \(v_{\mathfrak p}(a_{\mathfrak p})=0\) for almost
all \(\mathfrak p\).

The map \(\Phi\) is a group homomorphism, since valuations are additive:
\[
v_{\mathfrak p}(a_{\mathfrak p}b_{\mathfrak p})
=
v_{\mathfrak p}(a_{\mathfrak p})+v_{\mathfrak p}(b_{\mathfrak p}).
\]

We first show that \(\Phi\) is surjective. Let \(\mathfrak a\) be a fractional
ideal of \(\mathcal O_F\). Write
\[
\mathfrak a=\prod_{\mathfrak p\notin S_0}\mathfrak p^{n_{\mathfrak p}},
\]
where all but finitely many \(n_{\mathfrak p}\) are zero. Choose a uniformizer
\(\pi_{\mathfrak p}\in F_{\mathfrak p}\) for each finite place
\(\mathfrak p\), and define an idele \(a=(a_{\mathfrak p})\) by
\(a_{\mathfrak p}=\pi_{\mathfrak p}^{n_{\mathfrak p}}\) for finite
\(\mathfrak p\), and \(a_{\mathfrak p}=1\) for \(\mathfrak p\in S_0\). Then
\[
\Phi(a)=[\mathfrak a].
\]
Thus \(\Phi\) is surjective.

Now compute the kernel. If \(u\in I_{F,S_0}\), then for every finite place
\(\mathfrak p\notin S_0\), one has \(u_{\mathfrak p}\in
\mathcal O_{\mathfrak p}^{\times}\), hence
\(v_{\mathfrak p}(u_{\mathfrak p})=0\). Therefore \(\Phi(u)=1\). Also, if
\(x\in F^\times\) is diagonally embedded into \(I_F\), then
\[
\prod_{\mathfrak p\notin S_0}
\mathfrak p^{v_{\mathfrak p}(x)}
=(x),
\]
so \(\Phi(x)=1\) in the ideal class group. Hence
\[
F^\times I_{F,S_0}\subseteq \ker \Phi.
\]

Conversely, suppose \(a=(a_{\mathfrak p})\in\ker\Phi\). Then
\[
\prod_{\mathfrak p\notin S_0}
\mathfrak p^{v_{\mathfrak p}(a_{\mathfrak p})}
\]
is principal, say equal to \((x)\) for some \(x\in F^\times\). Therefore
\[
v_{\mathfrak p}(a_{\mathfrak p})
=
v_{\mathfrak p}(x)
\]
for every finite place \(\mathfrak p\notin S_0\). Hence
\[
v_{\mathfrak p}(a_{\mathfrak p}x^{-1})=0
\]
for every finite \(\mathfrak p\notin S_0\), so
\[
a x^{-1}\in I_{F,S_0}.
\]
Thus \(a\in F^\times I_{F,S_0}\), and we get
\[
\ker\Phi=F^\times I_{F,S_0}.
\]

By the first isomorphism theorem,
\[
I_F/F^\times I_{F,S_0}\cong \operatorname{Cl}(\mathcal O_F).
\]
\end{solution}


\subsection*{Exercise A.15.6}
Prove that the condition $\mathbf{I}_F = \theta(\operatorname{O}^+(V))\mathbf{I}_{F, S_0} \mathbf{I}_F^2$ in Thm. 5.2.12 holds if $\mathcal{O}_F = \mathbb{Z}$ and $\dim V \ge 3$.
\begin{solution}
Since \(\mathcal O_F=\mathbb Z\), we have \(F=\mathbb Q\). Also
\(S_0=\{\infty\}\). By Exercise A.15.5 and the fact that
\(\operatorname{Cl}(\mathbb Z)=1\), we have
\[
I_{\mathbb Q}=\mathbb Q^\times I_{\mathbb Q,S_0}.
\]
Thus it is enough to show that every positive rational square class belongs to
\(\theta(O^+(V))\).

Let \(a\in\mathbb Q_{>0}^{\times}\). We prove that
\(a\in \theta(O^+(V))\) modulo \(\mathbb Q^{\times 2}\).

First suppose that \(V\) is isotropic over \(\mathbb Q\). Then \(V\) contains a
hyperbolic plane. A hyperbolic plane represents every element of
\(\mathbb Q\), so we may choose \(x,y\in V\) such that
\[
Q(x)=a,\qquad Q(y)=1.
\]
Then \(\tau_x\tau_y\in O^+(V)\), and
\[
\theta(\tau_x\tau_y)=Q(x)Q(y)=a
\quad \text{in } \mathbb Q^\times/\mathbb Q^{\times 2}.
\]

Now suppose that \(V\) is anisotropic over \(\mathbb Q\). Consider the quadratic
space
\[
V\oplus V,\qquad (u,v)\longmapsto Q(u)-aQ(v).
\]
Its dimension is \(2n\ge 6\). For every finite prime \(p\), every quadratic
form over \(\mathbb Q_p\) of dimension at least \(5\) is isotropic. Hence
\(Q\oplus(-aQ)\) is isotropic over \(\mathbb Q_p\) for every finite \(p\).
Over \(\mathbb R\), it is also isotropic because \(a>0\): if \(v\in V_\mathbb R\)
with \(Q(v)\neq 0\), then
\[
Q(\sqrt a\,v)-aQ(v)=0.
\]
Therefore, by the Hasse--Minkowski theorem, \(Q\oplus(-aQ)\) is isotropic over
\(\mathbb Q\). Hence there exist \(x,y\in V\), not both zero, such that
\[
Q(x)=aQ(y).
\]
Since \(V\) is anisotropic, \(x,y\neq 0\) and \(Q(x),Q(y)\neq 0\). Thus
\[
\theta(\tau_x\tau_y)=Q(x)Q(y)=aQ(y)^2
\equiv a \pmod{\mathbb Q^{\times 2}}.
\]
So every positive rational square class lies in \(\theta(O^+(V))\).

Now let \(t\in I_{\mathbb Q}\). Since
\(I_{\mathbb Q}=\mathbb Q^\times I_{\mathbb Q,S_0}\), write
\[
t=b u,\qquad b\in\mathbb Q^\times,\quad u\in I_{\mathbb Q,S_0}.
\]
Write \(b=\varepsilon a\), where \(\varepsilon=\pm 1\) and
\(a\in\mathbb Q_{>0}^{\times}\). Since \(\varepsilon\) is a \(p\)-adic unit for
every finite prime \(p\), we have \(\varepsilon\in I_{\mathbb Q,S_0}\). Hence
\[
t=a(\varepsilon u),
\qquad \varepsilon u\in I_{\mathbb Q,S_0}.
\]
By the previous paragraph, \(a\in \theta(O^+(V)) I_{\mathbb Q}^2\). Therefore
\[
t\in \theta(O^+(V))I_{\mathbb Q,S_0}I_{\mathbb Q}^2.
\]
Since \(t\in I_{\mathbb Q}\) was arbitrary, we get
\[
I_{\mathbb Q}
=
\theta(O^+(V))I_{\mathbb Q,S_0}I_{\mathbb Q}^2.
\]
\end{solution}
\subsection*{Exercise A.15.7}
Let $\mathfrak{p}$ be a non-archimedean place of $F$. Let $L$ be a maximal lattice in $V$.
\begin{enumerate}
    \item Prove that $\operatorname{O}(L_{\mathfrak{p}}) = \operatorname{O}(V_{\mathfrak{p}})$ if $V_{\mathfrak{p}}$ is anisotropic.
    \item Prove that if $\dim V \ge 3$, then $\theta_{\mathfrak{p}}(\operatorname{O}^+(L_{\mathfrak{p}})) \supseteq \mathcal{O}_{\mathfrak{p}}^\times$.
\end{enumerate}
\begin{solution}
Put \(K=F_{\mathfrak p}\), \(\mathcal O=\mathcal O_{\mathfrak p}\), and
\(\Lambda=L_{\mathfrak p}\).

\begin{enumerate}
    \item Suppose first that \(V_{\mathfrak p}\) is anisotropic. Let
    \[
    \Lambda_0:=\{x\in V_{\mathfrak p}\mid Q(x)\in \mathcal O\}.
    \]
    Since \(V_{\mathfrak p}\) is anisotropic, \(\Lambda_0\) is an
    \(\mathcal O\)-lattice. It is integral, and every integral lattice in
    \(V_{\mathfrak p}\) is contained in \(\Lambda_0\). Hence \(\Lambda_0\) is
    the unique maximal lattice in \(V_{\mathfrak p}\). Since \(\Lambda\) is
    maximal, we have \(\Lambda=\Lambda_0\).

    Now let \(\sigma\in O(V_{\mathfrak p})\). Since \(\sigma\) preserves \(Q\),
    one has
    \[
    Q(\sigma x)=Q(x)
    \]
    for every \(x\in V_{\mathfrak p}\). Therefore
    \[
    x\in \Lambda_0
    \quad\Longleftrightarrow\quad
    \sigma x\in \Lambda_0.
    \]
    Thus \(\sigma(\Lambda)=\Lambda\), so \(\sigma\in O(\Lambda)\). Hence
    \[
    O(\Lambda)=O(V_{\mathfrak p}).
    \]

    \item We prove that
    \[
    \theta_{\mathfrak p}(O^+(\Lambda))\supseteq \mathcal O^\times.
    \]

    If \(V_{\mathfrak p}\) is anisotropic, then by part (1),
    \[
    O^+(\Lambda)=O^+(V_{\mathfrak p}).
    \]
    Since \(\dim V\ge 3\), Exercise A.15.3 gives
    \[
    \theta_{\mathfrak p}(O^+(V_{\mathfrak p}))
    =
    K^\times/K^{\times 2}.
    \]
    Hence in particular
    \[
    \theta_{\mathfrak p}(O^+(\Lambda))
    \supseteq \mathcal O^\times K^{\times 2}/K^{\times 2}.
    \]

    It remains to treat the isotropic case. Since \(V_{\mathfrak p}\) is
    isotropic and \(\Lambda\) is maximal, the local splitting theorem for
    maximal lattices gives an orthogonal decomposition
    \[
    \Lambda=H\perp \Lambda',
    \]
    where \(H=\mathcal Oe\oplus \mathcal Of\) is a hyperbolic plane with
    \[
    Q(e)=Q(f)=0,
    \qquad
    B(e,f)=1.
    \]
    Hence
    \[
    Q(ae+bf)=ab.
    \]

    Let \(u\in\mathcal O^\times\). Set
    \[
    x=e+uf,
    \qquad
    y=e+f.
    \]
    Then
    \[
    Q(x)=u,
    \qquad
    Q(y)=1.
    \]
    Since \(Q(x)\) and \(Q(y)\) are units, and since \(B(\Lambda,x)\subseteq
    \mathcal O\) and \(B(\Lambda,y)\subseteq\mathcal O\), the symmetries
    \(\tau_x\) and \(\tau_y\) preserve \(\Lambda\). Thus
    \[
    \tau_x,\tau_y\in O(\Lambda),
    \qquad
    \tau_x\tau_y\in O^+(\Lambda).
    \]
    The spinor norm is
    \[
    \theta_{\mathfrak p}(\tau_x\tau_y)
    =
    Q(x)Q(y)
    =
    u
    \quad \text{in } K^\times/K^{\times 2}.
    \]
    Since \(u\in\mathcal O^\times\) was arbitrary, we obtain
    \[
    \theta_{\mathfrak p}(O^+(\Lambda))
    \supseteq \mathcal O^\times.
    \]
\end{enumerate}
\end{solution}
\subsection*{Exercise A.15.8}
Let $U, W$ be nonsingular quadratic spaces over $F$. Prove that there is a representation of quadratic spaces $U \to W$ over $F$ if and only if $U_{\mathfrak{p}} \to W_{\mathfrak{p}}$ for every $\mathfrak{p} \in \Omega_F$.
\begin{solution}
The necessity is immediate. If \(U\to W\) over \(F\), then after scalar
extension to every completion \(F_{\mathfrak p}\), we get
\[
U_{\mathfrak p}\to W_{\mathfrak p}.
\]

We prove the converse by induction on \(m=\dim U\). The case \(m=0\) is
trivial. Assume \(m\ge 1\), and suppose that
\[
U_{\mathfrak p}\to W_{\mathfrak p}
\]
for every place \(\mathfrak p\).

Suppose that $U$ represent $a$, then we can assume that 
\begin{align*}
    U \simeq \langle a \rangle \perp U'.
\end{align*}
since $\langle a \rangle$ is not singular. Hence, for every $\mathfrak{p}$ we can assume that 
\begin{align*}
    W_{\mathfrak{p}} \simeq\langle a \rangle \perp U'_{\mathfrak{p}} \perp W'_{\mathfrak{p}}.
\end{align*}
By Hasse–Minkowski theorem we obtain that $W$ represents $a$ then we can suppose that 
\begin{align*}
    W\simeq \langle a \rangle \perp  W''.
\end{align*}
Therefore, we have 
\begin{align*}
    \langle a \rangle \perp  W''_{\mathfrak{p}} \simeq \langle a \rangle \perp U'_{\mathfrak{p}} \perp W'_{\mathfrak{p}}.
\end{align*}
By Witt's cancellation, we have 
\begin{align*}
    W''_{\mathfrak{p}} \simeq U'_{\mathfrak{p}} \perp W'_{\mathfrak{p}}
\end{align*}
It menas that $W''$ represent $U'$ at every place. Hence, by Hasse–Minkowski theroem, we obtain that $W''$ represent $U'$ by our induction.
Therefore, we have $W$ represents $U$.
\end{solution}
\subsection*{Exercise A.15.9}
Prove the following weak approximation theorem for $\operatorname{O}^+(V)$: Let $T$ be a finite subset of $\Omega_F$. Suppose that an element $\varphi_{\mathfrak{p}} \in \operatorname{O}^+(V_{\mathfrak{p}})$ is given for each $\mathfrak{p} \in T$. Then for any real number $\varepsilon > 0$ there exists $\sigma \in \operatorname{O}^+(V)$ such that
\[ \| \sigma - \varphi_{\mathfrak{p}} \|_{\mathfrak{p}} < \varepsilon \quad \text{for every } \mathfrak{p} \in T . \]
(\textit{Hint}: Use the Cartan--Dieudonné Theorem 2.3.2).
\begin{solution}
Fix an \(F\)-basis of \(V\), and use it to identify \(V\) with \(F^n\) and
\(\operatorname{End}(V)\) with \(F^{n^2}\). Let \(Q\) be the quadratic form on
\(V\), and let \(\tau_x\) denote the symmetry attached to a vector \(x\) with
\(Q(x)\neq 0\). Thus
\[
\tau_x(y)=y-\frac{B(y,x)}{Q(x)}x .
\]
For every place \(\mathfrak p\), the map
\[
V_{\mathfrak p}\setminus \{Q=0\}\longrightarrow O(V_{\mathfrak p}),
\qquad
x\longmapsto \tau_x
\]
is continuous.

Let \(T\subset \Omega_F\) be finite, and suppose that
\(\varphi_{\mathfrak p}\in O^+(V_{\mathfrak p})\) is given for every
\(\mathfrak p\in T\). By the Cartan--Dieudonne theorem, each
\(\varphi_{\mathfrak p}\) is a product of symmetries. Since
\(\varphi_{\mathfrak p}\in O^+(V_{\mathfrak p})\), it may be written as an even
product:
\[
\varphi_{\mathfrak p}
=
\tau_{x_{\mathfrak p,1}}\tau_{x_{\mathfrak p,2}}
\cdots
\tau_{x_{\mathfrak p,2r_{\mathfrak p}}},
\qquad
Q(x_{\mathfrak p,i})\neq 0.
\]
Choose \(r\) larger than all \(r_{\mathfrak p}\). By inserting trivial pairs
\(\tau_z\tau_z=1\), we may assume that all decompositions have the same length
\(2r\):
\[
\varphi_{\mathfrak p}
=
\tau_{x_{\mathfrak p,1}}\cdots \tau_{x_{\mathfrak p,2r}} .
\]

Consider the product map
\[
\mu_{\mathfrak p}:
\bigl(V_{\mathfrak p}\setminus\{Q=0\}\bigr)^{2r}
\longrightarrow O(V_{\mathfrak p}),
\qquad
(y_1,\dots,y_{2r})
\longmapsto
\tau_{y_1}\cdots \tau_{y_{2r}} .
\]
This map is continuous. Hence, for each \(\mathfrak p\in T\), there exists a
neighborhood of
\[
(x_{\mathfrak p,1},\dots,x_{\mathfrak p,2r})
\]
such that if \(y_i\) lies in that neighborhood, then
\[
\left\|
\tau_{y_1}\cdots\tau_{y_{2r}}-\varphi_{\mathfrak p}
\right\|_{\mathfrak p}<\varepsilon .
\]

By weak approximation for the vector space \(V\), we may choose global vectors
\(x_1,\dots,x_{2r}\in V\) such that, for every \(\mathfrak p\in T\), each
\(x_i\) is sufficiently close to \(x_{\mathfrak p,i}\). Since \(Q(x_{\mathfrak p,i})\neq 0\)
and the condition \(Q(x_i)\neq 0\) is open locally, we may also assume
\[
Q(x_i)\neq 0
\]
for every \(i\).

Define
\[
\sigma=\tau_{x_1}\tau_{x_2}\cdots\tau_{x_{2r}}\in O(V).
\]
Since \(\sigma\) is a product of an even number of symmetries, we have
\[
\sigma\in O^+(V).
\]
By the choice of the \(x_i\)'s and the continuity of the product map, for every
\(\mathfrak p\in T\),
\[
\|\sigma-\varphi_{\mathfrak p}\|_{\mathfrak p}<\varepsilon .
\]
This proves the weak approximation theorem for \(O^+(V)\).
\end{solution}
\subsection*{Exercise A.15.10}
Let $K$ and $L$ be full lattices in $V$ such that $\operatorname{gen}(K) = \operatorname{gen}(L)$. Let $T$ be a finite subset of $\operatorname{Spm}(\mathcal{O}_F)$. Prove that there is a full lattice $K'$ in $V$ such that
\[ K' \in \operatorname{cls}^+(K) \quad \text{and} \quad K'_{\mathfrak{p}} = L_{\mathfrak{p}} \text{ for all } \mathfrak{p} \in T . \]
\begin{solution}
Since \(\operatorname{gen}(K)=\operatorname{gen}(L)\), for every
\(\mathfrak p\in T\) we may choose
\[
\varphi_{\mathfrak p}\in O^+(V_{\mathfrak p})
\]
such that
\[
\varphi_{\mathfrak p}K_{\mathfrak p}=L_{\mathfrak p}.
\]

For each \(\mathfrak p\in T\), the stabilizer
\[
O(K_{\mathfrak p})
=
\{g\in O(V_{\mathfrak p})\mid gK_{\mathfrak p}=K_{\mathfrak p}\}
\]
is an open subgroup of \(O(V_{\mathfrak p})\). Hence
\[
\varphi_{\mathfrak p}O(K_{\mathfrak p})
\]
is an open neighborhood of \(\varphi_{\mathfrak p}\) in \(O(V_{\mathfrak p})\).
By the weak approximation theorem for \(O^+(V)\), there exists
\[
\sigma\in O^+(V)
\]
such that, for every \(\mathfrak p\in T\),
\[
\sigma\in \varphi_{\mathfrak p}O(K_{\mathfrak p}).
\]
Equivalently, for every \(\mathfrak p\in T\),
\[
\varphi_{\mathfrak p}^{-1}\sigma\in O(K_{\mathfrak p}).
\]
Therefore
\[
\sigma K_{\mathfrak p}
=
\varphi_{\mathfrak p}K_{\mathfrak p}
=
L_{\mathfrak p}
\]
for every \(\mathfrak p\in T\).

Now set
\[
K'=\sigma K.
\]
Since \(\sigma\in O^+(V)\), we have
\[
K'\in \operatorname{cls}^+(K).
\]
Moreover, for every \(\mathfrak p\in T\),
\[
K'_{\mathfrak p}
=
(\sigma K)_{\mathfrak p}
=
\sigma K_{\mathfrak p}
=
L_{\mathfrak p}.
\]
Thus the required lattice \(K'\) exists.
\end{solution}
\subsection*{Exercise A.15.11}
Let $L$ be a full lattice in $V$ and let $K \subseteq V$ be a nonsingular lattice. Suppose that for every $\mathfrak{p} \in \operatorname{Spm}(\mathcal{O}_F)$, there is a representation $K_{\mathfrak{p}} \to L_{\mathfrak{p}}$. Prove that there is a representation $K \to L'$ of $K$ into a lattice $L'$ in $\operatorname{gen}(L)$.
\begin{solution}
Let \(U=FK\). Since \(K\subset V\) is nonsingular, \(U\) is a nonsingular
subspace of \(V\), and \(K\) is a full lattice in \(U\).

First note that
\[
K_{\mathfrak p}\subseteq L_{\mathfrak p}
\]
for almost all \(\mathfrak p\). Indeed, after choosing finitely many generators
of \(K\) and an \(\mathcal O_F\)-basis of \(L\), only finitely many primes occur
in the denominators of these generators. Hence outside a finite set of primes,
all elements of \(K\) are integral with respect to \(L\).

Let \(S\subset \operatorname{Spm}(\mathcal O_F)\) be a finite set containing
all primes for which \(K_{\mathfrak p}\not\subseteq L_{\mathfrak p}\). For
\(\mathfrak p\notin S\), set
\[
L'_{\mathfrak p}=L_{\mathfrak p}.
\]

Now fix \(\mathfrak p\in S\). By assumption, there is a representation
\[
\rho_{\mathfrak p}:K_{\mathfrak p}\longrightarrow L_{\mathfrak p}.
\]
Extending scalars, \(\rho_{\mathfrak p}\) gives an isometric embedding
\[
U_{\mathfrak p}\longrightarrow V_{\mathfrak p}.
\]
Since \(U_{\mathfrak p}\subset V_{\mathfrak p}\) is already a nonsingular
subspace, Witt's extension theorem gives an element
\[
\sigma_{\mathfrak p}\in O(V_{\mathfrak p})
\]
such that
\[
\sigma_{\mathfrak p}|_{U_{\mathfrak p}}=\rho_{\mathfrak p}.
\]
Hence
\[
\sigma_{\mathfrak p}K_{\mathfrak p}
=
\rho_{\mathfrak p}K_{\mathfrak p}
\subseteq L_{\mathfrak p}.
\]
Set
\[
L'_{\mathfrak p}:=\sigma_{\mathfrak p}^{-1}L_{\mathfrak p}.
\]
Then
\[
K_{\mathfrak p}\subseteq L'_{\mathfrak p},
\]
and \(L'_{\mathfrak p}\) is isometric to \(L_{\mathfrak p}\).

Thus we have constructed local lattices \(L'_{\mathfrak p}\subset V_{\mathfrak p}\)
such that \(L'_{\mathfrak p}=L_{\mathfrak p}\) for almost all \(\mathfrak p\),
and such that \(L'_{\mathfrak p}\cong L_{\mathfrak p}\) for every
\(\mathfrak p\). Define
\[
L'=\{x\in V\mid x\in L'_{\mathfrak p}\text{ for every }
\mathfrak p\in\operatorname{Spm}(\mathcal O_F)\}.
\]
By the standard localization property of lattices, \(L'\) is a full lattice in
\(V\), and its localization at \(\mathfrak p\) is precisely
\(L'_{\mathfrak p}\).

Since \(L'_{\mathfrak p}\cong L_{\mathfrak p}\) for every \(\mathfrak p\), we
have
\[
L'\in \operatorname{gen}(L).
\]
Moreover, for every \(\mathfrak p\),
\[
K_{\mathfrak p}\subseteq L'_{\mathfrak p}.
\]
Therefore \(K\subseteq L'\). Hence the inclusion gives a representation
\[
K\longrightarrow L'
\]
with \(L'\in \operatorname{gen}(L)\), as required.
\end{solution}
\subsection*{Exercise A.15.12}
Assume $\mathcal{O}_F = \mathbb{Z}$. Let $L = \mathbb{Z}e_1 \perp \mathbb{Z}e_2$ in the quadratic space $V = \mathbb{Q}e_1 \perp \mathbb{Q}e_2$ with $Q(e_1) = 1$ and $Q(e_2) = 11$.
\begin{enumerate}
    \item Prove that $3 \in Q(V)$ and $3 \in Q(L_{\mathfrak{p}})$ for every $\mathfrak{p} \in \operatorname{Spm}(\mathbb{Z})$.
    \item Prove that $3 \notin Q(L)$.
\end{enumerate}
\begin{solution}
Here \(V=\mathbb Q e_1\perp \mathbb Q e_2\) has quadratic form
\[
Q(xe_1+ye_2)=x^2+11y^2,
\]
and \(L=\mathbb Z e_1\perp \mathbb Z e_2\).

First, \(3\in Q(V)\). Indeed,
\[
Q\left(\frac{1}{2}e_1+\frac{1}{2}e_2\right)
=
\frac{1}{4}+\frac{11}{4}=3.
\]

We now prove that \(3\in Q(L_p)\) for every prime \(p\), where
\(L_p=\mathbb Z_p e_1\perp \mathbb Z_p e_2\).

For \(p=2\), it is enough to show that \(3/11\) is a square in
\(\mathbb Z_2^\times\). Since
\[
\frac{3}{11}\equiv 3\cdot 3\equiv 1 \pmod 8,
\]
and a \(2\)-adic unit is a square if and only if it is congruent to \(1\)
modulo \(8\), there exists \(y\in \mathbb Z_2^\times\) such that
\(y^2=3/11\). Hence
\[
Q(0\cdot e_1+y e_2)=11y^2=3.
\]

For \(p=3\), consider
\[
f(X,Y)=X^2+11Y^2-3.
\]
Modulo \(3\), we have
\[
f(1,1)=1+11-3=9\equiv 0\pmod 3,
\]
and
\[
\frac{\partial f}{\partial X}(1,1)=2\not\equiv 0\pmod 3.
\]
By Hensel's lemma, this solution lifts to a solution in \(\mathbb Z_3^2\).
Thus \(3\in Q(L_3)\).

For \(p=11\), we have
\[
f(5,0)=25-3=22\equiv 0\pmod{11},
\]
and
\[
\frac{\partial f}{\partial X}(5,0)=10\not\equiv 0\pmod{11}.
\]
Again Hensel's lemma gives a solution in \(\mathbb Z_{11}^2\). Hence
\(3\in Q(L_{11})\).

Now let \(p\neq 2,3,11\). Then the reduction
\[
\bar Q(X,Y)=X^2+11Y^2
\]
is a nonsingular binary quadratic form over \(\mathbb F_p\). A nonsingular
binary quadratic form over a finite field of odd characteristic represents
every element of the field. Hence there exist \(\bar x,\bar y\in\mathbb F_p\)
such that
\[
\bar x^2+11\bar y^2=3.
\]
Since \(3\neq 0\) in \(\mathbb F_p\), the vector \((\bar x,\bar y)\) is nonzero.
Because \(p\neq 2,11\), at least one of the partial derivatives
\[
2X,\qquad 22Y
\]
is nonzero at \((\bar x,\bar y)\). Therefore Hensel's lemma lifts this solution
to a solution in \(\mathbb Z_p^2\). Thus \(3\in Q(L_p)\) for every prime \(p\).

Finally, we show that \(3\notin Q(L)\). If \(3=Q(xe_1+ye_2)\) for some
\(x,y\in\mathbb Z\), then
\[
x^2+11y^2=3.
\]
If \(y\neq 0\), then \(11y^2\ge 11>3\), impossible. Hence \(y=0\), and then
\(x^2=3\), also impossible in \(\mathbb Z\). Therefore
\[
3\notin Q(L).
\]
\end{solution}
\end{document}